% EDIT THIS FILE to write your own answers. No answers are supplied.
% Find each \begin{solution}[...]{QUESTION-ID}, and type on the blank line inside.
% Do not put a % before your answer: that would hide it from the PDF.
% Compile with pdfLaTeX twice. Upload this file and/or its PDF for review.
% Preserve question IDs and MC option order. Do not regenerate after adding answers.

% Options for packages loaded elsewhere
\PassOptionsToPackage{unicode}{hyperref}
\PassOptionsToPackage{hyphens}{url}
\PassOptionsToPackage{dvipsnames,svgnames,x11names}{xcolor}
\documentclass[
  10pt,
]{article}
\usepackage{xcolor}
\usepackage[a4paper,margin=20mm,headheight=15pt]{geometry}
\usepackage{amsmath,amssymb}
\setcounter{secnumdepth}{-\maxdimen} % remove section numbering
\usepackage{iftex}
\ifPDFTeX
  \usepackage[T1]{fontenc}
  \usepackage[utf8]{inputenc}
  \usepackage{textcomp} % provide euro and other symbols
\else % if luatex or xetex
  \usepackage{unicode-math} % this also loads fontspec
  \defaultfontfeatures{Scale=MatchLowercase}
  \defaultfontfeatures[\rmfamily]{Ligatures=TeX,Scale=1}
\fi
\usepackage{lmodern}
\ifPDFTeX\else
  % xetex/luatex font selection
\fi
% Use upquote if available, for straight quotes in verbatim environments
\IfFileExists{upquote.sty}{\usepackage{upquote}}{}
\IfFileExists{microtype.sty}{% use microtype if available
  \usepackage[]{microtype}
  \UseMicrotypeSet[protrusion]{basicmath} % disable protrusion for tt fonts
}{}
\makeatletter
\@ifundefined{KOMAClassName}{% if non-KOMA class
  \IfFileExists{parskip.sty}{%
    \usepackage{parskip}
  }{% else
    \setlength{\parindent}{0pt}
    \setlength{\parskip}{6pt plus 2pt minus 1pt}}
}{% if KOMA class
  \KOMAoptions{parskip=half}}
\makeatother
\setlength{\emergencystretch}{3em} % prevent overfull lines
\providecommand{\tightlist}{%
  \setlength{\itemsep}{0pt}\setlength{\parskip}{0pt}}
\usepackage{lmodern}
\usepackage{enumitem}
\usepackage{fancyhdr}
\usepackage{titlesec}
\usepackage{xcolor}
\definecolor{examnavy}{HTML}{163653}
\definecolor{examteal}{HTML}{007F86}
\pagestyle{fancy}
\fancyhf{}
\fancyhead[L]{\small\sffamily\color{examnavy}Artificial Intelligence}
\fancyhead[R]{\small\sffamily\color{examnavy}Exam practice · Questions only}
\fancyfoot[C]{\small\thepage}
\renewcommand{\headrulewidth}{0.3pt}
\titleformat{\section}{\Large\sffamily\bfseries\color{examnavy}}{}{0pt}{}
\titleformat{\subsection}{\large\sffamily\bfseries\color{examteal}}{}{0pt}{}
\titleformat{\subsubsection}{\normalsize\sffamily\bfseries\color{examteal}}{}{0pt}{}
\titlespacing*{\subsubsection}{0pt}{5pt}{4pt}
\titlespacing*{\subsection}{0pt}{10pt}{5pt}
\setlist[itemize]{label={},leftmargin=1.4em,itemsep=2pt,topsep=3pt,parsep=0pt}
\setlength{\parindent}{0pt}
\setlength{\parskip}{5pt}
\setlength{\emergencystretch}{3em}
% This preamble is embedded in the exported workbook: no local style files are needed.
\usepackage{xparse}
\fancyhead[R]{\small\sffamily\color{examnavy}Practice workbook}
\AtBeginDocument{\hypersetup{linkcolor=examnavy,filecolor=examteal}}

% Type your response between \begin{solution} and \end{solution}.
% The optional dimension controls handwriting space ONLY when the answer is empty.
% Typed answers replace that space and may contain paragraphs and displayed maths.
\ExplSyntaxOn
\NewDocumentEnvironment{solution}{O{20mm} m +b}
  {
    \par\nopagebreak[4]\smallskip
    \tl_set:Nn \l_tmpa_tl {#3}
    \tl_remove_all:Nn \l_tmpa_tl {\par}
    \tl_if_blank:VTF \l_tmpa_tl
      {
        \noindent\begin{minipage}{\linewidth}
        {\small\sffamily\bfseries\color{examteal}My~answer~(#2):}\par
        \vspace*{#1}
        \end{minipage}\par
      }
      {
        {\small\sffamily\bfseries\color{examteal}My~answer~(#2):}\par\nopagebreak[4]
        #3\par
      }
    \medskip
  }
  {}
\ExplSyntaxOff
\usepackage{bookmark}
\IfFileExists{xurl.sty}{\usepackage{xurl}}{} % add URL line breaks if available
\urlstyle{same}
% fallback for those not using the hyperref driver hyperxmp:
\makeatletter
\@ifundefined{xmpquote}{\newcommand{\xmpquote}[1]{#1}}{}
\makeatother
\hypersetup{
  pdftitle={Artificial Intelligence: Practice Workbook},
  pdfauthor={Slide-by-slide practice},
  colorlinks=true,
  linkcolor={Maroon},
  filecolor={Maroon},
  citecolor={Blue},
  urlcolor={examteal},
  pdfcreator={LaTeX via pandoc}}

\title{Artificial Intelligence: Practice Workbook}
\usepackage{etoolbox}
\makeatletter
\providecommand{\subtitle}[1]{% add subtitle to \maketitle
  \apptocmd{\@title}{\par {\large #1 \par}}{}{}
}
\makeatother
\subtitle{380 questions with editable answer spaces}
\author{Slide-by-slide practice}
\date{}

\begin{document}
\maketitle

{
\setcounter{tocdepth}{1}
\tableofcontents
}
\clearpage

\section{Using this workbook}\label{using-this-workbook}

This workbook contains the same \textbf{380 questions}, multiple-choice
option order, question IDs, slide titles, and source page references as
\textbf{AI\_EXAM\_PRACTICE\_QUESTIONS.pdf}. It includes 18 teaching PDFs
and the separate Lab 4 supplement. \textbf{0. Presentation and 0.1. Math
background are excluded.} There are no supplied answers or worked
solutions.

\textbf{To type your solutions:} open
\textbf{AI\_EXAM\_PRACTICE\_WORKBOOK.tex} in a LaTeX editor. Each
question is followed by a \texttt{solution} environment labelled with
its ID. Type your response on the blank line between
\texttt{\textbackslash{}begin\{solution\}} and
\texttt{\textbackslash{}end\{solution\}}. Lines starting with
\texttt{\%} are comments and do not appear in the PDF: put your answer
on its own line, without \texttt{\%}. Keep the question IDs and option
order unchanged so your responses can be checked reliably.

\begin{itemize}
\tightlist
\item
  \textbf{Multiple choice:} enter A, B, C, or D. A brief explanation is
  optional.
\item
  \textbf{Short answers:} usually write 1--3 sentences or a few brief
  points.
\item
  \textbf{Calculations:} show the short calculation and your final
  result. Use \texttt{\$...\$} for inline mathematics or
  \texttt{\textbackslash{}{[}...\textbackslash{}{]}} for a displayed
  equation.
\end{itemize}

Empty answer blocks provide handwriting space. Once you type a response,
that space is replaced by your text. You can change the optional
dimension in square brackets to adjust the blank space if you prefer
printing the workbook.

\textbf{Compile with pdfLaTeX}, twice to update the table of contents.
This is a standalone \texttt{.tex} file: you can upload it directly to
Overleaf or use your usual LaTeX editor. The course slide PDFs are
needed only if you want the local source links to open; their absence
does not prevent compilation. Source references count PDF pages from 1,
including title pages.

\textbf{For feedback:} upload your completed \texttt{.tex} file,
preferably together with its compiled PDF. You may submit one section at
a time. Leave unattempted questions blank so they can be distinguished
from your answers. Save a copy before editing; do not regenerate the
workbook after entering solutions, because regeneration replaces its
contents.

This editable file is for preparation before the exam. For timed
practice, follow the announced exam rules: answer on paper, consult
offline PDFs, and do not type or use text search. The full workbook is a
question bank, not a single one-hour paper.

\clearpage

\section{1.1. Fundamentals}\label{fundamentals}

\textbf{Slides:}
\href{../Lecture\%20Slides/1.1.\%20Fundamentals.pdf}{1.1.
Fundamentals.pdf}. \textbf{20 questions: 10 MC · 7 short answers · 3
easy calculations.}

\subsection{Multiple choice}\label{multiple-choice}

% QUESTION 1.1-M01
\noindent\begin{minipage}{\linewidth}

\subsubsection{1.1-M01 · MC}\label{m01-mc}

\emph{Source:
\href{../Lecture\%20Slides/1.1.\%20Fundamentals.pdf\#page=10}{PDF
pp.~10--14}.}

Predicting a house's selling price from examples with known prices is:

\begin{itemize}
\tightlist
\item
  A. Reinforcement learning.
\item
  B. Unsupervised clustering.
\item
  C. Supervised regression.
\item
  D. Supervised classification.
\end{itemize}

\begin{solution}[8mm]{1.1-M01}
% Write A, B, C, or D here. You may add a brief reason.
C
\end{solution}
\end{minipage}\par\medskip

% QUESTION 1.1-M02
\noindent\begin{minipage}{\linewidth}

\subsubsection{1.1-M02 · MC}\label{m02-mc}

\emph{Source:
\href{../Lecture\%20Slides/1.1.\%20Fundamentals.pdf\#page=15}{PDF
pp.~15--18}.}

Which task is unsupervised learning?

\begin{itemize}
\tightlist
\item
  A. Fitting a curve to known targets.
\item
  B. Predicting labelled disease categories.
\item
  C. Learning actions from rewards.
\item
  D. Grouping observations without predefined group labels.
\end{itemize}

\begin{solution}[8mm]{1.1-M02}
% Write A, B, C, or D here. You may add a brief reason.
D
\end{solution}
\end{minipage}\par\medskip

% QUESTION 1.1-M03
\noindent\begin{minipage}{\linewidth}

\subsubsection{1.1-M03 · MC}\label{m03-mc}

\emph{Source:
\href{../Lecture\%20Slides/1.1.\%20Fundamentals.pdf\#page=18}{PDF
pp.~18--21}.}

What distinguishes reinforcement learning from ordinary supervised
learning?

\begin{itemize}
\tightlist
\item
  A. It uses adjustable model parameters.
\item
  B. Actions affect future experience, and feedback takes the form of
  rewards.
\item
  C. It always receives the correct action label.
\item
  D. It can use numerical inputs.
\end{itemize}

\begin{solution}[8mm]{1.1-M03}
% Write A, B, C, or D here. You may add a brief reason.
B
\end{solution}
\end{minipage}\par\medskip

% QUESTION 1.1-M04
\noindent\begin{minipage}{\linewidth}

\subsubsection{1.1-M04 · MC}\label{m04-mc}

\emph{Source:
\href{../Lecture\%20Slides/1.1.\%20Fundamentals.pdf\#page=6}{PDF
pp.~6}.}

An inductive bias is:

\begin{itemize}
\tightlist
\item
  A. The average training error.
\item
  B. The scalar intercept of any model.
\item
  C. An error caused only by unbalanced data.
\item
  D. An assumption that favours certain patterns or solutions.
\end{itemize}

\begin{solution}[8mm]{1.1-M04}
% Write A, B, C, or D here. You may add a brief reason.
D
\end{solution}
\end{minipage}\par\medskip

% QUESTION 1.1-M05
\noindent\begin{minipage}{\linewidth}

\subsubsection{1.1-M05 · MC}\label{m05-mc}

\emph{Source:
\href{../Lecture\%20Slides/1.1.\%20Fundamentals.pdf\#page=31}{PDF
pp.~31--32, 41}.}

Which split should guide hyperparameter selection while preserving an
independent final evaluation?

\begin{itemize}
\tightlist
\item
  A. The validation set.
\item
  B. All data merged together.
\item
  C. The training labels alone.
\item
  D. The test set after every experiment.
\end{itemize}

\begin{solution}[8mm]{1.1-M05}
% Write A, B, C, or D here. You may add a brief reason.
A
\end{solution}
\end{minipage}\par\medskip

% QUESTION 1.1-M06
\noindent\begin{minipage}{\linewidth}

\subsubsection{1.1-M06 · MC}\label{m06-mc}

\emph{Source:
\href{../Lecture\%20Slides/1.1.\%20Fundamentals.pdf\#page=31}{PDF
pp.~31--32, 41}.}

What is the main purpose of a held-out test set?

\begin{itemize}
\tightlist
\item
  A. Fit model weights.
\item
  B. Estimate performance on unseen data after model selection.
\item
  C. Increase training-set size during evaluation.
\item
  D. Select preprocessing settings repeatedly.
\end{itemize}

\begin{solution}[8mm]{1.1-M06}
% Write A, B, C, or D here. You may add a brief reason.
B
\end{solution}
\end{minipage}\par\medskip

% QUESTION 1.1-M07
\noindent\begin{minipage}{\linewidth}

\subsubsection{1.1-M07 · MC}\label{m07-mc}

\emph{Source:
\href{../Lecture\%20Slides/1.1.\%20Fundamentals.pdf\#page=33}{PDF
pp.~33--36}.}

Increasing \(k\) in k-NN classification usually makes predictions:

\begin{itemize}
\tightlist
\item
  A. Independent of the training data.
\item
  B. Equivalent to a linear classifier.
\item
  C. More sensitive to single noisy neighbours.
\item
  D. Smoother and less sensitive to individual examples.
\end{itemize}

\begin{solution}[8mm]{1.1-M07}
% Write A, B, C, or D here. You may add a brief reason.
D
\end{solution}
\end{minipage}\par\medskip

% QUESTION 1.1-M08
\noindent\begin{minipage}{\linewidth}

\subsubsection{1.1-M08 · MC}\label{m08-mc}

\emph{Source:
\href{../Lecture\%20Slides/1.1.\%20Fundamentals.pdf\#page=33}{PDF
pp.~33--34}.}

What does standard k-NN need at prediction time?

\begin{itemize}
\tightlist
\item
  A. Only one learned slope and intercept.
\item
  B. A reward function.
\item
  C. Stored training examples and a distance rule.
\item
  D. Only the class frequencies.
\end{itemize}

\begin{solution}[8mm]{1.1-M08}
% Write A, B, C, or D here. You may add a brief reason.
C
\end{solution}
\end{minipage}\par\medskip

% QUESTION 1.1-M09
\noindent\begin{minipage}{\linewidth}

\subsubsection{1.1-M09 · MC}\label{m09-mc}

\emph{Source:
\href{../Lecture\%20Slides/1.1.\%20Fundamentals.pdf\#page=26}{PDF
pp.~26}.}

A model finds a strong correlation between two variables. What follows?

\begin{itemize}
\tightlist
\item
  A. The association alone does not establish causation.
\item
  B. Confounding variables can be ruled out.
\item
  C. The first variable necessarily causes the second.
\item
  D. The second variable necessarily causes the first.
\end{itemize}

\begin{solution}[8mm]{1.1-M09}
% Write A, B, C, or D here. You may add a brief reason.
A
\end{solution}
\end{minipage}\par\medskip

% QUESTION 1.1-M10
\noindent\begin{minipage}{\linewidth}

\subsubsection{1.1-M10 · MC}\label{m10-mc}

\emph{Source:
\href{../Lecture\%20Slides/1.1.\%20Fundamentals.pdf\#page=27}{PDF
pp.~27--30}.}

Why might a very large training dataset still give poor deployment
performance?

\begin{itemize}
\tightlist
\item
  A. Its examples may poorly represent the deployment population.
\item
  B. Generalization depends only on sample count.
\item
  C. Large datasets cannot contain labels.
\item
  D. More examples necessarily increase estimation error.
\end{itemize}

\begin{solution}[8mm]{1.1-M10}
% Write A, B, C, or D here. You may add a brief reason.
A
\end{solution}
\end{minipage}\par\medskip

\subsection{Short answers}\label{short-answers}

% QUESTION 1.1-S01
\noindent\begin{minipage}{\linewidth}

\subsubsection{1.1-S01 · Short}\label{s01-short}

\emph{Source:
\href{../Lecture\%20Slides/1.1.\%20Fundamentals.pdf\#page=10}{PDF
pp.~10--17}.}

Distinguish classification from regression using the type of output.
Give one example of each.

\begin{solution}[20mm]{1.1-S01}
% Write your answer here, usually 1--3 sentences.
Classification's output values are discrete, $y\in\{0,...,n\}$ and Regression's output values come from a continuous function, $y\in\mathbb{R}$.

Regression: predict the price of a house

Classification: check if an image contains a car or not
\end{solution}
\end{minipage}\par\medskip

% QUESTION 1.1-S02
\noindent\begin{minipage}{\linewidth}

\subsubsection{1.1-S02 · Short}\label{s02-short}

\emph{Source:
\href{../Lecture\%20Slides/1.1.\%20Fundamentals.pdf\#page=31}{PDF
pp.~31--32, 41}.}

What are the roles of training, validation, and test data? Use one short
phrase per split.

\begin{solution}[20mm]{1.1-S02}
% Write your answer here, usually 1--3 sentences.
Training: learn the model's parameters values
Validation: compare results of the different models
Test: evaluate which was considered the best model. Only use this set after deciding model using validation, if not, we are introducing bias.
\end{solution}
\end{minipage}\par\medskip

% QUESTION 1.1-S03
\noindent\begin{minipage}{\linewidth}

\subsubsection{1.1-S03 · Short}\label{s03-short}

\emph{Source:
\href{../Lecture\%20Slides/1.1.\%20Fundamentals.pdf\#page=31}{PDF
pp.~31--32, 41}.}

Why can repeatedly choosing models using test results make reported test
performance misleading?

\begin{solution}[20mm]{1.1-S03}
% Write your answer here, usually 1--3 sentences.
Because if you do so, you are deciding the model depending on the test result. The model may not be learning correctly how to predict in our specific problem, and may just adjust the predictions to what he saw in the test split.
\end{solution}
\end{minipage}\par\medskip

% QUESTION 1.1-S04
\noindent\begin{minipage}{\linewidth}

\subsubsection{1.1-S04 · Short}\label{s04-short}

\emph{Source:
\href{../Lecture\%20Slides/1.1.\%20Fundamentals.pdf\#page=6}{PDF pp.~6,
30, 33--36}.}

Distinguish a learned parameter from a hyperparameter. In k-NN, which is
\(k\)?

\begin{solution}[20mm]{1.1-S04}
% Write your answer here, usually 1--3 sentences.
Learned parameter: it is used by the model to make the predictions, its value is updated while training process. For example, weights in a regression.

Hyperparameter: we choose their value depending on the behaviour we want to induct to our models. These are regularization parameters, learning rates, etc.

k in the k-NN algorithm is an hyperparameter, the model does not learn it, we choose it beforehand.
\end{solution}
\end{minipage}\par\medskip

% QUESTION 1.1-S05
\noindent\begin{minipage}{\linewidth}

\subsubsection{1.1-S05 · Short}\label{s05-short}

\emph{Source:
\href{../Lecture\%20Slides/1.1.\%20Fundamentals.pdf\#page=33}{PDF
pp.~33--36}.}

A k-NN model uses height in metres and income in euros with Euclidean
distance. Why might scaling matter?

\begin{solution}[20mm]{1.1-S05}
% Write your answer here, usually 1--3 sentences.
We scale the data because we are interested in the "relative distance" between our sample points, not in its absolute value. If we don't scale, in this example, height may have less variability than income so it can affect the result of k-NN algorithm.
\end{solution}
\end{minipage}\par\medskip

% QUESTION 1.1-S06
\noindent\begin{minipage}{\linewidth}

\subsubsection{1.1-S06 · Short}\label{s06-short}

\emph{Source:
\href{../Lecture\%20Slides/1.1.\%20Fundamentals.pdf\#page=6}{PDF
pp.~6}.}

What practical lesson does the No Free Lunch idea give about choosing a
learning algorithm before examining the problem and data?

\begin{solution}[20mm]{1.1-S06}
% Write your answer here, usually 1--3 sentences.
The lesson is that there is not such thing as best model, their behaviour and results will depend on how they fit to our data and its structure. We should try with different reasonable models making reasonable assumptions about our data, and the final model chosen may depend on the nature of the problem we are studying.
\end{solution}
\end{minipage}\par\medskip

% QUESTION 1.1-S07
\noindent\begin{minipage}{\linewidth}

\subsubsection{1.1-S07 · Short}\label{s07-short}

\emph{Source:
\href{../Lecture\%20Slides/1.1.\%20Fundamentals.pdf\#page=27}{PDF
pp.~27--30, 42--44}.}

A dataset contains many near-duplicate images from one camera. Why might
adding varied examples help more than adding more such duplicates?

\begin{solution}[20mm]{1.1-S07}
% Write your answer here, usually 1--3 sentences.
Because we want our model to train with different images, and with the more variability the best. If not, our model would overfit to the features that it learnt (the duplicate images ones) and when trying to use other images as input the results may be really bad. If we train with all samples being similar, our model will only be able to make good predictions for those kind of samples.
\end{solution}
\end{minipage}\par\medskip

\subsection{Easy calculations}\label{easy-calculations}

% QUESTION 1.1-C01
\noindent\begin{minipage}{\linewidth}

\subsubsection{1.1-C01 · Calculation}\label{c01-calculation}

\emph{Source:
\href{../Lecture\%20Slides/1.1.\%20Fundamentals.pdf\#page=33}{PDF
pp.~33--34}.}

The five nearest neighbours have labels A, B, A, A, B. What class does
unweighted 5-NN predict, and what fraction of neighbours support it?

\begin{solution}[26mm]{1.1-C01}
% Write your short calculation and final result here. Use $...$ or \[...\] for maths.
prob(A) = 3/5, prob(B) = 2/5
Our model predicts: A
\end{solution}
\end{minipage}\par\medskip

% QUESTION 1.1-C02
\noindent\begin{minipage}{\linewidth}

\subsubsection{1.1-C02 · Calculation}\label{c02-calculation}

\emph{Source:
\href{../Lecture\%20Slides/1.1.\%20Fundamentals.pdf\#page=33}{PDF
pp.~33--34}.}

For k-NN regression, suppose the rule is to average neighbours' targets.
The three nearest targets are 2, 4, and 6. Compute the prediction.

\begin{solution}[26mm]{1.1-C02}
% Write your short calculation and final result here. Use $...$ or \[...\] for maths.
With \(k=3\):
\[
\hat y=\frac{2+4+6}{3}=\boxed{4}
\]
\end{solution}
\end{minipage}\par\medskip

% QUESTION 1.1-C03
\noindent\begin{minipage}{\linewidth}

\subsubsection{1.1-C03 · Calculation}\label{c03-calculation}

\emph{Source:
\href{../Lecture\%20Slides/1.1.\%20Fundamentals.pdf\#page=31}{PDF
pp.~31--32}.}

Split 1,000 examples into 70\% training, 15\% validation, and the
remainder for testing. How many examples are in each split?

\begin{solution}[26mm]{1.1-C03}
% Write your short calculation and final result here. Use $...$ or \[...\] for maths.
Training: 700 samples
Validation: 150 samples
Test: 150 samples
\end{solution}
\end{minipage}\par\medskip

\clearpage

\section{1.2. Shallow Learning}\label{shallow-learning}

\textbf{Slides:}
\href{../Lecture\%20Slides/1.2.\%20Shallow\%20Learning.pdf}{1.2. Shallow
Learning.pdf}. \textbf{20 questions: 10 MC · 7 short answers · 3 easy
calculations.}

\subsection{Multiple choice}\label{multiple-choice-1}

% QUESTION 1.2-M01
\noindent\begin{minipage}{\linewidth}

\subsubsection{1.2-M01 · MC}\label{m01-mc-1}

\emph{Source:
\href{../Lecture\%20Slides/1.2.\%20Shallow\%20Learning.pdf\#page=5}{PDF
pp.~5, 11--12}.}

Why is \(\hat y=w_0+w_1x+w_2x^2\) linear regression with polynomial
features?

\begin{itemize}
\tightlist
\item
  A. It has only one coefficient.
\item
  B. It is a straight line in \(x\).
\item
  C. It needs no loss function.
\item
  D. It is linear in the fitted coefficients.
\end{itemize}

\begin{solution}[8mm]{1.2-M01}
% Write A, B, C, or D here. You may add a brief reason.
D
\end{solution}
\end{minipage}\par\medskip

% QUESTION 1.2-M02
\noindent\begin{minipage}{\linewidth}

\subsubsection{1.2-M02 · MC}\label{m02-mc-1}

\emph{Source:
\href{../Lecture\%20Slides/1.2.\%20Shallow\%20Learning.pdf\#page=7}{PDF
pp.~7--8}.}

Least-squares regression corresponds to maximum likelihood under which
noise assumption?

\begin{itemize}
\tightlist
\item
  A. Uniform noise with unknown support.
\item
  B. Independent Gaussian noise with fixed common variance.
\item
  C. Noise that is always zero.
\item
  D. Bernoulli noise on continuous targets.
\end{itemize}

\begin{solution}[8mm]{1.2-M02}
% Write A, B, C, or D here. You may add a brief reason.
B
\end{solution}
\end{minipage}\par\medskip

% QUESTION 1.2-M03
\noindent\begin{minipage}{\linewidth}

\subsubsection{1.2-M03 · MC}\label{m03-mc-1}

\emph{Source:
\href{../Lecture\%20Slides/1.2.\%20Shallow\%20Learning.pdf\#page=16}{PDF
pp.~16--20}.}

Increasing \(\lambda\) in \(L(w)+\lambda\|w\|_2^2\) directly:

\begin{itemize}
\tightlist
\item
  A. Strengthens the penalty on large coefficient magnitudes.
\item
  B. Removes the data-fit term.
\item
  C. Increases the learning rate.
\item
  D. Adds polynomial features.
\end{itemize}

\begin{solution}[8mm]{1.2-M03}
% Write A, B, C, or D here. You may add a brief reason.
A
\end{solution}
\end{minipage}\par\medskip

% QUESTION 1.2-M04
\noindent\begin{minipage}{\linewidth}

\subsubsection{1.2-M04 · MC}\label{m04-mc-1}

\emph{Source:
\href{../Lecture\%20Slides/1.2.\%20Shallow\%20Learning.pdf\#page=19}{PDF
pp.~19}.}

Which penalty is particularly associated with exactly zero coefficients?

\begin{itemize}
\tightlist
\item
  A. A squared residual penalty alone.
\item
  B. A larger validation set.
\item
  C. An L1 weight penalty.
\item
  D. An L2 weight penalty.
\end{itemize}

\begin{solution}[8mm]{1.2-M04}
% Write A, B, C, or D here. You may add a brief reason.
C
\end{solution}
\end{minipage}\par\medskip

% QUESTION 1.2-M05
\noindent\begin{minipage}{\linewidth}

\subsubsection{1.2-M05 · MC}\label{m05-mc-1}

\emph{Source:
\href{../Lecture\%20Slides/1.2.\%20Shallow\%20Learning.pdf\#page=49}{PDF
pp.~49--53}.}

Logistic regression outputs \(p(y=1\mid x)=0.8\). What does this
express?

\begin{itemize}
\tightlist
\item
  A. The dataset's overall accuracy.
\item
  B. A guarantee that the true class is 1.
\item
  C. The proportion of correctly fitted coefficients.
\item
  D. The model's probability for class 1 on this input.
\end{itemize}

\begin{solution}[8mm]{1.2-M05}
% Write A, B, C, or D here. You may add a brief reason.
D
\end{solution}
\end{minipage}\par\medskip

% QUESTION 1.2-M06
\noindent\begin{minipage}{\linewidth}

\subsubsection{1.2-M06 · MC}\label{m06-mc-1}

\emph{Source:
\href{../Lecture\%20Slides/1.2.\%20Shallow\%20Learning.pdf\#page=60}{PDF
pp.~60--65}.}

On data with 95\% negatives, a classifier predicts every example as
negative. Which is correct?

\begin{itemize}
\tightlist
\item
  A. Accuracy is 95\%, but positive-class recall is 0\%.
\item
  B. Positive-class recall is 95\%.
\item
  C. Positive-class F1 is 95\%.
\item
  D. It detects most positive examples.
\end{itemize}

\begin{solution}[8mm]{1.2-M06}
% Write A, B, C, or D here. You may add a brief reason.
A
\end{solution}
\end{minipage}\par\medskip

% QUESTION 1.2-M07
\noindent\begin{minipage}{\linewidth}

\subsubsection{1.2-M07 · MC}\label{m07-mc-1}

\emph{Source:
\href{../Lecture\%20Slides/1.2.\%20Shallow\%20Learning.pdf\#page=63}{PDF
pp.~63--64}.}

For fixed scores, lowering the positive-class decision threshold:

\begin{itemize}
\tightlist
\item
  A. Decreases both true positives and false positives.
\item
  B. Guarantees increased precision.
\item
  C. Changes the learned coefficients.
\item
  D. Cannot decrease recall, but can increase false positives.
\end{itemize}

\begin{solution}[8mm]{1.2-M07}
% Write A, B, C, or D here. You may add a brief reason.
D
\end{solution}
\end{minipage}\par\medskip

% QUESTION 1.2-M08
\noindent\begin{minipage}{\linewidth}

\subsubsection{1.2-M08 · MC}\label{m08-mc-1}

\emph{Source:
\href{../Lecture\%20Slides/1.2.\%20Shallow\%20Learning.pdf\#page=72}{PDF
pp.~72--76}.}

The hard-margin SVM chooses a separating hyperplane that:

\begin{itemize}
\tightlist
\item
  A. Maximizes the margin to the closest examples.
\item
  B. Minimizes the margin to the closest examples.
\item
  C. Maximizes the number of features.
\item
  D. Passes through the class means.
\end{itemize}

\begin{solution}[8mm]{1.2-M08}
% Write A, B, C, or D here. You may add a brief reason.
A
\end{solution}
\end{minipage}\par\medskip

% QUESTION 1.2-M09
\noindent\begin{minipage}{\linewidth}

\subsubsection{1.2-M09 · MC}\label{m09-mc-1}

\emph{Source:
\href{../Lecture\%20Slides/1.2.\%20Shallow\%20Learning.pdf\#page=78}{PDF
pp.~78--79}.}

The kernel trick allows an SVM to:

\begin{itemize}
\tightlist
\item
  A. Convert every task into clustering.
\item
  B. Avoid any similarity computation.
\item
  C. Learn without training examples.
\item
  D. Compute feature-space inner products without explicitly
  constructing the feature vectors.
\end{itemize}

\begin{solution}[8mm]{1.2-M09}
% Write A, B, C, or D here. You may add a brief reason.
D
\end{solution}
\end{minipage}\par\medskip

% QUESTION 1.2-M10
\noindent\begin{minipage}{\linewidth}

\subsubsection{1.2-M10 · MC}\label{m10-mc-1}

\emph{Source:
\href{../Lecture\%20Slides/1.2.\%20Shallow\%20Learning.pdf\#page=82}{PDF
pp.~82--84}.}

Which distinction between bagging and boosting is correct?

\begin{itemize}
\tightlist
\item
  A. Bagging uses only one model.
\item
  B. Random forests are linear regressors.
\item
  C. Boosting requires independent learners.
\item
  D. Bagging combines resampled learners; boosting builds learners
  sequentially to address remaining errors.
\end{itemize}

\begin{solution}[8mm]{1.2-M10}
% Write A, B, C, or D here. You may add a brief reason.
D
\end{solution}
\end{minipage}\par\medskip

\subsection{Short answers}\label{short-answers-1}

% QUESTION 1.2-S01
\noindent\begin{minipage}{\linewidth}

\subsubsection{1.2-S01 · Short}\label{s01-short-1}

\emph{Source:
\href{../Lecture\%20Slides/1.2.\%20Shallow\%20Learning.pdf\#page=13}{PDF
pp.~13--15, 21--25}.}

Training error is low but validation error is high. Name the likely
problem and one appropriate change to try.

\begin{solution}[20mm]{1.2-S01}
% Write your answer here, usually 1--3 sentences.
Overfitting, the model is "learning by heart" the features of the training dataset and is not generalizing well. One possible attempt to avoid overfitting would be implementing regularization. We could also try to use K fold cross-validation.
\end{solution}
\end{minipage}\par\medskip

% QUESTION 1.2-S02
\noindent\begin{minipage}{\linewidth}

\subsubsection{1.2-S02 · Short}\label{s02-short-1}

\emph{Source:
\href{../Lecture\%20Slides/1.2.\%20Shallow\%20Learning.pdf\#page=34}{PDF
pp.~34--39}.}

Why are squared residuals sensitive to outliers? Name one robust
alternative from the slides.

\begin{solution}[20mm]{1.2-S02}
% Write your answer here, usually 1--3 sentences.
Squared residuals give much larger penalties to large errors because the error is squared. Therefore, outliers can disproportionately influence the fitted model. Huber loss is a robust alternative, behaving quadratically for small errors and linearly for large errors.
\end{solution}
\end{minipage}\par\medskip

% QUESTION 1.2-S03
\noindent\begin{minipage}{\linewidth}

\subsubsection{1.2-S03 · Short}\label{s03-short-1}

\emph{Source:
\href{../Lecture\%20Slides/1.2.\%20Shallow\%20Learning.pdf\#page=40}{PDF
pp.~40}.}

Using \(R^2=1-\mathrm{SSE}/\mathrm{SST}\) on a nonconstant test target,
interpret \(R^2=1\), \(R^2=0\), and \(R^2<0\) relative to the test-mean
predictor.

\begin{solution}[20mm]{1.2-S03}
% Write your answer here, usually 1--3 sentences.
$R^2 = 1$: Our model predicted correctly all the samples. Prediction = Known Label for all of them.

$R^2 = 0$: Our model is exactly as good as predicting the mean value all time.

$R^2 < 0$: Our model is worse than predicting always the mean value.

\end{solution}
\end{minipage}\par\medskip

% QUESTION 1.2-S04
\noindent\begin{minipage}{\linewidth}

\subsubsection{1.2-S04 · Short}\label{s04-short-1}

\emph{Source:
\href{../Lecture\%20Slides/1.2.\%20Shallow\%20Learning.pdf\#page=60}{PDF
pp.~60--65}.}

Distinguish precision from recall in words. Which measures the fraction
of actual positives detected?

\begin{solution}[20mm]{1.2-S04}
% Write your answer here, usually 1--3 sentences.
Precision measures how much of the values predicted as positive were real positives. Recall measures how many positives we detected from the total positives that we had in our data.
\end{solution}
\end{minipage}\par\medskip

% QUESTION 1.2-S05
\noindent\begin{minipage}{\linewidth}

\subsubsection{1.2-S05 · Short}\label{s05-short-1}

\emph{Source:
\href{../Lecture\%20Slides/1.2.\%20Shallow\%20Learning.pdf\#page=65}{PDF
pp.~65}.}

What two metrics does F1 combine? Why is a large value of only one
insufficient for a high F1?

\begin{solution}[20mm]{1.2-S05}
% Write your answer here, usually 1--3 sentences.
F1 combines precision and recall. $F1 = 2\frac{Precision \cdot Recall}{Precision + Recall}$. Then it is clear that for a high value of $F1$ we need high precision and high recall.
\end{solution}
\end{minipage}\par\medskip

% QUESTION 1.2-S06
\noindent\begin{minipage}{\linewidth}

\subsubsection{1.2-S06 · Short}\label{s06-short-1}

\emph{Source:
\href{../Lecture\%20Slides/1.2.\%20Shallow\%20Learning.pdf\#page=68}{PDF
pp.~68--71}.}

Distinguish the goals of PCA and k-means. Does either require class
labels for its usual fitting procedure?

\begin{solution}[20mm]{1.2-S06}
% Write your answer here, usually 1--3 sentences.
PCA is an unsupervised dimensionality-reduction method that projects data onto directions of maximum variance. K-means is an unsupervised clustering algorithm that groups observations into \(k\) clusters based on distance to their centroids. Neither requires class labels.
\end{solution}
\end{minipage}\par\medskip

% QUESTION 1.2-S07
\noindent\begin{minipage}{\linewidth}

\subsubsection{1.2-S07 · Short}\label{s07-short-1}

\emph{Source:
\href{../Lecture\%20Slides/1.2.\%20Shallow\%20Learning.pdf\#page=81}{PDF
pp.~81--84}.}

What does a decision-tree split do? Why can a deep tree overfit, and how
can averaging randomized trees help?

\begin{solution}[20mm]{1.2-S07}
% Write your answer here, usually 1--3 sentences.
A decision-tree split divides the data according to a feature and a threshold, aiming to create more homogeneous subsets. Deep trees can overfit because they learn noise and very specific patterns from the training data. Random forests reduce variance by averaging predictions from multiple randomized trees, improving generalization.
\end{solution}
\end{minipage}\par\medskip

\subsection{Easy calculations}\label{easy-calculations-1}

% QUESTION 1.2-C01
\noindent\begin{minipage}{\linewidth}

\subsubsection{1.2-C01 · Calculation}\label{c01-calculation-1}

\emph{Source:
\href{../Lecture\%20Slides/1.2.\%20Shallow\%20Learning.pdf\#page=60}{PDF
pp.~60--61}.}

\begin{solution}[26mm]{1.2-C01}
% Write your short calculation and final result here. Use $...$ or \[...\] for maths.
Precision $= TP / (TP + FP) = 30 / (30 + 10) = 0.75$

Recall $= TP / (TP + FN) = 30 / (30 + 20) = 0.6$
\end{solution}
\end{minipage}\par\medskip

% QUESTION 1.2-C02
\noindent\begin{minipage}{\linewidth}

\subsubsection{1.2-C02 · Calculation}\label{c02-calculation-1}

\emph{Source:
\href{../Lecture\%20Slides/1.2.\%20Shallow\%20Learning.pdf\#page=40}{PDF
pp.~40}.}

Two predictions have residuals \(y-\hat y\) equal to 1 and \(-3\).
Compute their mean absolute error and mean squared error.

\begin{solution}[26mm]{1.2-C02}
% Write your short calculation and final result here. Use $...$ or \[...\] for maths.
MAE $= \frac{1}{2}(|1| + |-3|) = 2$

MSE $= \frac{1}{2}((1)^2 + (-3)^2) = 5$
\end{solution}
\end{minipage}\par\medskip

% QUESTION 1.2-C03
\noindent\begin{minipage}{\linewidth}

\subsubsection{1.2-C03 · Calculation}\label{c03-calculation-1}

\emph{Source:
\href{../Lecture\%20Slides/1.2.\%20Shallow\%20Learning.pdf\#page=16}{PDF
pp.~16--18}.}

The unregularized loss is 5, \(\|w\|_2^2=4\), and \(\lambda=0.5\).
Compute \(J=L+\lambda\|w\|_2^2\). If \(\lambda\) increases while \(w\)
stays fixed, does this objective rise or fall?

\begin{solution}[26mm]{1.2-C03}
% Write your short calculation and final result here. Use $...$ or \[...\] for maths.
$J = 5 + 0.5\cdot 4 = 7$. If $\lambda$ increases, $J$ increases too.
\end{solution}
\end{minipage}\par\medskip

\clearpage

\section{1.3. Bayesian Learning}\label{bayesian-learning}

\textbf{Slides:}
\href{../Lecture\%20Slides/1.3.\%20Bayesian\%20Learning.pdf}{1.3.
Bayesian Learning.pdf}. \textbf{20 questions: 10 MC · 7 short answers ·
3 easy calculations.}

\subsection{Multiple choice}\label{multiple-choice-2}

% QUESTION 1.3-M01
\noindent\begin{minipage}{\linewidth}

\subsubsection{1.3-M01 · MC}\label{m01-mc-2}

\emph{Source:
\href{../Lecture\%20Slides/1.3.\%20Bayesian\%20Learning.pdf\#page=5}{PDF
pp.~5--8}.}

In Bayesian parameter learning, the posterior represents:

\begin{itemize}
\tightlist
\item
  A. Fixed labels for future data.
\item
  B. Updated parameter uncertainty after observing data.
\item
  C. The optimization step size.
\item
  D. Only the input probability before fitting.
\end{itemize}

\begin{solution}[8mm]{1.3-M01}
% Write A, B, C, or D here. You may add a brief reason.

\end{solution}
\end{minipage}\par\medskip

% QUESTION 1.3-M02
\noindent\begin{minipage}{\linewidth}

\subsubsection{1.3-M02 · MC}\label{m02-mc-2}

\emph{Source:
\href{../Lecture\%20Slides/1.3.\%20Bayesian\%20Learning.pdf\#page=7}{PDF
pp.~7--8}.}

Compared with full Bayesian inference, MAP estimation returns:

\begin{itemize}
\tightlist
\item
  A. The entire predictive distribution by definition.
\item
  B. Every parameter value with equal weight.
\item
  C. A single parameter value maximizing the posterior density.
\item
  D. The likelihood with no prior contribution.
\end{itemize}

\begin{solution}[8mm]{1.3-M02}
% Write A, B, C, or D here. You may add a brief reason.

\end{solution}
\end{minipage}\par\medskip

% QUESTION 1.3-M03
\noindent\begin{minipage}{\linewidth}

\subsubsection{1.3-M03 · MC}\label{m03-mc-2}

\emph{Source:
\href{../Lecture\%20Slides/1.3.\%20Bayesian\%20Learning.pdf\#page=9}{PDF
pp.~9--14}.}

Gaussian regression likelihoods and Gaussian weight priors are
convenient because:

\begin{itemize}
\tightlist
\item
  A. The weight posterior can also be obtained in Gaussian form.
\item
  B. They guarantee nonlinear decision boundaries.
\item
  C. They eliminate uncertainty.
\item
  D. They need no observed targets.
\end{itemize}

\begin{solution}[8mm]{1.3-M03}
% Write A, B, C, or D here. You may add a brief reason.

\end{solution}
\end{minipage}\par\medskip

% QUESTION 1.3-M04
\noindent\begin{minipage}{\linewidth}

\subsubsection{1.3-M04 · MC}\label{m04-mc-2}

\emph{Source:
\href{../Lecture\%20Slides/1.3.\%20Bayesian\%20Learning.pdf\#page=13}{PDF
pp.~13}.}

A zero-mean isotropic Gaussian weight prior gives a MAP estimate
associated with:

\begin{itemize}
\tightlist
\item
  A. A tree-depth penalty.
\item
  B. An L2 penalty.
\item
  C. An entropy-only objective.
\item
  D. An L1 penalty.
\end{itemize}

\begin{solution}[8mm]{1.3-M04}
% Write A, B, C, or D here. You may add a brief reason.

\end{solution}
\end{minipage}\par\medskip

% QUESTION 1.3-M05
\noindent\begin{minipage}{\linewidth}

\subsubsection{1.3-M05 · MC}\label{m05-mc-2}

\emph{Source:
\href{../Lecture\%20Slides/1.3.\%20Bayesian\%20Learning.pdf\#page=14}{PDF
pp.~14, 19--20}.}

Two predictive distributions share a mean. The one with larger variance
expresses:

\begin{itemize}
\tightlist
\item
  A. Greater error on every sample.
\item
  B. More observations at that input.
\item
  C. A larger predicted mean target.
\item
  D. Greater predictive uncertainty under the model.
\end{itemize}

\begin{solution}[8mm]{1.3-M05}
% Write A, B, C, or D here. You may add a brief reason.

\end{solution}
\end{minipage}\par\medskip

% QUESTION 1.3-M06
\noindent\begin{minipage}{\linewidth}

\subsubsection{1.3-M06 · MC}\label{m06-mc-2}

\emph{Source:
\href{../Lecture\%20Slides/1.3.\%20Bayesian\%20Learning.pdf\#page=30}{PDF
pp.~30--35}.}

A Gaussian process specifies that:

\begin{itemize}
\tightlist
\item
  A. Function values at any finite collection of inputs have a joint
  Gaussian distribution.
\item
  B. All function values are independent.
\item
  C. Input features must be Gaussian.
\item
  D. Every sampled function is a Gaussian bell curve.
\end{itemize}

\begin{solution}[8mm]{1.3-M06}
% Write A, B, C, or D here. You may add a brief reason.

\end{solution}
\end{minipage}\par\medskip

% QUESTION 1.3-M07
\noindent\begin{minipage}{\linewidth}

\subsubsection{1.3-M07 · MC}\label{m07-mc-2}

\emph{Source:
\href{../Lecture\%20Slides/1.3.\%20Bayesian\%20Learning.pdf\#page=31}{PDF
pp.~31--38}.}

What does a GP kernel \(k(x,x')\) describe?

\begin{itemize}
\tightlist
\item
  A. The posterior mean at both inputs.
\item
  B. Covariance between the corresponding function values.
\item
  C. The observation count at \(x'\).
\item
  D. The chance of selecting \(x\) during SGD.
\end{itemize}

\begin{solution}[8mm]{1.3-M07}
% Write A, B, C, or D here. You may add a brief reason.

\end{solution}
\end{minipage}\par\medskip

% QUESTION 1.3-M08
\noindent\begin{minipage}{\linewidth}

\subsubsection{1.3-M08 · MC}\label{m08-mc-2}

\emph{Source:
\href{../Lecture\%20Slides/1.3.\%20Bayesian\%20Learning.pdf\#page=38}{PDF
pp.~38, 46--47}.}

A larger RBF length scale generally favours:

\begin{itemize}
\tightlist
\item
  A. Higher noise variance by definition.
\item
  B. More rapid variation across nearby inputs.
\item
  C. Correlations extending over larger input distances.
\item
  D. Periodicity with a shorter period.
\end{itemize}

\begin{solution}[8mm]{1.3-M08}
% Write A, B, C, or D here. You may add a brief reason.

\end{solution}
\end{minipage}\par\medskip

% QUESTION 1.3-M09
\noindent\begin{minipage}{\linewidth}

\subsubsection{1.3-M09 · MC}\label{m09-mc-2}

\emph{Source:
\href{../Lecture\%20Slides/1.3.\%20Bayesian\%20Learning.pdf\#page=42}{PDF
pp.~42--45}.}

Which kernel component suits a repeating seasonal pattern?

\begin{itemize}
\tightlist
\item
  A. Only independent white noise.
\item
  B. A linear kernel alone.
\item
  C. A constant kernel alone.
\item
  D. A periodic kernel.
\end{itemize}

\begin{solution}[8mm]{1.3-M09}
% Write A, B, C, or D here. You may add a brief reason.

\end{solution}
\end{minipage}\par\medskip

% QUESTION 1.3-M10
\noindent\begin{minipage}{\linewidth}

\subsubsection{1.3-M10 · MC}\label{m10-mc-2}

\emph{Source:
\href{../Lecture\%20Slides/1.3.\%20Bayesian\%20Learning.pdf\#page=49}{PDF
pp.~49--52}.}

GP kernel hyperparameters can be selected by optimizing:

\begin{itemize}
\tightlist
\item
  A. Test error after every update.
\item
  B. The number of target labels.
\item
  C. The marginal likelihood of the observed data.
\item
  D. Only the unconditioned prior mean.
\end{itemize}

\begin{solution}[8mm]{1.3-M10}
% Write A, B, C, or D here. You may add a brief reason.

\end{solution}
\end{minipage}\par\medskip

\subsection{Short answers}\label{short-answers-2}

% QUESTION 1.3-S01
\noindent\begin{minipage}{\linewidth}

\subsubsection{1.3-S01 · Short}\label{s01-short-2}

\emph{Source:
\href{../Lecture\%20Slides/1.3.\%20Bayesian\%20Learning.pdf\#page=5}{PDF
pp.~5--8}.}

Identify the prior, likelihood, and posterior in
\(p(w\mid D)\propto p(D\mid w)p(w)\).

\begin{solution}[20mm]{1.3-S01}
% Write your answer here, usually 1--3 sentences.

\end{solution}
\end{minipage}\par\medskip

% QUESTION 1.3-S02
\noindent\begin{minipage}{\linewidth}

\subsubsection{1.3-S02 · Short}\label{s02-short-2}

\emph{Source:
\href{../Lecture\%20Slides/1.3.\%20Bayesian\%20Learning.pdf\#page=7}{PDF
pp.~7, 14}.}

How does Bayesian prediction account for weight uncertainty instead of
using just one fitted weight vector?

\begin{solution}[20mm]{1.3-S02}
% Write your answer here, usually 1--3 sentences.

\end{solution}
\end{minipage}\par\medskip

% QUESTION 1.3-S03
\noindent\begin{minipage}{\linewidth}

\subsubsection{1.3-S03 · Short}\label{s03-short-2}

\emph{Source:
\href{../Lecture\%20Slides/1.3.\%20Bayesian\%20Learning.pdf\#page=14}{PDF
pp.~14}.}

Distinguish uncertainty about the latent function from observation
noise. Which can often be reduced by additional informative data?

\begin{solution}[20mm]{1.3-S03}
% Write your answer here, usually 1--3 sentences.

\end{solution}
\end{minipage}\par\medskip

% QUESTION 1.3-S04
\noindent\begin{minipage}{\linewidth}

\subsubsection{1.3-S04 · Short}\label{s04-short-2}

\emph{Source:
\href{../Lecture\%20Slides/1.3.\%20Bayesian\%20Learning.pdf\#page=36}{PDF
pp.~36}.}

Why is a GP called nonparametric even though its kernel has
hyperparameters?

\begin{solution}[20mm]{1.3-S04}
% Write your answer here, usually 1--3 sentences.

\end{solution}
\end{minipage}\par\medskip

% QUESTION 1.3-S05
\noindent\begin{minipage}{\linewidth}

\subsubsection{1.3-S05 · Short}\label{s05-short-2}

\emph{Source:
\href{../Lecture\%20Slides/1.3.\%20Bayesian\%20Learning.pdf\#page=46}{PDF
pp.~46--48}.}

With comparably scaled input dimensions, what does a very large ARD
length scale for one dimension suggest about its relevance?

\begin{solution}[20mm]{1.3-S05}
% Write your answer here, usually 1--3 sentences.

\end{solution}
\end{minipage}\par\medskip

% QUESTION 1.3-S06
\noindent\begin{minipage}{\linewidth}

\subsubsection{1.3-S06 · Short}\label{s06-short-2}

\emph{Source:
\href{../Lecture\%20Slides/1.3.\%20Bayesian\%20Learning.pdf\#page=44}{PDF
pp.~44--45}.}

What pattern might a sum of a linear kernel and a periodic kernel model?
Give a one-sentence example.

\begin{solution}[20mm]{1.3-S06}
% Write your answer here, usually 1--3 sentences.

\end{solution}
\end{minipage}\par\medskip

% QUESTION 1.3-S07
\noindent\begin{minipage}{\linewidth}

\subsubsection{1.3-S07 · Short}\label{s07-short-2}

\emph{Source:
\href{../Lecture\%20Slides/1.3.\%20Bayesian\%20Learning.pdf\#page=51}{PDF
pp.~51--52}.}

Why solve a linear system using a Cholesky factor rather than explicitly
form a matrix inverse in GP regression?

\begin{solution}[20mm]{1.3-S07}
% Write your answer here, usually 1--3 sentences.

\end{solution}
\end{minipage}\par\medskip

\subsection{Easy calculations}\label{easy-calculations-2}

% QUESTION 1.3-C01
\noindent\begin{minipage}{\linewidth}

\subsubsection{1.3-C01 · Calculation}\label{c01-calculation-2}

\emph{Source:
\href{../Lecture\%20Slides/1.3.\%20Bayesian\%20Learning.pdf\#page=5}{PDF
pp.~5--6}.}

Two exhaustive hypotheses have prior probabilities 0.5 each. Data
likelihoods are 0.8 under \(H_1\) and 0.2 under \(H_2\). Compute
\(P(H_1\mid D)\) by normalizing the prior-times-likelihood scores.

\begin{solution}[26mm]{1.3-C01}
% Write your short calculation and final result here. Use $...$ or \[...\] for maths.

\end{solution}
\end{minipage}\par\medskip

% QUESTION 1.3-C02
\noindent\begin{minipage}{\linewidth}

\subsubsection{1.3-C02 · Calculation}\label{c02-calculation-2}

\emph{Source:
\href{../Lecture\%20Slides/1.3.\%20Bayesian\%20Learning.pdf\#page=13}{PDF
pp.~13--14}.}

For one scalar weight, prior precision is 1 and data precision is 3.
Posterior variance is \(1/(1+3)\) and posterior mean is
\((1\cdot0+3\cdot4)/(1+3)\). Compute both posterior quantities.

\begin{solution}[26mm]{1.3-C02}
% Write your short calculation and final result here. Use $...$ or \[...\] for maths.

\end{solution}
\end{minipage}\par\medskip

% QUESTION 1.3-C03
\noindent\begin{minipage}{\linewidth}

\subsubsection{1.3-C03 · Calculation}\label{c03-calculation-2}

\emph{Source:
\href{../Lecture\%20Slides/1.3.\%20Bayesian\%20Learning.pdf\#page=14}{PDF
pp.~14}.}

Latent-function predictive variance is 0.25 and independent
observation-noise variance is 0.75. Compute the variance and standard
deviation of the noisy prediction.

\begin{solution}[26mm]{1.3-C03}
% Write your short calculation and final result here. Use $...$ or \[...\] for maths.

\end{solution}
\end{minipage}\par\medskip

\clearpage

\section{2.1. Introduction to Reinforcement
Learning}\label{introduction-to-reinforcement-learning}

\textbf{Slides:}
\href{../Lecture\%20Slides/2.1.\%20Introduction\%20to\%20Reinforcement\%20Learning.pdf}{2.1.
Introduction to Reinforcement Learning.pdf}. \textbf{20 questions: 10 MC
· 7 short answers · 3 easy calculations.}

\subsection{Multiple choice}\label{multiple-choice-3}

% QUESTION 2.1-M01
\noindent\begin{minipage}{\linewidth}

\subsubsection{2.1-M01 · MC}\label{m01-mc-3}

\emph{Source:
\href{../Lecture\%20Slides/2.1.\%20Introduction\%20to\%20Reinforcement\%20Learning.pdf\#page=3}{PDF
pp.~3--5}.}

What feedback does a standard RL agent receive?

\begin{itemize}
\tightlist
\item
  A. The gradient of the environment's hidden equations.
\item
  B. The complete optimal policy.
\item
  C. A scalar reward signal.
\item
  D. The correct action at every step.
\end{itemize}

\begin{solution}[8mm]{2.1-M01}
% Write A, B, C, or D here. You may add a brief reason.

\end{solution}
\end{minipage}\par\medskip

% QUESTION 2.1-M02
\noindent\begin{minipage}{\linewidth}

\subsubsection{2.1-M02 · MC}\label{m02-mc-3}

\emph{Source:
\href{../Lecture\%20Slides/2.1.\%20Introduction\%20to\%20Reinforcement\%20Learning.pdf\#page=5}{PDF
pp.~5}.}

The usual RL objective is to maximize:

\begin{itemize}
\tightlist
\item
  A. Expected cumulative reward.
\item
  B. Visited-state count regardless of reward.
\item
  C. The size of the action set.
\item
  D. Only the first action's reward.
\end{itemize}

\begin{solution}[8mm]{2.1-M02}
% Write A, B, C, or D here. You may add a brief reason.

\end{solution}
\end{minipage}\par\medskip

% QUESTION 2.1-M03
\noindent\begin{minipage}{\linewidth}

\subsubsection{2.1-M03 · MC}\label{m03-mc-3}

\emph{Source:
\href{../Lecture\%20Slides/2.1.\%20Introduction\%20to\%20Reinforcement\%20Learning.pdf\#page=7}{PDF
pp.~7}.}

Which describes the interaction loop?

\begin{itemize}
\tightlist
\item
  A. The agent acts; the environment supplies a successor state and
  reward.
\item
  B. The environment selects the agent's policy and ends training.
\item
  C. The agent directly changes transition probabilities.
\item
  D. The agent supplies a correct label and receives examples.
\end{itemize}

\begin{solution}[8mm]{2.1-M03}
% Write A, B, C, or D here. You may add a brief reason.

\end{solution}
\end{minipage}\par\medskip

% QUESTION 2.1-M04
\noindent\begin{minipage}{\linewidth}

\subsubsection{2.1-M04 · MC}\label{m04-mc-3}

\emph{Source:
\href{../Lecture\%20Slides/2.1.\%20Introduction\%20to\%20Reinforcement\%20Learning.pdf\#page=8}{PDF
pp.~8--9}.}

In stochastic gridworld, choosing North means:

\begin{itemize}
\tightlist
\item
  A. The state is necessarily unobservable.
\item
  B. The successor cell must be North.
\item
  C. The policy is necessarily random.
\item
  D. Actual movement is sampled from a transition distribution.
\end{itemize}

\begin{solution}[8mm]{2.1-M04}
% Write A, B, C, or D here. You may add a brief reason.

\end{solution}
\end{minipage}\par\medskip

% QUESTION 2.1-M05
\noindent\begin{minipage}{\linewidth}

\subsubsection{2.1-M05 · MC}\label{m05-mc-3}

\emph{Source:
\href{../Lecture\%20Slides/2.1.\%20Introduction\%20to\%20Reinforcement\%20Learning.pdf\#page=8}{PDF
pp.~8}.}

What happens at a terminal state in an episodic task?

\begin{itemize}
\tightlist
\item
  A. The policy is retrained from scratch.
\item
  B. Earlier rewards are erased.
\item
  C. That episode ends.
\item
  D. Its reward must be negative.
\end{itemize}

\begin{solution}[8mm]{2.1-M05}
% Write A, B, C, or D here. You may add a brief reason.

\end{solution}
\end{minipage}\par\medskip

% QUESTION 2.1-M06
\noindent\begin{minipage}{\linewidth}

\subsubsection{2.1-M06 · MC}\label{m06-mc-3}

\emph{Source:
\href{../Lecture\%20Slides/2.1.\%20Introduction\%20to\%20Reinforcement\%20Learning.pdf\#page=11}{PDF
pp.~11--12}.}

A policy describes:

\begin{itemize}
\tightlist
\item
  A. Rewards attached to terminal states.
\item
  B. The environment's dynamics.
\item
  C. The agent's behaviour: which actions to take in each state.
\item
  D. Sensor noise.
\end{itemize}

\begin{solution}[8mm]{2.1-M06}
% Write A, B, C, or D here. You may add a brief reason.

\end{solution}
\end{minipage}\par\medskip

% QUESTION 2.1-M07
\noindent\begin{minipage}{\linewidth}

\subsubsection{2.1-M07 · MC}\label{m07-mc-3}

\emph{Source:
\href{../Lecture\%20Slides/2.1.\%20Introduction\%20to\%20Reinforcement\%20Learning.pdf\#page=11}{PDF
pp.~11}.}

A transition model predicts:

\begin{itemize}
\tightlist
\item
  A. The best action without a state.
\item
  B. Only accumulated past reward.
\item
  C. The final trained weights.
\item
  D. Successor-state probabilities given the current state and action.
\end{itemize}

\begin{solution}[8mm]{2.1-M07}
% Write A, B, C, or D here. You may add a brief reason.

\end{solution}
\end{minipage}\par\medskip

% QUESTION 2.1-M08
\noindent\begin{minipage}{\linewidth}

\subsubsection{2.1-M08 · MC}\label{m08-mc-3}

\emph{Source:
\href{../Lecture\%20Slides/2.1.\%20Introduction\%20to\%20Reinforcement\%20Learning.pdf\#page=13}{PDF
pp.~13}.}

A state's value under a policy concerns:

\begin{itemize}
\tightlist
\item
  A. The available-action count.
\item
  B. Expected future return starting there and following that policy.
\item
  C. Only rewards already collected.
\item
  D. Its chance of being the initial state.
\end{itemize}

\begin{solution}[8mm]{2.1-M08}
% Write A, B, C, or D here. You may add a brief reason.

\end{solution}
\end{minipage}\par\medskip

% QUESTION 2.1-M09
\noindent\begin{minipage}{\linewidth}

\subsubsection{2.1-M09 · MC}\label{m09-mc-3}

\emph{Source:
\href{../Lecture\%20Slides/2.1.\%20Introduction\%20to\%20Reinforcement\%20Learning.pdf\#page=14}{PDF
pp.~14--16}.}

What distinguishes planning with a known model from learning by
interaction?

\begin{itemize}
\tightlist
\item
  A. Planning has no objective.
\item
  B. Planning always needs human action labels.
\item
  C. Learning needs deterministic transitions.
\item
  D. Planning computes consequences using a model; learning obtains
  information from experience.
\end{itemize}

\begin{solution}[8mm]{2.1-M09}
% Write A, B, C, or D here. You may add a brief reason.

\end{solution}
\end{minipage}\par\medskip

% QUESTION 2.1-M10
\noindent\begin{minipage}{\linewidth}

\subsubsection{2.1-M10 · MC}\label{m10-mc-3}

\emph{Source:
\href{../Lecture\%20Slides/2.1.\%20Introduction\%20to\%20Reinforcement\%20Learning.pdf\#page=16}{PDF
pp.~16}.}

Trying an unfamiliar action to learn whether it is useful is:

\begin{itemize}
\tightlist
\item
  A. Known-model policy evaluation.
\item
  B. Exploitation.
\item
  C. Reward normalization.
\item
  D. Exploration.
\end{itemize}

\begin{solution}[8mm]{2.1-M10}
% Write A, B, C, or D here. You may add a brief reason.

\end{solution}
\end{minipage}\par\medskip

\subsection{Short answers}\label{short-answers-3}

% QUESTION 2.1-S01
\noindent\begin{minipage}{\linewidth}

\subsubsection{2.1-S01 · Short}\label{s01-short-3}

\emph{Source:
\href{../Lecture\%20Slides/2.1.\%20Introduction\%20to\%20Reinforcement\%20Learning.pdf\#page=3}{PDF
pp.~3}.}

Why does an RL agent's behaviour influence the data it later learns
from?

\begin{solution}[20mm]{2.1-S01}
% Write your answer here, usually 1--3 sentences.

\end{solution}
\end{minipage}\par\medskip

% QUESTION 2.1-S02
\noindent\begin{minipage}{\linewidth}

\subsubsection{2.1-S02 · Short}\label{s02-short-3}

\emph{Source:
\href{../Lecture\%20Slides/2.1.\%20Introduction\%20to\%20Reinforcement\%20Learning.pdf\#page=5}{PDF
pp.~5--6}.}

State the relationship between maximizing reward and minimizing cost.
Give one reward or cost suitable for a walking robot.

\begin{solution}[20mm]{2.1-S02}
% Write your answer here, usually 1--3 sentences.

\end{solution}
\end{minipage}\par\medskip

% QUESTION 2.1-S03
\noindent\begin{minipage}{\linewidth}

\subsubsection{2.1-S03 · Short}\label{s03-short-3}

\emph{Source:
\href{../Lecture\%20Slides/2.1.\%20Introduction\%20to\%20Reinforcement\%20Learning.pdf\#page=7}{PDF
pp.~7, 10}.}

Distinguish the environment's true state from the agent's observation
when sensing is incomplete.

\begin{solution}[20mm]{2.1-S03}
% Write your answer here, usually 1--3 sentences.

\end{solution}
\end{minipage}\par\medskip

% QUESTION 2.1-S04
\noindent\begin{minipage}{\linewidth}

\subsubsection{2.1-S04 · Short}\label{s04-short-3}

\emph{Source:
\href{../Lecture\%20Slides/2.1.\%20Introduction\%20to\%20Reinforcement\%20Learning.pdf\#page=12}{PDF
pp.~12}.}

Distinguish deterministic and stochastic policies. Does a deterministic
policy imply deterministic transitions?

\begin{solution}[20mm]{2.1-S04}
% Write your answer here, usually 1--3 sentences.

\end{solution}
\end{minipage}\par\medskip

% QUESTION 2.1-S05
\noindent\begin{minipage}{\linewidth}

\subsubsection{2.1-S05 · Short}\label{s05-short-3}

\emph{Source:
\href{../Lecture\%20Slides/2.1.\%20Introduction\%20to\%20Reinforcement\%20Learning.pdf\#page=8}{PDF
pp.~8, 12}.}

Why can a more negative per-step living reward encourage a gridworld
agent to finish sooner?

\begin{solution}[20mm]{2.1-S05}
% Write your answer here, usually 1--3 sentences.

\end{solution}
\end{minipage}\par\medskip

% QUESTION 2.1-S06
\noindent\begin{minipage}{\linewidth}

\subsubsection{2.1-S06 · Short}\label{s06-short-3}

\emph{Source:
\href{../Lecture\%20Slides/2.1.\%20Introduction\%20to\%20Reinforcement\%20Learning.pdf\#page=13}{PDF
pp.~13}.}

What does the discount factor control? How does reducing it change the
importance of delayed rewards?

\begin{solution}[20mm]{2.1-S06}
% Write your answer here, usually 1--3 sentences.

\end{solution}
\end{minipage}\par\medskip

% QUESTION 2.1-S07
\noindent\begin{minipage}{\linewidth}

\subsubsection{2.1-S07 · Short}\label{s07-short-3}

\emph{Source:
\href{../Lecture\%20Slides/2.1.\%20Introduction\%20to\%20Reinforcement\%20Learning.pdf\#page=14}{PDF
pp.~14--16}.}

Why might always choosing the currently best-known action prevent
discovery of a better policy?

\begin{solution}[20mm]{2.1-S07}
% Write your answer here, usually 1--3 sentences.

\end{solution}
\end{minipage}\par\medskip

\subsection{Easy calculations}\label{easy-calculations-3}

% QUESTION 2.1-C01
\noindent\begin{minipage}{\linewidth}

\subsubsection{2.1-C01 · Calculation}\label{c01-calculation-3}

\emph{Source:
\href{../Lecture\%20Slides/2.1.\%20Introduction\%20to\%20Reinforcement\%20Learning.pdf\#page=8}{PDF
pp.~8}.}

When North is selected, movement is North with probability 0.8, East
with 0.1, and West with 0.1. Assume no walls. What is the probability of
sideways movement and its expected count in 10 such actions?

\begin{solution}[26mm]{2.1-C01}
% Write your short calculation and final result here. Use $...$ or \[...\] for maths.

\end{solution}
\end{minipage}\par\medskip

% QUESTION 2.1-C02
\noindent\begin{minipage}{\linewidth}

\subsubsection{2.1-C02 · Calculation}\label{c02-calculation-3}

\emph{Source:
\href{../Lecture\%20Slides/2.1.\%20Introduction\%20to\%20Reinforcement\%20Learning.pdf\#page=5}{PDF
pp.~5, 8}.}

An episode has four nonterminal moves with reward \(-1\) each, followed
by a terminal transition with reward \(+10\). With no discounting,
compute total reward.

\begin{solution}[26mm]{2.1-C02}
% Write your short calculation and final result here. Use $...$ or \[...\] for maths.

\end{solution}
\end{minipage}\par\medskip

% QUESTION 2.1-C03
\noindent\begin{minipage}{\linewidth}

\subsubsection{2.1-C03 · Calculation}\label{c03-calculation-3}

\emph{Source:
\href{../Lecture\%20Slides/2.1.\%20Introduction\%20to\%20Reinforcement\%20Learning.pdf\#page=13}{PDF
pp.~13}.}

The next two rewards are 2 and 4, after which the episode ends. With
\(\gamma=0.5\), compute \(G=2+\gamma\,4\).

\begin{solution}[26mm]{2.1-C03}
% Write your short calculation and final result here. Use $...$ or \[...\] for maths.

\end{solution}
\end{minipage}\par\medskip

\clearpage

\section{2.2. Markov decision
processes}\label{markov-decision-processes}

\textbf{Slides:}
\href{../Lecture\%20Slides/2.2.\%20Markov\%20decision\%20processes.pdf}{2.2.
Markov decision processes.pdf}. \textbf{20 questions: 10 MC · 7 short
answers · 3 easy calculations.}

\subsection{Multiple choice}\label{multiple-choice-4}

% QUESTION 2.2-M01
\noindent\begin{minipage}{\linewidth}

\subsubsection{2.2-M01 · MC}\label{m01-mc-4}

\emph{Source:
\href{../Lecture\%20Slides/2.2.\%20Markov\%20decision\%20processes.pdf\#page=2}{PDF
pp.~2--3}.}

The Markov property means that, given the current state and action:

\begin{itemize}
\tightlist
\item
  A. The history adds no information about the next-state distribution.
\item
  B. The next state must be deterministic.
\item
  C. All actions have equal value.
\item
  D. Rewards must be positive.
\end{itemize}

\begin{solution}[8mm]{2.2-M01}
% Write A, B, C, or D here. You may add a brief reason.

\end{solution}
\end{minipage}\par\medskip

% QUESTION 2.2-M02
\noindent\begin{minipage}{\linewidth}

\subsubsection{2.2-M02 · MC}\label{m02-mc-4}

\emph{Source:
\href{../Lecture\%20Slides/2.2.\%20Markov\%20decision\%20processes.pdf\#page=3}{PDF
pp.~3}.}

Which tuple describes the main ingredients of the discounted MDP in the
slides?

\begin{itemize}
\tightlist
\item
  A. Prior, likelihood, posterior, kernel, noise.
\item
  B. Features, labels, loss, optimizer, batch size.
\item
  C. States, actions, transitions, rewards, discount factor.
\item
  D. Inputs, hidden layers, outputs, weights, biases.
\end{itemize}

\begin{solution}[8mm]{2.2-M02}
% Write A, B, C, or D here. You may add a brief reason.

\end{solution}
\end{minipage}\par\medskip

% QUESTION 2.2-M03
\noindent\begin{minipage}{\linewidth}

\subsubsection{2.2-M03 · MC}\label{m03-mc-4}

\emph{Source:
\href{../Lecture\%20Slides/2.2.\%20Markov\%20decision\%20processes.pdf\#page=6}{PDF
pp.~6}.}

A stochastic policy \(\pi(a\mid s)\) gives:

\begin{itemize}
\tightlist
\item
  A. Only the reward of action \(a\).
\item
  B. A probability distribution over actions in state \(s\).
\item
  C. The probability of the next state.
\item
  D. The optimal return from \(s\).
\end{itemize}

\begin{solution}[8mm]{2.2-M03}
% Write A, B, C, or D here. You may add a brief reason.

\end{solution}
\end{minipage}\par\medskip

% QUESTION 2.2-M04
\noindent\begin{minipage}{\linewidth}

\subsubsection{2.2-M04 · MC}\label{m04-mc-4}

\emph{Source:
\href{../Lecture\%20Slides/2.2.\%20Markov\%20decision\%20processes.pdf\#page=7}{PDF
pp.~7--8}.}

If \(\gamma=0\), the discounted-return objective considers:

\begin{itemize}
\tightlist
\item
  A. Only the next immediate reward.
\item
  B. Only rewards before the current state.
\item
  C. No rewards at all.
\item
  D. All future rewards equally.
\end{itemize}

\begin{solution}[8mm]{2.2-M04}
% Write A, B, C, or D here. You may add a brief reason.

\end{solution}
\end{minipage}\par\medskip

% QUESTION 2.2-M05
\noindent\begin{minipage}{\linewidth}

\subsubsection{2.2-M05 · MC}\label{m05-mc-4}

\emph{Source:
\href{../Lecture\%20Slides/2.2.\%20Markov\%20decision\%20processes.pdf\#page=8}{PDF
pp.~8--10}.}

With bounded rewards, using \(0\leq\gamma<1\) in an infinite-horizon
problem helps ensure:

\begin{itemize}
\tightlist
\item
  A. A finite discounted return.
\item
  B. Deterministic transitions.
\item
  C. Every policy has the same value.
\item
  D. An episode ends after exactly \(1/\gamma\) steps.
\end{itemize}

\begin{solution}[8mm]{2.2-M05}
% Write A, B, C, or D here. You may add a brief reason.

\end{solution}
\end{minipage}\par\medskip

% QUESTION 2.2-M06
\noindent\begin{minipage}{\linewidth}

\subsubsection{2.2-M06 · MC}\label{m06-mc-4}

\emph{Source:
\href{../Lecture\%20Slides/2.2.\%20Markov\%20decision\%20processes.pdf\#page=11}{PDF
pp.~11}.}

What does \(Q^\pi(s,a)\) represent?

\begin{itemize}
\tightlist
\item
  A. The probability that \(a\) is optimal.
\item
  B. Only the immediate reward.
\item
  C. Expected return after taking \(a\) in \(s\) and then following
  \(\pi\).
\item
  D. The number of visits to \((s,a)\).
\end{itemize}

\begin{solution}[8mm]{2.2-M06}
% Write A, B, C, or D here. You may add a brief reason.

\end{solution}
\end{minipage}\par\medskip

% QUESTION 2.2-M07
\noindent\begin{minipage}{\linewidth}

\subsubsection{2.2-M07 · MC}\label{m07-mc-4}

\emph{Source:
\href{../Lecture\%20Slides/2.2.\%20Markov\%20decision\%20processes.pdf\#page=12}{PDF
pp.~12--15}.}

The Bellman expectation equation decomposes value into:

\begin{itemize}
\tightlist
\item
  A. Reward divided by the number of actions.
\item
  B. Expected immediate reward plus discounted successor value.
\item
  C. Prior probability and likelihood.
\item
  D. Training error and validation error.
\end{itemize}

\begin{solution}[8mm]{2.2-M07}
% Write A, B, C, or D here. You may add a brief reason.

\end{solution}
\end{minipage}\par\medskip

% QUESTION 2.2-M08
\noindent\begin{minipage}{\linewidth}

\subsubsection{2.2-M08 · MC}\label{m08-mc-4}

\emph{Source:
\href{../Lecture\%20Slides/2.2.\%20Markov\%20decision\%20processes.pdf\#page=16}{PDF
pp.~16--19}.}

An optimal policy is greedy with respect to:

\begin{itemize}
\tightlist
\item
  A. The largest transition probability alone.
\item
  B. Immediate rewards alone, for every discount factor.
\item
  C. Any randomly initialized value estimates.
\item
  D. The optimal action-value function.
\end{itemize}

\begin{solution}[8mm]{2.2-M08}
% Write A, B, C, or D here. You may add a brief reason.

\end{solution}
\end{minipage}\par\medskip

% QUESTION 2.2-M09
\noindent\begin{minipage}{\linewidth}

\subsubsection{2.2-M09 · MC}\label{m09-mc-4}

\emph{Source:
\href{../Lecture\%20Slides/2.2.\%20Markov\%20decision\%20processes.pdf\#page=20}{PDF
pp.~20--22}.}

Policy evaluation computes:

\begin{itemize}
\tightlist
\item
  A. The value of a given policy.
\item
  B. The supervised classification error.
\item
  C. Only an environment model from samples.
\item
  D. A new policy without using the current one.
\end{itemize}

\begin{solution}[8mm]{2.2-M09}
% Write A, B, C, or D here. You may add a brief reason.

\end{solution}
\end{minipage}\par\medskip

% QUESTION 2.2-M10
\noindent\begin{minipage}{\linewidth}

\subsubsection{2.2-M10 · MC}\label{m10-mc-4}

\emph{Source:
\href{../Lecture\%20Slides/2.2.\%20Markov\%20decision\%20processes.pdf\#page=20}{PDF
pp.~20--24}.}

Policy iteration repeatedly alternates:

\begin{itemize}
\tightlist
\item
  A. Clustering and dimensionality reduction.
\item
  B. Policy evaluation and policy improvement.
\item
  C. Random action selection and restarting all values.
\item
  D. Reward deletion and state deletion.
\end{itemize}

\begin{solution}[8mm]{2.2-M10}
% Write A, B, C, or D here. You may add a brief reason.

\end{solution}
\end{minipage}\par\medskip

\subsection{Short answers}\label{short-answers-4}

% QUESTION 2.2-S01
\noindent\begin{minipage}{\linewidth}

\subsubsection{2.2-S01 · Short}\label{s01-short-4}

\emph{Source:
\href{../Lecture\%20Slides/2.2.\%20Markov\%20decision\%20processes.pdf\#page=2}{PDF
pp.~2--3}.}

Why might a robot's position alone fail to be Markov if its velocity
affects what happens next?

\begin{solution}[20mm]{2.2-S01}
% Write your answer here, usually 1--3 sentences.

\end{solution}
\end{minipage}\par\medskip

% QUESTION 2.2-S02
\noindent\begin{minipage}{\linewidth}

\subsubsection{2.2-S02 · Short}\label{s02-short-4}

\emph{Source:
\href{../Lecture\%20Slides/2.2.\%20Markov\%20decision\%20processes.pdf\#page=11}{PDF
pp.~11--15}.}

Distinguish \(V^\pi(s)\) from \(Q^\pi(s,a)\). How does a stochastic
policy relate the two?

\begin{solution}[20mm]{2.2-S02}
% Write your answer here, usually 1--3 sentences.

\end{solution}
\end{minipage}\par\medskip

% QUESTION 2.2-S03
\noindent\begin{minipage}{\linewidth}

\subsubsection{2.2-S03 · Short}\label{s03-short-4}

\emph{Source:
\href{../Lecture\%20Slides/2.2.\%20Markov\%20decision\%20processes.pdf\#page=16}{PDF
pp.~16--19}.}

What changes between a Bellman expectation backup for a fixed policy and
a Bellman optimality backup?

\begin{solution}[20mm]{2.2-S03}
% Write your answer here, usually 1--3 sentences.

\end{solution}
\end{minipage}\par\medskip

% QUESTION 2.2-S04
\noindent\begin{minipage}{\linewidth}

\subsubsection{2.2-S04 · Short}\label{s04-short-4}

\emph{Source:
\href{../Lecture\%20Slides/2.2.\%20Markov\%20decision\%20processes.pdf\#page=20}{PDF
pp.~20--24}.}

How does value iteration differ from policy iteration? State the central
operation of each.

\begin{solution}[20mm]{2.2-S04}
% Write your answer here, usually 1--3 sentences.

\end{solution}
\end{minipage}\par\medskip

% QUESTION 2.2-S05
\noindent\begin{minipage}{\linewidth}

\subsubsection{2.2-S05 · Short}\label{s05-short-4}

\emph{Source:
\href{../Lecture\%20Slides/2.2.\%20Markov\%20decision\%20processes.pdf\#page=10}{PDF
pp.~10, 22}.}

Why is the continuation value after reaching a terminal state usually
zero? Does this remove the reward on the transition into it?

\begin{solution}[20mm]{2.2-S05}
% Write your answer here, usually 1--3 sentences.

\end{solution}
\end{minipage}\par\medskip

% QUESTION 2.2-S06
\noindent\begin{minipage}{\linewidth}

\subsubsection{2.2-S06 · Short}\label{s06-short-4}

\emph{Source:
\href{../Lecture\%20Slides/2.2.\%20Markov\%20decision\%20processes.pdf\#page=6}{PDF
pp.~6, 10}.}

Why might an optimal action in a finite-horizon task depend on the time
remaining as well as the state?

\begin{solution}[20mm]{2.2-S06}
% Write your answer here, usually 1--3 sentences.

\end{solution}
\end{minipage}\par\medskip

% QUESTION 2.2-S07
\noindent\begin{minipage}{\linewidth}

\subsubsection{2.2-S07 · Short}\label{s07-short-4}

\emph{Source:
\href{../Lecture\%20Slides/2.2.\%20Markov\%20decision\%20processes.pdf\#page=21}{PDF
pp.~21--24}.}

In synchronous value iteration, why should all updates in one sweep use
the preceding sweep's values?

\begin{solution}[20mm]{2.2-S07}
% Write your answer here, usually 1--3 sentences.

\end{solution}
\end{minipage}\par\medskip

\subsection{Easy calculations}\label{easy-calculations-4}

% QUESTION 2.2-C01
\noindent\begin{minipage}{\linewidth}

\subsubsection{2.2-C01 · Calculation}\label{c01-calculation-4}

\emph{Source:
\href{../Lecture\%20Slides/2.2.\%20Markov\%20decision\%20processes.pdf\#page=7}{PDF
pp.~7--9}.}

Option A gives reward 3 immediately. Option B gives 0 immediately and 8
one step later, then ends. With \(\gamma=0.5\), compute both returns and
choose the better option.

\begin{solution}[26mm]{2.2-C01}
% Write your short calculation and final result here. Use $...$ or \[...\] for maths.

\end{solution}
\end{minipage}\par\medskip

% QUESTION 2.2-C02
\noindent\begin{minipage}{\linewidth}

\subsubsection{2.2-C02 · Calculation}\label{c02-calculation-4}

\emph{Source:
\href{../Lecture\%20Slides/2.2.\%20Markov\%20decision\%20processes.pdf\#page=13}{PDF
pp.~13--15}.}

A policy selects action A with probability 0.25 and B with probability
0.75. Given \(Q^\pi(s,A)=8\) and \(Q^\pi(s,B)=4\), compute \(V^\pi(s)\).

\begin{solution}[26mm]{2.2-C02}
% Write your short calculation and final result here. Use $...$ or \[...\] for maths.

\end{solution}
\end{minipage}\par\medskip

% QUESTION 2.2-C03
\noindent\begin{minipage}{\linewidth}

\subsubsection{2.2-C03 · Calculation}\label{c03-calculation-4}

\emph{Source:
\href{../Lecture\%20Slides/2.2.\%20Markov\%20decision\%20processes.pdf\#page=18}{PDF
pp.~18--24}.}

For one value-iteration backup, action A gives reward 1 and
deterministically reaches a state with old value 6; action B terminates
with reward 3. With \(\gamma=0.5\), compute the new value of the current
state and its greedy action.

\begin{solution}[26mm]{2.2-C03}
% Write your short calculation and final result here. Use $...$ or \[...\] for maths.

\end{solution}
\end{minipage}\par\medskip

\clearpage

\section{2.3. Q-learning}\label{q-learning}

\textbf{Slides:} \href{../Lecture\%20Slides/2.3.\%20Q-learning.pdf}{2.3.
Q-learning.pdf}. \textbf{20 questions: 10 MC · 7 short answers · 3 easy
calculations.}

\subsection{Multiple choice}\label{multiple-choice-5}

% QUESTION 2.3-M01
\noindent\begin{minipage}{\linewidth}

\subsubsection{2.3-M01 · MC}\label{m01-mc-5}

\emph{Source:
\href{../Lecture\%20Slides/2.3.\%20Q-learning.pdf\#page=6}{PDF
pp.~6--10}.}

Which description is model-free RL?

\begin{itemize}
\tightlist
\item
  A. Planning only from known transition probabilities.
\item
  B. Learning values or policies from experience without requiring an
  explicit transition model.
\item
  C. Learning with no parameters of any kind.
\item
  D. Acting without observing rewards.
\end{itemize}

\begin{solution}[8mm]{2.3-M01}
% Write A, B, C, or D here. You may add a brief reason.

\end{solution}
\end{minipage}\par\medskip

% QUESTION 2.3-M02
\noindent\begin{minipage}{\linewidth}

\subsubsection{2.3-M02 · MC}\label{m02-mc-5}

\emph{Source:
\href{../Lecture\%20Slides/2.3.\%20Q-learning.pdf\#page=11}{PDF
pp.~11--16}.}

Monte Carlo policy evaluation normally uses:

\begin{itemize}
\tightlist
\item
  A. Exact expectations over a known model.
\item
  B. Only the next immediate reward.
\item
  C. Complete sampled returns from episodes.
\item
  D. Only predicted successor values.
\end{itemize}

\begin{solution}[8mm]{2.3-M02}
% Write A, B, C, or D here. You may add a brief reason.

\end{solution}
\end{minipage}\par\medskip

% QUESTION 2.3-M03
\noindent\begin{minipage}{\linewidth}

\subsubsection{2.3-M03 · MC}\label{m03-mc-5}

\emph{Source:
\href{../Lecture\%20Slides/2.3.\%20Q-learning.pdf\#page=19}{PDF
pp.~19--26}.}

Bootstrapping in TD learning means:

\begin{itemize}
\tightlist
\item
  A. Resampling a supervised dataset with replacement.
\item
  B. Initializing every value to zero.
\item
  C. Updating an estimate using another current estimate.
\item
  D. Waiting for all episodes to finish.
\end{itemize}

\begin{solution}[8mm]{2.3-M03}
% Write A, B, C, or D here. You may add a brief reason.

\end{solution}
\end{minipage}\par\medskip

% QUESTION 2.3-M04
\noindent\begin{minipage}{\linewidth}

\subsubsection{2.3-M04 · MC}\label{m04-mc-5}

\emph{Source:
\href{../Lecture\%20Slides/2.3.\%20Q-learning.pdf\#page=19}{PDF
pp.~19--24}.}

In \(V\leftarrow V+\alpha(\text{target}-V)\), \(\alpha=1\) means:

\begin{itemize}
\tightlist
\item
  A. Set the target to zero.
\item
  B. Average all past targets equally.
\item
  C. Replace the current estimate by this target.
\item
  D. Ignore the new target.
\end{itemize}

\begin{solution}[8mm]{2.3-M04}
% Write A, B, C, or D here. You may add a brief reason.

\end{solution}
\end{minipage}\par\medskip

% QUESTION 2.3-M05
\noindent\begin{minipage}{\linewidth}

\subsubsection{2.3-M05 · MC}\label{m05-mc-5}

\emph{Source:
\href{../Lecture\%20Slides/2.3.\%20Q-learning.pdf\#page=30}{PDF
pp.~30--32}.}

Under the usual \(\epsilon\)-greedy rule, increasing \(\epsilon\)
increases:

\begin{itemize}
\tightlist
\item
  A. The frequency of random action selection.
\item
  B. The number of possible states.
\item
  C. The value-update step size.
\item
  D. The discount factor.
\end{itemize}

\begin{solution}[8mm]{2.3-M05}
% Write A, B, C, or D here. You may add a brief reason.

\end{solution}
\end{minipage}\par\medskip

% QUESTION 2.3-M06
\noindent\begin{minipage}{\linewidth}

\subsubsection{2.3-M06 · MC}\label{m06-mc-5}

\emph{Source:
\href{../Lecture\%20Slides/2.3.\%20Q-learning.pdf\#page=35}{PDF
pp.~35--38}.}

The standard Q-learning target uses:

\begin{itemize}
\tightlist
\item
  A. The maximum next-state action value.
\item
  B. The average value of all states.
\item
  C. The actual sampled next action's value.
\item
  D. The minimum current-state action value.
\end{itemize}

\begin{solution}[8mm]{2.3-M06}
% Write A, B, C, or D here. You may add a brief reason.

\end{solution}
\end{minipage}\par\medskip

% QUESTION 2.3-M07
\noindent\begin{minipage}{\linewidth}

\subsubsection{2.3-M07 · MC}\label{m07-mc-5}

\emph{Source:
\href{../Lecture\%20Slides/2.3.\%20Q-learning.pdf\#page=29}{PDF pp.~29,
37--38}.}

SARSA is on-policy because its update uses:

\begin{itemize}
\tightlist
\item
  A. The next action sampled from the policy being followed and
  evaluated.
\item
  B. The best next action regardless of behaviour.
\item
  C. A known environment model.
\item
  D. Only the immediate reward.
\end{itemize}

\begin{solution}[8mm]{2.3-M07}
% Write A, B, C, or D here. You may add a brief reason.

\end{solution}
\end{minipage}\par\medskip

% QUESTION 2.3-M08
\noindent\begin{minipage}{\linewidth}

\subsubsection{2.3-M08 · MC}\label{m08-mc-5}

\emph{Source:
\href{../Lecture\%20Slides/2.3.\%20Q-learning.pdf\#page=35}{PDF
pp.~35--38}.}

Q-learning is called off-policy because:

\begin{itemize}
\tightlist
\item
  A. It requires a separate labelled dataset.
\item
  B. Its target can evaluate greedy behaviour while data come from an
  exploratory policy.
\item
  C. It cannot collect data online.
\item
  D. It never explores.
\end{itemize}

\begin{solution}[8mm]{2.3-M08}
% Write A, B, C, or D here. You may add a brief reason.

\end{solution}
\end{minipage}\par\medskip

% QUESTION 2.3-M09
\noindent\begin{minipage}{\linewidth}

\subsubsection{2.3-M09 · MC}\label{m09-mc-5}

\emph{Source:
\href{../Lecture\%20Slides/2.3.\%20Q-learning.pdf\#page=39}{PDF
pp.~39--44}.}

Function approximation helps with large state spaces by:

\begin{itemize}
\tightlist
\item
  A. Making rewards deterministic.
\item
  B. Sharing parameters and generalizing between states.
\item
  C. Storing a separate independent entry for every state.
\item
  D. Removing the need for experience.
\end{itemize}

\begin{solution}[8mm]{2.3-M09}
% Write A, B, C, or D here. You may add a brief reason.

\end{solution}
\end{minipage}\par\medskip

% QUESTION 2.3-M10
\noindent\begin{minipage}{\linewidth}

\subsubsection{2.3-M10 · MC}\label{m10-mc-5}

\emph{Source:
\href{../Lecture\%20Slides/2.3.\%20Q-learning.pdf\#page=47}{PDF
pp.~47}.}

Which three ingredients form the deadly triad associated with
instability?

\begin{itemize}
\tightlist
\item
  A. Discounting, rewards, terminal states.
\item
  B. Function approximation, bootstrapping, off-policy learning.
\item
  C. Exploration, deterministic policies, finite horizons.
\item
  D. Labels, minibatches, cross-entropy.
\end{itemize}

\begin{solution}[8mm]{2.3-M10}
% Write A, B, C, or D here. You may add a brief reason.

\end{solution}
\end{minipage}\par\medskip

\subsection{Short answers}\label{short-answers-5}

% QUESTION 2.3-S01
\noindent\begin{minipage}{\linewidth}

\subsubsection{2.3-S01 · Short}\label{s01-short-5}

\emph{Source:
\href{../Lecture\%20Slides/2.3.\%20Q-learning.pdf\#page=11}{PDF
pp.~11--26}.}

Why can one-step TD update before an episode finishes, whereas ordinary
Monte Carlo return estimation waits for the episode's outcome?

\begin{solution}[20mm]{2.3-S01}
% Write your answer here, usually 1--3 sentences.

\end{solution}
\end{minipage}\par\medskip

% QUESTION 2.3-S02
\noindent\begin{minipage}{\linewidth}

\subsubsection{2.3-S02 · Short}\label{s02-short-5}

\emph{Source:
\href{../Lecture\%20Slides/2.3.\%20Q-learning.pdf\#page=19}{PDF
pp.~19--24}.}

What is a TD error? Interpret a positive error relative to the current
value estimate.

\begin{solution}[20mm]{2.3-S02}
% Write your answer here, usually 1--3 sentences.

\end{solution}
\end{minipage}\par\medskip

% QUESTION 2.3-S03
\noindent\begin{minipage}{\linewidth}

\subsubsection{2.3-S03 · Short}\label{s03-short-5}

\emph{Source:
\href{../Lecture\%20Slides/2.3.\%20Q-learning.pdf\#page=30}{PDF
pp.~30--32}.}

Why can a purely greedy learner get stuck choosing an inferior action?
What problem arises if it keeps exploring almost all the time?

\begin{solution}[20mm]{2.3-S03}
% Write your answer here, usually 1--3 sentences.

\end{solution}
\end{minipage}\par\medskip

% QUESTION 2.3-S04
\noindent\begin{minipage}{\linewidth}

\subsubsection{2.3-S04 · Short}\label{s04-short-5}

\emph{Source:
\href{../Lecture\%20Slides/2.3.\%20Q-learning.pdf\#page=29}{PDF pp.~29,
35--38}.}

At the next state, the behaviour policy chooses a low-value exploratory
action. How can the SARSA and Q-learning targets differ?

\begin{solution}[20mm]{2.3-S04}
% Write your answer here, usually 1--3 sentences.

\end{solution}
\end{minipage}\par\medskip

% QUESTION 2.3-S05
\noindent\begin{minipage}{\linewidth}

\subsubsection{2.3-S05 · Short}\label{s05-short-5}

\emph{Source:
\href{../Lecture\%20Slides/2.3.\%20Q-learning.pdf\#page=35}{PDF
pp.~35--36}.}

For a truly terminal successor state, what continuation term is used in
the Q-learning target, and what happens to the observed terminal reward?

\begin{solution}[20mm]{2.3-S05}
% Write your answer here, usually 1--3 sentences.

\end{solution}
\end{minipage}\par\medskip

% QUESTION 2.3-S06
\noindent\begin{minipage}{\linewidth}

\subsubsection{2.3-S06 · Short}\label{s06-short-5}

\emph{Source:
\href{../Lecture\%20Slides/2.3.\%20Q-learning.pdf\#page=42}{PDF
pp.~42--44}.}

In \(Q(s,a)=\sum_i w_i f_i(s,a)\), distinguish features from learned
weights. Why can one weight update affect several state--action pairs?

\begin{solution}[20mm]{2.3-S06}
% Write your answer here, usually 1--3 sentences.

\end{solution}
\end{minipage}\par\medskip

% QUESTION 2.3-S07
\noindent\begin{minipage}{\linewidth}

\subsubsection{2.3-S07 · Short}\label{s07-short-5}

\emph{Source:
\href{../Lecture\%20Slides/2.3.\%20Q-learning.pdf\#page=31}{PDF
pp.~31}.}

What does regret measure relative to a reference best strategy? Does one
positive immediate reward necessarily mean zero regret?

\begin{solution}[20mm]{2.3-S07}
% Write your answer here, usually 1--3 sentences.

\end{solution}
\end{minipage}\par\medskip

\subsection{Easy calculations}\label{easy-calculations-5}

% QUESTION 2.3-C01
\noindent\begin{minipage}{\linewidth}

\subsubsection{2.3-C01 · Calculation}\label{c01-calculation-5}

\emph{Source:
\href{../Lecture\%20Slides/2.3.\%20Q-learning.pdf\#page=35}{PDF
pp.~35--36}.}

Perform one Q-learning update with \(Q(s,a)=2\), \(r=1\),
\(\gamma=0.5\), \(\alpha=0.5\), and \(\max_{a'}Q(s',a')=6\). First
compute the target, then \(Q_{new}=Q+\alpha(\text{target}-Q)\).

\begin{solution}[26mm]{2.3-C01}
% Write your short calculation and final result here. Use $...$ or \[...\] for maths.

\end{solution}
\end{minipage}\par\medskip

% QUESTION 2.3-C02
\noindent\begin{minipage}{\linewidth}

\subsubsection{2.3-C02 · Calculation}\label{c02-calculation-5}

\emph{Source:
\href{../Lecture\%20Slides/2.3.\%20Q-learning.pdf\#page=30}{PDF
pp.~30--32}.}

An \(\epsilon\)-greedy policy has four actions, one uniquely greedy
action, and \(\epsilon=0.2\). The random branch samples uniformly from
all four actions. Compute the probability of choosing the greedy action.

\begin{solution}[26mm]{2.3-C02}
% Write your short calculation and final result here. Use $...$ or \[...\] for maths.

\end{solution}
\end{minipage}\par\medskip

% QUESTION 2.3-C03
\noindent\begin{minipage}{\linewidth}

\subsubsection{2.3-C03 · Calculation}\label{c03-calculation-5}

\emph{Source:
\href{../Lecture\%20Slides/2.3.\%20Q-learning.pdf\#page=42}{PDF
pp.~42--44}.}

A linear approximation uses \(Q(s,a)=2f_1(s,a)-f_2(s,a)\). If \(f_1=3\)
and \(f_2=2\), compute \(Q(s,a)\). Is this value a probability?

\begin{solution}[26mm]{2.3-C03}
% Write your short calculation and final result here. Use $...$ or \[...\] for maths.

\end{solution}
\end{minipage}\par\medskip

\clearpage

\section{2.4. Policy Search}\label{policy-search}

\textbf{Slides:}
\href{../Lecture\%20Slides/2.4.\%20Policy\%20Search.pdf}{2.4. Policy
Search.pdf}. \textbf{20 questions: 10 MC · 7 short answers · 3 easy
calculations.}

\subsection{Multiple choice}\label{multiple-choice-6}

% QUESTION 2.4-M01
\noindent\begin{minipage}{\linewidth}

\subsubsection{2.4-M01 · MC}\label{m01-mc-6}

\emph{Source:
\href{../Lecture\%20Slides/2.4.\%20Policy\%20Search.pdf\#page=2}{PDF
pp.~2}.}

Policy-based RL directly parameterizes:

\begin{itemize}
\tightlist
\item
  A. Only a table of optimal Q-values.
\item
  B. Only the transition matrix.
\item
  C. The action-selection policy.
\item
  D. The test-set labels.
\end{itemize}

\begin{solution}[8mm]{2.4-M01}
% Write A, B, C, or D here. You may add a brief reason.

\end{solution}
\end{minipage}\par\medskip

% QUESTION 2.4-M02
\noindent\begin{minipage}{\linewidth}

\subsubsection{2.4-M02 · MC}\label{m02-mc-6}

\emph{Source:
\href{../Lecture\%20Slides/2.4.\%20Policy\%20Search.pdf\#page=3}{PDF
pp.~3}.}

One reason to use policy search is that it:

\begin{itemize}
\tightlist
\item
  A. Always finds a global optimum.
\item
  B. Eliminates sampling variance.
\item
  C. Does not require reward feedback.
\item
  D. Naturally supports stochastic policies and continuous actions.
\end{itemize}

\begin{solution}[8mm]{2.4-M02}
% Write A, B, C, or D here. You may add a brief reason.

\end{solution}
\end{minipage}\par\medskip

% QUESTION 2.4-M03
\noindent\begin{minipage}{\linewidth}

\subsubsection{2.4-M03 · MC}\label{m03-mc-6}

\emph{Source:
\href{../Lecture\%20Slides/2.4.\%20Policy\%20Search.pdf\#page=6}{PDF
pp.~6--7}.}

The policy objective \(J(\theta)\) is typically:

\begin{itemize}
\tightlist
\item
  A. The negative transition probability.
\item
  B. The number of policy parameters.
\item
  C. Expected return under the parameterized policy.
\item
  D. A supervised class-accuracy score by definition.
\end{itemize}

\begin{solution}[8mm]{2.4-M03}
% Write A, B, C, or D here. You may add a brief reason.

\end{solution}
\end{minipage}\par\medskip

% QUESTION 2.4-M04
\noindent\begin{minipage}{\linewidth}

\subsubsection{2.4-M04 · MC}\label{m04-mc-6}

\emph{Source:
\href{../Lecture\%20Slides/2.4.\%20Policy\%20Search.pdf\#page=7}{PDF
pp.~7}.}

To maximize \(J(\theta)\) using its gradient, the update direction is:

\begin{itemize}
\tightlist
\item
  A. Independent of \(J\).
\item
  B. \(-\nabla_\theta J\).
\item
  C. Always towards zero parameters.
\item
  D. \(+\nabla_\theta J\).
\end{itemize}

\begin{solution}[8mm]{2.4-M04}
% Write A, B, C, or D here. You may add a brief reason.

\end{solution}
\end{minipage}\par\medskip

% QUESTION 2.4-M05
\noindent\begin{minipage}{\linewidth}

\subsubsection{2.4-M05 · MC}\label{m05-mc-6}

\emph{Source:
\href{../Lecture\%20Slides/2.4.\%20Policy\%20Search.pdf\#page=11}{PDF
pp.~11}.}

A softmax policy converts action scores into:

\begin{itemize}
\tightlist
\item
  A. Rewards guaranteed to be positive.
\item
  B. State-transition probabilities.
\item
  C. Nonnegative action probabilities that sum to 1.
\item
  D. A single continuous action without a distribution.
\end{itemize}

\begin{solution}[8mm]{2.4-M05}
% Write A, B, C, or D here. You may add a brief reason.

\end{solution}
\end{minipage}\par\medskip

% QUESTION 2.4-M06
\noindent\begin{minipage}{\linewidth}

\subsubsection{2.4-M06 · MC}\label{m06-mc-6}

\emph{Source:
\href{../Lecture\%20Slides/2.4.\%20Policy\%20Search.pdf\#page=9}{PDF
pp.~9--14}.}

In REINFORCE, sampled returns primarily:

\begin{itemize}
\tightlist
\item
  A. Guarantee the best trajectory is repeated exactly.
\item
  B. Weight gradients of log action probabilities.
\item
  C. Replace the environment model with correct labels.
\item
  D. Remove the need to sample actions.
\end{itemize}

\begin{solution}[8mm]{2.4-M06}
% Write A, B, C, or D here. You may add a brief reason.

\end{solution}
\end{minipage}\par\medskip

% QUESTION 2.4-M07
\noindent\begin{minipage}{\linewidth}

\subsubsection{2.4-M07 · MC}\label{m07-mc-6}

\emph{Source:
\href{../Lecture\%20Slides/2.4.\%20Policy\%20Search.pdf\#page=9}{PDF
pp.~9, 14--15}.}

REINFORCE estimates the expected policy gradient using:

\begin{itemize}
\tightlist
\item
  A. Monte Carlo samples of trajectories or returns.
\item
  B. A supervised test-set confusion matrix.
\item
  C. Only the initial state distribution.
\item
  D. Exact enumeration of every possible trajectory.
\end{itemize}

\begin{solution}[8mm]{2.4-M07}
% Write A, B, C, or D here. You may add a brief reason.

\end{solution}
\end{minipage}\par\medskip

% QUESTION 2.4-M08
\noindent\begin{minipage}{\linewidth}

\subsubsection{2.4-M08 · MC}\label{m08-mc-6}

\emph{Source:
\href{../Lecture\%20Slides/2.4.\%20Policy\%20Search.pdf\#page=17}{PDF
pp.~17}.}

In actor--critic, the critic learns:

\begin{itemize}
\tightlist
\item
  A. Only a transition model.
\item
  B. A value estimate that helps guide policy updates.
\item
  C. The input-feature normalization.
\item
  D. A replacement for the action space.
\end{itemize}

\begin{solution}[8mm]{2.4-M08}
% Write A, B, C, or D here. You may add a brief reason.

\end{solution}
\end{minipage}\par\medskip

% QUESTION 2.4-M09
\noindent\begin{minipage}{\linewidth}

\subsubsection{2.4-M09 · MC}\label{m09-mc-6}

\emph{Source:
\href{../Lecture\%20Slides/2.4.\%20Policy\%20Search.pdf\#page=17}{PDF
pp.~17}.}

Compared with plain Monte Carlo policy gradients, a learned critic can:

\begin{itemize}
\tightlist
\item
  A. Make the reward function unnecessary.
\item
  B. Remove every source of bias and variance.
\item
  C. Guarantee a globally optimal policy.
\item
  D. Reduce variance while introducing approximation or bootstrapping
  bias.
\end{itemize}

\begin{solution}[8mm]{2.4-M09}
% Write A, B, C, or D here. You may add a brief reason.

\end{solution}
\end{minipage}\par\medskip

% QUESTION 2.4-M10
\noindent\begin{minipage}{\linewidth}

\subsubsection{2.4-M10 · MC}\label{m10-mc-6}

\emph{Source:
\href{../Lecture\%20Slides/2.4.\%20Policy\%20Search.pdf\#page=18}{PDF
pp.~18--19}.}

Which policy-search approach can be used without policy derivatives?

\begin{itemize}
\tightlist
\item
  A. Exact differentiation of \(J\).
\item
  B. An evolutionary optimization method.
\item
  C. Backpropagated policy gradient.
\item
  D. Gradient ascent through the policy network.
\end{itemize}

\begin{solution}[8mm]{2.4-M10}
% Write A, B, C, or D here. You may add a brief reason.

\end{solution}
\end{minipage}\par\medskip

\subsection{Short answers}\label{short-answers-6}

% QUESTION 2.4-S01
\noindent\begin{minipage}{\linewidth}

\subsubsection{2.4-S01 · Short}\label{s01-short-6}

\emph{Source:
\href{../Lecture\%20Slides/2.4.\%20Policy\%20Search.pdf\#page=2}{PDF
pp.~2--3}.}

Contrast obtaining a policy from Q-values with directly learning a
parameterized policy.

\begin{solution}[20mm]{2.4-S01}
% Write your answer here, usually 1--3 sentences.

\end{solution}
\end{minipage}\par\medskip

% QUESTION 2.4-S02
\noindent\begin{minipage}{\linewidth}

\subsubsection{2.4-S02 · Short}\label{s02-short-6}

\emph{Source:
\href{../Lecture\%20Slides/2.4.\%20Policy\%20Search.pdf\#page=5}{PDF
pp.~5}.}

What factors determine the probability of a sampled trajectory? Name the
roles of the initial state, policy, and environment transitions.

\begin{solution}[20mm]{2.4-S02}
% Write your answer here, usually 1--3 sentences.

\end{solution}
\end{minipage}\par\medskip

% QUESTION 2.4-S03
\noindent\begin{minipage}{\linewidth}

\subsubsection{2.4-S03 · Short}\label{s03-short-6}

\emph{Source:
\href{../Lecture\%20Slides/2.4.\%20Policy\%20Search.pdf\#page=8}{PDF
pp.~8--10}.}

Why can sampling trajectories avoid the need to enumerate all possible
futures when estimating a policy gradient?

\begin{solution}[20mm]{2.4-S03}
% Write your answer here, usually 1--3 sentences.

\end{solution}
\end{minipage}\par\medskip

% QUESTION 2.4-S04
\noindent\begin{minipage}{\linewidth}

\subsubsection{2.4-S04 · Short}\label{s04-short-6}

\emph{Source:
\href{../Lecture\%20Slides/2.4.\%20Policy\%20Search.pdf\#page=12}{PDF
pp.~12--14}.}

Why does the REINFORCE loss weight log action probabilities by returns
instead of simply imitating all sampled actions equally?

\begin{solution}[20mm]{2.4-S04}
% Write your answer here, usually 1--3 sentences.

\end{solution}
\end{minipage}\par\medskip

% QUESTION 2.4-S05
\noindent\begin{minipage}{\linewidth}

\subsubsection{2.4-S05 · Short}\label{s05-short-6}

\emph{Source:
\href{../Lecture\%20Slides/2.4.\%20Policy\%20Search.pdf\#page=14}{PDF
pp.~14--15}.}

Why can an estimate based on only one episode be noisy even when the
gradient calculation for that episode is exact?

\begin{solution}[20mm]{2.4-S05}
% Write your answer here, usually 1--3 sentences.

\end{solution}
\end{minipage}\par\medskip

% QUESTION 2.4-S06
\noindent\begin{minipage}{\linewidth}

\subsubsection{2.4-S06 · Short}\label{s06-short-6}

\emph{Source:
\href{../Lecture\%20Slides/2.4.\%20Policy\%20Search.pdf\#page=17}{PDF
pp.~17}.}

State the actor's and critic's roles in one sentence each. Which one
directly determines action probabilities?

\begin{solution}[20mm]{2.4-S06}
% Write your answer here, usually 1--3 sentences.

\end{solution}
\end{minipage}\par\medskip

% QUESTION 2.4-S07
\noindent\begin{minipage}{\linewidth}

\subsubsection{2.4-S07 · Short}\label{s07-short-6}

\emph{Source:
\href{../Lecture\%20Slides/2.4.\%20Policy\%20Search.pdf\#page=3}{PDF
pp.~3, 7, 18}.}

Does policy-gradient ascent guarantee the globally best policy? Give one
reason why optimizing a parameterized policy can be difficult.

\begin{solution}[20mm]{2.4-S07}
% Write your answer here, usually 1--3 sentences.

\end{solution}
\end{minipage}\par\medskip

\subsection{Easy calculations}\label{easy-calculations-6}

% QUESTION 2.4-C01
\noindent\begin{minipage}{\linewidth}

\subsubsection{2.4-C01 · Calculation}\label{c01-calculation-6}

\emph{Source:
\href{../Lecture\%20Slides/2.4.\%20Policy\%20Search.pdf\#page=7}{PDF
pp.~7}.}

A scalar policy parameter is \(\theta=1\), the estimated return gradient
is 2, and the ascent step size is 0.1. Compute the updated parameter.

\begin{solution}[26mm]{2.4-C01}
% Write your short calculation and final result here. Use $...$ or \[...\] for maths.

\end{solution}
\end{minipage}\par\medskip

% QUESTION 2.4-C02
\noindent\begin{minipage}{\linewidth}

\subsubsection{2.4-C02 · Calculation}\label{c02-calculation-6}

\emph{Source:
\href{../Lecture\%20Slides/2.4.\%20Policy\%20Search.pdf\#page=11}{PDF
pp.~11}.}

Two action scores have exponentials 1 and 3. Compute their softmax
probabilities without evaluating any exponentials yourself.

\begin{solution}[26mm]{2.4-C02}
% Write your short calculation and final result here. Use $...$ or \[...\] for maths.

\end{solution}
\end{minipage}\par\medskip

% QUESTION 2.4-C03
\noindent\begin{minipage}{\linewidth}

\subsubsection{2.4-C03 · Calculation}\label{c03-calculation-6}

\emph{Source:
\href{../Lecture\%20Slides/2.4.\%20Policy\%20Search.pdf\#page=6}{PDF
pp.~6, 9}.}

Three sampled episodes under the same policy have undiscounted returns
2, 4, and 6. Compute the Monte Carlo estimate of expected return. Does
this sample average prove the policy is optimal?

\begin{solution}[26mm]{2.4-C03}
% Write your short calculation and final result here. Use $...$ or \[...\] for maths.

\end{solution}
\end{minipage}\par\medskip

\clearpage

\section{3.1. Introduction to Deep
Learning}\label{introduction-to-deep-learning}

\textbf{Slides:}
\href{../Lecture\%20Slides/3.1.\%20Introduction\%20to\%20Deep\%20Learning.pdf}{3.1.
Introduction to Deep Learning.pdf}. \textbf{20 questions: 10 MC · 7
short answers · 3 easy calculations.}

\subsection{Multiple choice}\label{multiple-choice-7}

% QUESTION 3.1-M01
\noindent\begin{minipage}{\linewidth}

\subsubsection{3.1-M01 · MC}\label{m01-mc-7}

\emph{Source:
\href{../Lecture\%20Slides/3.1.\%20Introduction\%20to\%20Deep\%20Learning.pdf\#page=6}{PDF
pp.~6--10}.}

Deep learning is best described as:

\begin{itemize}
\tightlist
\item
  A. A guarantee of human-like understanding.
\item
  B. A family of machine-learning methods using multiple learned
  representation layers.
\item
  C. Every program with many lines of code.
\item
  D. Any method using a large dataset.
\end{itemize}

\begin{solution}[8mm]{3.1-M01}
% Write A, B, C, or D here. You may add a brief reason.

\end{solution}
\end{minipage}\par\medskip

% QUESTION 3.1-M02
\noindent\begin{minipage}{\linewidth}

\subsubsection{3.1-M02 · MC}\label{m02-mc-7}

\emph{Source:
\href{../Lecture\%20Slides/3.1.\%20Introduction\%20to\%20Deep\%20Learning.pdf\#page=6}{PDF
pp.~6--10}.}

Which relationship is correct?

\begin{itemize}
\tightlist
\item
  A. AI is a subset of deep learning.
\item
  B. Deep learning is part of machine learning, which is part of AI.
\item
  C. Machine learning and AI are disjoint.
\item
  D. Every AI system is a neural network.
\end{itemize}

\begin{solution}[8mm]{3.1-M02}
% Write A, B, C, or D here. You may add a brief reason.

\end{solution}
\end{minipage}\par\medskip

% QUESTION 3.1-M03
\noindent\begin{minipage}{\linewidth}

\subsubsection{3.1-M03 · MC}\label{m03-mc-7}

\emph{Source:
\href{../Lecture\%20Slides/3.1.\%20Introduction\%20to\%20Deep\%20Learning.pdf\#page=9}{PDF
pp.~9--10, 38--39}.}

Representation learning means:

\begin{itemize}
\tightlist
\item
  A. Learning useful intermediate features from data.
\item
  B. Choosing a plotting colour for every class.
\item
  C. Manually writing all useful features.
\item
  D. Storing only class labels.
\end{itemize}

\begin{solution}[8mm]{3.1-M03}
% Write A, B, C, or D here. You may add a brief reason.

\end{solution}
\end{minipage}\par\medskip

% QUESTION 3.1-M04
\noindent\begin{minipage}{\linewidth}

\subsubsection{3.1-M04 · MC}\label{m04-mc-7}

\emph{Source:
\href{../Lecture\%20Slides/3.1.\%20Introduction\%20to\%20Deep\%20Learning.pdf\#page=38}{PDF
pp.~38--39}.}

What changes in an end-to-end deep classifier compared with a fixed
handcrafted-feature pipeline?

\begin{itemize}
\tightlist
\item
  A. Feature extraction can be optimized jointly with the prediction
  stage.
\item
  B. It no longer needs a prediction objective.
\item
  C. Evaluation becomes unnecessary.
\item
  D. It no longer needs training data.
\end{itemize}

\begin{solution}[8mm]{3.1-M04}
% Write A, B, C, or D here. You may add a brief reason.

\end{solution}
\end{minipage}\par\medskip

% QUESTION 3.1-M05
\noindent\begin{minipage}{\linewidth}

\subsubsection{3.1-M05 · MC}\label{m05-mc-7}

\emph{Source:
\href{../Lecture\%20Slides/3.1.\%20Introduction\%20to\%20Deep\%20Learning.pdf\#page=13}{PDF
pp.~13--25}.}

Which statement about deep-learning applications matches the slides?

\begin{itemize}
\tightlist
\item
  A. They are limited to image classification.
\item
  B. They require only categorical outputs.
\item
  C. They include language, vision, control, and scientific tasks.
\item
  D. They exclude sequential data.
\end{itemize}

\begin{solution}[8mm]{3.1-M05}
% Write A, B, C, or D here. You may add a brief reason.

\end{solution}
\end{minipage}\par\medskip

% QUESTION 3.1-M06
\noindent\begin{minipage}{\linewidth}

\subsubsection{3.1-M06 · MC}\label{m06-mc-7}

\emph{Source:
\href{../Lecture\%20Slides/3.1.\%20Introduction\%20to\%20Deep\%20Learning.pdf\#page=15}{PDF
pp.~15--16}.}

Face detection and face recognition differ because:

\begin{itemize}
\tightlist
\item
  A. They always have identical outputs.
\item
  B. Detection locates faces; recognition concerns identity.
\item
  C. Recognition cannot use learned representations.
\item
  D. Detection identifies a person; recognition finds image coordinates
  only.
\end{itemize}

\begin{solution}[8mm]{3.1-M06}
% Write A, B, C, or D here. You may add a brief reason.

\end{solution}
\end{minipage}\par\medskip

% QUESTION 3.1-M07
\noindent\begin{minipage}{\linewidth}

\subsubsection{3.1-M07 · MC}\label{m07-mc-7}

\emph{Source:
\href{../Lecture\%20Slides/3.1.\%20Introduction\%20to\%20Deep\%20Learning.pdf\#page=26}{PDF
pp.~26--34}.}

Which is the most useful first check for a large training dataset?

\begin{itemize}
\tightlist
\item
  A. Whether its examples, labels, and variability suit the intended
  task.
\item
  B. Whether every image has the same filename length.
\item
  C. Whether the network can memorize its row order.
\item
  D. Whether the file is larger than all previous datasets.
\end{itemize}

\begin{solution}[8mm]{3.1-M07}
% Write A, B, C, or D here. You may add a brief reason.

\end{solution}
\end{minipage}\par\medskip

% QUESTION 3.1-M08
\noindent\begin{minipage}{\linewidth}

\subsubsection{3.1-M08 · MC}\label{m08-mc-7}

\emph{Source:
\href{../Lecture\%20Slides/3.1.\%20Introduction\%20to\%20Deep\%20Learning.pdf\#page=31}{PDF
pp.~31--37}.}

Which evaluation practice remains necessary in deep learning?

\begin{itemize}
\tightlist
\item
  A. Select everything on test performance.
\item
  B. Separate fitting, model selection, and final evaluation.
\item
  C. Assume depth guarantees generalization.
\item
  D. Evaluate only on training examples.
\end{itemize}

\begin{solution}[8mm]{3.1-M08}
% Write A, B, C, or D here. You may add a brief reason.

\end{solution}
\end{minipage}\par\medskip

% QUESTION 3.1-M09
\noindent\begin{minipage}{\linewidth}

\subsubsection{3.1-M09 · MC}\label{m09-mc-7}

\emph{Source:
\href{../Lecture\%20Slides/3.1.\%20Introduction\%20to\%20Deep\%20Learning.pdf\#page=34}{PDF
pp.~34--35}.}

Training an object detector normally needs which richer annotation than
image classification?

\begin{itemize}
\tightlist
\item
  A. Only image resolution.
\item
  B. Only the camera's manufacturer.
\item
  C. Only one label for the entire image.
\item
  D. Object locations as well as categories.
\end{itemize}

\begin{solution}[8mm]{3.1-M09}
% Write A, B, C, or D here. You may add a brief reason.

\end{solution}
\end{minipage}\par\medskip

% QUESTION 3.1-M10
\noindent\begin{minipage}{\linewidth}

\subsubsection{3.1-M10 · MC}\label{m10-mc-7}

\emph{Source:
\href{../Lecture\%20Slides/3.1.\%20Introduction\%20to\%20Deep\%20Learning.pdf\#page=16}{PDF
pp.~16}.}

For face-identity embeddings, a useful training objective encourages:

\begin{itemize}
\tightlist
\item
  A. Same-identity examples to be maximally distant.
\item
  B. Embeddings to depend only on file size.
\item
  C. All people to share one embedding.
\item
  D. Same-identity examples to be close and different identities
  separated.
\end{itemize}

\begin{solution}[8mm]{3.1-M10}
% Write A, B, C, or D here. You may add a brief reason.

\end{solution}
\end{minipage}\par\medskip

\subsection{Short answers}\label{short-answers-7}

% QUESTION 3.1-S01
\noindent\begin{minipage}{\linewidth}

\subsubsection{3.1-S01 · Short}\label{s01-short-7}

\emph{Source:
\href{../Lecture\%20Slides/3.1.\%20Introduction\%20to\%20Deep\%20Learning.pdf\#page=6}{PDF
pp.~6--10}.}

Give one AI approach that need not learn from data. Explain why this
shows that AI and machine learning are not synonyms.

\begin{solution}[20mm]{3.1-S01}
% Write your answer here, usually 1--3 sentences.

\end{solution}
\end{minipage}\par\medskip

% QUESTION 3.1-S02
\noindent\begin{minipage}{\linewidth}

\subsubsection{3.1-S02 · Short}\label{s02-short-7}

\emph{Source:
\href{../Lecture\%20Slides/3.1.\%20Introduction\%20to\%20Deep\%20Learning.pdf\#page=9}{PDF
pp.~9--10, 38--39}.}

What is meant by a hierarchy of learned representations? Use a simple
image example.

\begin{solution}[20mm]{3.1-S02}
% Write your answer here, usually 1--3 sentences.

\end{solution}
\end{minipage}\par\medskip

% QUESTION 3.1-S03
\noindent\begin{minipage}{\linewidth}

\subsubsection{3.1-S03 · Short}\label{s03-short-7}

\emph{Source:
\href{../Lecture\%20Slides/3.1.\%20Introduction\%20to\%20Deep\%20Learning.pdf\#page=30}{PDF
pp.~30--39}.}

Contrast a handcrafted feature with a learned feature. Which part of the
pipeline is learned in each case?

\begin{solution}[20mm]{3.1-S03}
% Write your answer here, usually 1--3 sentences.

\end{solution}
\end{minipage}\par\medskip

% QUESTION 3.1-S04
\noindent\begin{minipage}{\linewidth}

\subsubsection{3.1-S04 · Short}\label{s04-short-7}

\emph{Source:
\href{../Lecture\%20Slides/3.1.\%20Introduction\%20to\%20Deep\%20Learning.pdf\#page=27}{PDF
pp.~27--34}.}

Why can adding more images from one narrow setting fail to solve poor
generalization to another setting?

\begin{solution}[20mm]{3.1-S04}
% Write your answer here, usually 1--3 sentences.

\end{solution}
\end{minipage}\par\medskip

% QUESTION 3.1-S05
\noindent\begin{minipage}{\linewidth}

\subsubsection{3.1-S05 · Short}\label{s05-short-7}

\emph{Source:
\href{../Lecture\%20Slides/3.1.\%20Introduction\%20to\%20Deep\%20Learning.pdf\#page=34}{PDF
pp.~34--35}.}

How do the required labels differ between image classification and
pixel-level segmentation?

\begin{solution}[20mm]{3.1-S05}
% Write your answer here, usually 1--3 sentences.

\end{solution}
\end{minipage}\par\medskip

% QUESTION 3.1-S06
\noindent\begin{minipage}{\linewidth}

\subsubsection{3.1-S06 · Short}\label{s06-short-7}

\emph{Source:
\href{../Lecture\%20Slides/3.1.\%20Introduction\%20to\%20Deep\%20Learning.pdf\#page=31}{PDF
pp.~31--37}.}

A deep model performs very well on training images but poorly on new
images. Give two plausible data or methodology issues to investigate.

\begin{solution}[20mm]{3.1-S06}
% Write your answer here, usually 1--3 sentences.

\end{solution}
\end{minipage}\par\medskip

% QUESTION 3.1-S07
\noindent\begin{minipage}{\linewidth}

\subsubsection{3.1-S07 · Short}\label{s07-short-7}

\emph{Source:
\href{../Lecture\%20Slides/3.1.\%20Introduction\%20to\%20Deep\%20Learning.pdf\#page=16}{PDF
pp.~16, 38--39}.}

Why might a useful embedding support more than one downstream task,
rather than only the task used to learn it?

\begin{solution}[20mm]{3.1-S07}
% Write your answer here, usually 1--3 sentences.

\end{solution}
\end{minipage}\par\medskip

\subsection{Easy calculations}\label{easy-calculations-7}

% QUESTION 3.1-C01
\noindent\begin{minipage}{\linewidth}

\subsubsection{3.1-C01 · Calculation}\label{c01-calculation-7}

\emph{Source:
\href{../Lecture\%20Slides/3.1.\%20Introduction\%20to\%20Deep\%20Learning.pdf\#page=30}{PDF
pp.~30--32, 38--39}.}

Represent a \(16\times16\) RGB image by flattening all three colour
channels. How many scalar input values does the resulting vector
contain?

\begin{solution}[26mm]{3.1-C01}
% Write your short calculation and final result here. Use $...$ or \[...\] for maths.

\end{solution}
\end{minipage}\par\medskip

% QUESTION 3.1-C02
\noindent\begin{minipage}{\linewidth}

\subsubsection{3.1-C02 · Calculation}\label{c02-calculation-7}

\emph{Source:
\href{../Lecture\%20Slides/3.1.\%20Introduction\%20to\%20Deep\%20Learning.pdf\#page=31}{PDF
pp.~31--32}.}

For a toy dataset of 2,000 images, allocate 80\% to training and 10\% to
validation. How many remain for testing?

\begin{solution}[26mm]{3.1-C02}
% Write your short calculation and final result here. Use $...$ or \[...\] for maths.

\end{solution}
\end{minipage}\par\medskip

% QUESTION 3.1-C03
\noindent\begin{minipage}{\linewidth}

\subsubsection{3.1-C03 · Calculation}\label{c03-calculation-7}

\emph{Source:
\href{../Lecture\%20Slides/3.1.\%20Introduction\%20to\%20Deep\%20Learning.pdf\#page=34}{PDF
pp.~34--35}.}

You have 100 images of size \(8\times8\). Compare the number of target
class labels needed for image classification (one per image) with
semantic segmentation (one per pixel).

\begin{solution}[26mm]{3.1-C03}
% Write your short calculation and final result here. Use $...$ or \[...\] for maths.

\end{solution}
\end{minipage}\par\medskip

\clearpage

\section{3.2. Shallow NN}\label{shallow-nn}

\textbf{Slides:}
\href{../Lecture\%20Slides/3.2.\%20Shallow\%20NN.pdf}{3.2. Shallow
NN.pdf}. \textbf{20 questions: 10 MC · 7 short answers · 3 easy
calculations.}

\subsection{Multiple choice}\label{multiple-choice-8}

% QUESTION 3.2-M01
\noindent\begin{minipage}{\linewidth}

\subsubsection{3.2-M01 · MC}\label{m01-mc-8}

\emph{Source:
\href{../Lecture\%20Slides/3.2.\%20Shallow\%20NN.pdf\#page=5}{PDF
pp.~5--9, 49--50}.}

In an affine unit \(z=wx+b\), the bias \(b\):

\begin{itemize}
\tightlist
\item
  A. Must equal zero after training.
\item
  B. Shifts the preactivation independently of the input value.
\item
  C. Is the same thing as the loss gradient.
\item
  D. Determines the number of input samples.
\end{itemize}

\begin{solution}[8mm]{3.2-M01}
% Write A, B, C, or D here. You may add a brief reason.

\end{solution}
\end{minipage}\par\medskip

% QUESTION 3.2-M02
\noindent\begin{minipage}{\linewidth}

\subsubsection{3.2-M02 · MC}\label{m02-mc-8}

\emph{Source:
\href{../Lecture\%20Slides/3.2.\%20Shallow\%20NN.pdf\#page=7}{PDF
pp.~7--8, 51}.}

The ReLU activation is:

\begin{itemize}
\tightlist
\item
  A. A probability distribution over classes.
\item
  B. \(z^2\).
\item
  C. \(\min(0,z)\).
\item
  D. \(\max(0,z)\).
\end{itemize}

\begin{solution}[8mm]{3.2-M02}
% Write A, B, C, or D here. You may add a brief reason.

\end{solution}
\end{minipage}\par\medskip

% QUESTION 3.2-M03
\noindent\begin{minipage}{\linewidth}

\subsubsection{3.2-M03 · MC}\label{m03-mc-8}

\emph{Source:
\href{../Lecture\%20Slides/3.2.\%20Shallow\%20NN.pdf\#page=11}{PDF
pp.~11--18}.}

A 1D shallow ReLU network with an affine output represents:

\begin{itemize}
\tightlist
\item
  A. Only discontinuous step functions.
\item
  B. A continuous piecewise-linear function.
\item
  C. Only periodic functions.
\item
  D. Only one straight line.
\end{itemize}

\begin{solution}[8mm]{3.2-M03}
% Write A, B, C, or D here. You may add a brief reason.

\end{solution}
\end{minipage}\par\medskip

% QUESTION 3.2-M04
\noindent\begin{minipage}{\linewidth}

\subsubsection{3.2-M04 · MC}\label{m04-mc-8}

\emph{Source:
\href{../Lecture\%20Slides/3.2.\%20Shallow\%20NN.pdf\#page=12}{PDF
pp.~12--18}.}

An activation pattern identifies:

\begin{itemize}
\tightlist
\item
  A. The training-example order.
\item
  B. Which hidden ReLU units are active for an input.
\item
  C. The optimizer's learning rate.
\item
  D. The sample's class label.
\end{itemize}

\begin{solution}[8mm]{3.2-M04}
% Write A, B, C, or D here. You may add a brief reason.

\end{solution}
\end{minipage}\par\medskip

% QUESTION 3.2-M05
\noindent\begin{minipage}{\linewidth}

\subsubsection{3.2-M05 · MC}\label{m05-mc-8}

\emph{Source:
\href{../Lecture\%20Slides/3.2.\%20Shallow\%20NN.pdf\#page=23}{PDF
pp.~23--24}.}

The universal approximation theorem establishes, under its usual
assumptions:

\begin{itemize}
\tightlist
\item
  A. SGD always finds the best approximation.
\item
  B. Perfect generalization from a finite dataset.
\item
  C. Sufficiently wide shallow networks can approximate continuous
  functions on compact domains arbitrarily well.
\item
  D. One hidden unit is enough for every function.
\end{itemize}

\begin{solution}[8mm]{3.2-M05}
% Write A, B, C, or D here. You may add a brief reason.

\end{solution}
\end{minipage}\par\medskip

% QUESTION 3.2-M06
\noindent\begin{minipage}{\linewidth}

\subsubsection{3.2-M06 · MC}\label{m06-mc-8}

\emph{Source:
\href{../Lecture\%20Slides/3.2.\%20Shallow\%20NN.pdf\#page=24}{PDF
pp.~24}.}

What does universal approximation alone NOT guarantee?

\begin{itemize}
\tightlist
\item
  A. That approximation accuracy can depend on width.
\item
  B. That training will discover suitable parameters efficiently.
\item
  C. That hidden units may be needed.
\item
  D. The existence of useful nonlinear approximations.
\end{itemize}

\begin{solution}[8mm]{3.2-M06}
% Write A, B, C, or D here. You may add a brief reason.

\end{solution}
\end{minipage}\par\medskip

% QUESTION 3.2-M07
\noindent\begin{minipage}{\linewidth}

\subsubsection{3.2-M07 · MC}\label{m07-mc-8}

\emph{Source:
\href{../Lecture\%20Slides/3.2.\%20Shallow\%20NN.pdf\#page=26}{PDF
pp.~26--28}.}

In a shallow network with two outputs and one shared hidden layer:

\begin{itemize}
\tightlist
\item
  A. The hidden units cannot be shared.
\item
  B. The two outputs require identical output biases.
\item
  C. Each output must equal the other.
\item
  D. Each output can use a different weighted combination of the same
  hidden activations.
\end{itemize}

\begin{solution}[8mm]{3.2-M07}
% Write A, B, C, or D here. You may add a brief reason.

\end{solution}
\end{minipage}\par\medskip

% QUESTION 3.2-M08
\noindent\begin{minipage}{\linewidth}

\subsubsection{3.2-M08 · MC}\label{m08-mc-8}

\emph{Source:
\href{../Lecture\%20Slides/3.2.\%20Shallow\%20NN.pdf\#page=30}{PDF
pp.~30--42}.}

For column-vector inputs \(x\in\mathbb{R}^{d}\) and \(h\) hidden units,
\(Wx+b\) uses a weight matrix of shape:

\begin{itemize}
\tightlist
\item
  A. \(h\times h\).
\item
  B. \(d\times h\) when used directly as \(Wx\).
\item
  C. \(d\times d\).
\item
  D. \(h\times d\).
\end{itemize}

\begin{solution}[8mm]{3.2-M08}
% Write A, B, C, or D here. You may add a brief reason.

\end{solution}
\end{minipage}\par\medskip

% QUESTION 3.2-M09
\noindent\begin{minipage}{\linewidth}

\subsubsection{3.2-M09 · MC}\label{m09-mc-8}

\emph{Source:
\href{../Lecture\%20Slides/3.2.\%20Shallow\%20NN.pdf\#page=39}{PDF
pp.~39--42}.}

For a fully connected layer with one bias per output, the parameter
count depends on:

\begin{itemize}
\tightlist
\item
  A. The random seed alone.
\item
  B. The number of evaluation examples.
\item
  C. Only the number of training batches.
\item
  D. Input width and output width.
\end{itemize}

\begin{solution}[8mm]{3.2-M09}
% Write A, B, C, or D here. You may add a brief reason.

\end{solution}
\end{minipage}\par\medskip

% QUESTION 3.2-M10
\noindent\begin{minipage}{\linewidth}

\subsubsection{3.2-M10 · MC}\label{m10-mc-8}

\emph{Source:
\href{../Lecture\%20Slides/3.2.\%20Shallow\%20NN.pdf\#page=18}{PDF
pp.~18, 30--46}.}

The linear regions of a ReLU network partition:

\begin{itemize}
\tightlist
\item
  A. Input space according to hidden activation patterns.
\item
  B. The training epochs into equal intervals.
\item
  C. The set of optimizer names.
\item
  D. Only the parameter file on disk.
\end{itemize}

\begin{solution}[8mm]{3.2-M10}
% Write A, B, C, or D here. You may add a brief reason.

\end{solution}
\end{minipage}\par\medskip

\subsection{Short answers}\label{short-answers-8}

% QUESTION 3.2-S01
\noindent\begin{minipage}{\linewidth}

\subsubsection{3.2-S01 · Short}\label{s01-short-8}

\emph{Source:
\href{../Lecture\%20Slides/3.2.\%20Shallow\%20NN.pdf\#page=5}{PDF
pp.~5--9, 12--16}.}

Describe the forward pass of a one-hidden-layer network in three steps,
from inputs to output.

\begin{solution}[20mm]{3.2-S01}
% Write your answer here, usually 1--3 sentences.

\end{solution}
\end{minipage}\par\medskip

% QUESTION 3.2-S02
\noindent\begin{minipage}{\linewidth}

\subsubsection{3.2-S02 · Short}\label{s02-short-8}

\emph{Source:
\href{../Lecture\%20Slides/3.2.\%20Shallow\%20NN.pdf\#page=7}{PDF
pp.~7--18}.}

Why are nonlinear activations needed if we want more than an affine
input--output mapping?

\begin{solution}[20mm]{3.2-S02}
% Write your answer here, usually 1--3 sentences.

\end{solution}
\end{minipage}\par\medskip

% QUESTION 3.2-S03
\noindent\begin{minipage}{\linewidth}

\subsubsection{3.2-S03 · Short}\label{s03-short-8}

\emph{Source:
\href{../Lecture\%20Slides/3.2.\%20Shallow\%20NN.pdf\#page=12}{PDF
pp.~12--18}.}

What happens to a ReLU hidden unit's contribution when its preactivation
becomes negative?

\begin{solution}[20mm]{3.2-S03}
% Write your answer here, usually 1--3 sentences.

\end{solution}
\end{minipage}\par\medskip

% QUESTION 3.2-S04
\noindent\begin{minipage}{\linewidth}

\subsubsection{3.2-S04 · Short}\label{s04-short-8}

\emph{Source:
\href{../Lecture\%20Slides/3.2.\%20Shallow\%20NN.pdf\#page=23}{PDF
pp.~23--24}.}

Distinguish a network's ability to represent a function from its ability
to learn that function from limited data.

\begin{solution}[20mm]{3.2-S04}
% Write your answer here, usually 1--3 sentences.

\end{solution}
\end{minipage}\par\medskip

% QUESTION 3.2-S05
\noindent\begin{minipage}{\linewidth}

\subsubsection{3.2-S05 · Short}\label{s05-short-8}

\emph{Source:
\href{../Lecture\%20Slides/3.2.\%20Shallow\%20NN.pdf\#page=26}{PDF
pp.~26--28, 37}.}

If a second output is added while the hidden layer stays fixed, which
parameters must be added? Do the existing input-to-hidden weights need
to be duplicated?

\begin{solution}[20mm]{3.2-S05}
% Write your answer here, usually 1--3 sentences.

\end{solution}
\end{minipage}\par\medskip

% QUESTION 3.2-S06
\noindent\begin{minipage}{\linewidth}

\subsubsection{3.2-S06 · Short}\label{s06-short-8}

\emph{Source:
\href{../Lecture\%20Slides/3.2.\%20Shallow\%20NN.pdf\#page=39}{PDF
pp.~39--42, 49--50}.}

Distinguish input dimension, hidden-layer width, and output dimension.
Give a situation where output dimension exceeds one.

\begin{solution}[20mm]{3.2-S06}
% Write your answer here, usually 1--3 sentences.

\end{solution}
\end{minipage}\par\medskip

% QUESTION 3.2-S07
\noindent\begin{minipage}{\linewidth}

\subsubsection{3.2-S07 · Short}\label{s07-short-8}

\emph{Source:
\href{../Lecture\%20Slides/3.2.\%20Shallow\%20NN.pdf\#page=18}{PDF
pp.~18, 44--46}.}

Why can a ReLU network be linear within each activation region but
nonlinear over the full input space?

\begin{solution}[20mm]{3.2-S07}
% Write your answer here, usually 1--3 sentences.

\end{solution}
\end{minipage}\par\medskip

\subsection{Easy calculations}\label{easy-calculations-8}

% QUESTION 3.2-C01
\noindent\begin{minipage}{\linewidth}

\subsubsection{3.2-C01 · Calculation}\label{c01-calculation-8}

\emph{Source:
\href{../Lecture\%20Slides/3.2.\%20Shallow\%20NN.pdf\#page=39}{PDF
pp.~39--42}.}

A fully connected shallow network has 2 inputs, 3 hidden units, and 1
output. Include one bias for every hidden and output unit. Compute the
total number of trainable parameters.

\begin{solution}[26mm]{3.2-C01}
% Write your short calculation and final result here. Use $...$ or \[...\] for maths.

\end{solution}
\end{minipage}\par\medskip

% QUESTION 3.2-C02
\noindent\begin{minipage}{\linewidth}

\subsubsection{3.2-C02 · Calculation}\label{c02-calculation-8}

\emph{Source:
\href{../Lecture\%20Slides/3.2.\%20Shallow\%20NN.pdf\#page=7}{PDF
pp.~7--16}.}

A hidden unit computes \(h=\max(0,2x-3)\). Evaluate \(h\) at \(x=1\) and
at \(x=2\).

\begin{solution}[26mm]{3.2-C02}
% Write your short calculation and final result here. Use $...$ or \[...\] for maths.

\end{solution}
\end{minipage}\par\medskip

% QUESTION 3.2-C03
\noindent\begin{minipage}{\linewidth}

\subsubsection{3.2-C03 · Calculation}\label{c03-calculation-8}

\emph{Source:
\href{../Lecture\%20Slides/3.2.\%20Shallow\%20NN.pdf\#page=17}{PDF
pp.~17, 22, 44--47}.}

A 1D shallow network has three ReLU hidden units with three distinct
breakpoints. Assume their slope changes survive in the output. How many
linear intervals can these breakpoints create?

\begin{solution}[26mm]{3.2-C03}
% Write your short calculation and final result here. Use $...$ or \[...\] for maths.

\end{solution}
\end{minipage}\par\medskip

\clearpage

\section{3.3. Deep Learning
Fundamentals}\label{deep-learning-fundamentals}

\textbf{Slides:}
\href{../Lecture\%20Slides/3.3.\%20Deep\%20Learning\%20Fundamentals.pdf}{3.3.
Deep Learning Fundamentals.pdf}. \textbf{20 questions: 10 MC · 7 short
answers · 3 easy calculations.}

\subsection{Multiple choice}\label{multiple-choice-9}

% QUESTION 3.3-M01
\noindent\begin{minipage}{\linewidth}

\subsubsection{3.3-M01 · MC}\label{m01-mc-9}

\emph{Source:
\href{../Lecture\%20Slides/3.3.\%20Deep\%20Learning\%20Fundamentals.pdf\#page=5}{PDF
pp.~5--8}.}

Network depth and width refer respectively to:

\begin{itemize}
\tightlist
\item
  A. Image height and image width in every model.
\item
  B. Learning rate and batch size.
\item
  C. Training-set size and test-set size.
\item
  D. Number of successive layers and number of units in a layer.
\end{itemize}

\begin{solution}[8mm]{3.3-M01}
% Write A, B, C, or D here. You may add a brief reason.

\end{solution}
\end{minipage}\par\medskip

% QUESTION 3.3-M02
\noindent\begin{minipage}{\linewidth}

\subsubsection{3.3-M02 · MC}\label{m02-mc-9}

\emph{Source:
\href{../Lecture\%20Slides/3.3.\%20Deep\%20Learning\%20Fundamentals.pdf\#page=5}{PDF
pp.~5--14}.}

Why can depth help represent some functions efficiently?

\begin{itemize}
\tightlist
\item
  A. It composes intermediate transformations instead of representing
  everything in one wide layer.
\item
  B. It guarantees lower test error with every layer.
\item
  C. It makes parameter learning unnecessary.
\item
  D. It removes nonlinearities.
\end{itemize}

\begin{solution}[8mm]{3.3-M02}
% Write A, B, C, or D here. You may add a brief reason.

\end{solution}
\end{minipage}\par\medskip

% QUESTION 3.3-M03
\noindent\begin{minipage}{\linewidth}

\subsubsection{3.3-M03 · MC}\label{m03-mc-9}

\emph{Source:
\href{../Lecture\%20Slides/3.3.\%20Deep\%20Learning\%20Fundamentals.pdf\#page=3}{PDF
pp.~3--10}.}

Composing affine layers without nonlinear activations gives:

\begin{itemize}
\tightlist
\item
  A. A kernel classifier.
\item
  B. An arbitrary nonlinear function.
\item
  C. Another affine mapping.
\item
  D. A piecewise-constant function.
\end{itemize}

\begin{solution}[8mm]{3.3-M03}
% Write A, B, C, or D here. You may add a brief reason.

\end{solution}
\end{minipage}\par\medskip

% QUESTION 3.3-M04
\noindent\begin{minipage}{\linewidth}

\subsubsection{3.3-M04 · MC}\label{m04-mc-9}

\emph{Source:
\href{../Lecture\%20Slides/3.3.\%20Deep\%20Learning\%20Fundamentals.pdf\#page=17}{PDF
pp.~17--24}.}

A raw class score or logit:

\begin{itemize}
\tightlist
\item
  A. Must lie in \([0,1]\).
\item
  B. Is already a calibrated probability.
\item
  C. Must sum to 1 across classes.
\item
  D. Can be any real value before probability normalization.
\end{itemize}

\begin{solution}[8mm]{3.3-M04}
% Write A, B, C, or D here. You may add a brief reason.

\end{solution}
\end{minipage}\par\medskip

% QUESTION 3.3-M05
\noindent\begin{minipage}{\linewidth}

\subsubsection{3.3-M05 · MC}\label{m05-mc-9}

\emph{Source:
\href{../Lecture\%20Slides/3.3.\%20Deep\%20Learning\%20Fundamentals.pdf\#page=24}{PDF
pp.~24--26}.}

Softmax transforms logits into:

\begin{itemize}
\tightlist
\item
  A. Independent unbounded scores.
\item
  B. Nonnegative probabilities summing to 1.
\item
  C. One-hot vectors in every forward pass.
\item
  D. Parameter gradients.
\end{itemize}

\begin{solution}[8mm]{3.3-M05}
% Write A, B, C, or D here. You may add a brief reason.

\end{solution}
\end{minipage}\par\medskip

% QUESTION 3.3-M06
\noindent\begin{minipage}{\linewidth}

\subsubsection{3.3-M06 · MC}\label{m06-mc-9}

\emph{Source:
\href{../Lecture\%20Slides/3.3.\%20Deep\%20Learning\%20Fundamentals.pdf\#page=24}{PDF
pp.~24--26}.}

For one-hot targets, cross-entropy \(-\log p_{true}\) is smaller when:

\begin{itemize}
\tightlist
\item
  A. The parameter count is larger, regardless of predictions.
\item
  B. The true class receives a larger probability.
\item
  C. The true class probability is ignored.
\item
  D. The true class receives a smaller probability.
\end{itemize}

\begin{solution}[8mm]{3.3-M06}
% Write A, B, C, or D here. You may add a brief reason.

\end{solution}
\end{minipage}\par\medskip

% QUESTION 3.3-M07
\noindent\begin{minipage}{\linewidth}

\subsubsection{3.3-M07 · MC}\label{m07-mc-9}

\emph{Source:
\href{../Lecture\%20Slides/3.3.\%20Deep\%20Learning\%20Fundamentals.pdf\#page=23}{PDF
pp.~23, 25--26}.}

Multiclass hinge loss encourages the correct-class score to:

\begin{itemize}
\tightlist
\item
  A. Exceed incorrect-class scores by a specified margin.
\item
  B. Equal zero for every example.
\item
  C. Be the smallest score.
\item
  D. Match every incorrect-class score.
\end{itemize}

\begin{solution}[8mm]{3.3-M07}
% Write A, B, C, or D here. You may add a brief reason.

\end{solution}
\end{minipage}\par\medskip

% QUESTION 3.3-M08
\noindent\begin{minipage}{\linewidth}

\subsubsection{3.3-M08 · MC}\label{m08-mc-9}

\emph{Source:
\href{../Lecture\%20Slides/3.3.\%20Deep\%20Learning\%20Fundamentals.pdf\#page=29}{PDF
pp.~29--35}.}

Mini-batch SGD estimates a gradient using:

\begin{itemize}
\tightlist
\item
  A. Every possible input in the world.
\item
  B. A subset of training examples at each update.
\item
  C. Only the model's architecture.
\item
  D. Only validation labels.
\end{itemize}

\begin{solution}[8mm]{3.3-M08}
% Write A, B, C, or D here. You may add a brief reason.

\end{solution}
\end{minipage}\par\medskip

% QUESTION 3.3-M09
\noindent\begin{minipage}{\linewidth}

\subsubsection{3.3-M09 · MC}\label{m09-mc-9}

\emph{Source:
\href{../Lecture\%20Slides/3.3.\%20Deep\%20Learning\%20Fundamentals.pdf\#page=38}{PDF
pp.~38--39}.}

Momentum uses:

\begin{itemize}
\tightlist
\item
  A. A separate label for every parameter.
\item
  B. A running combination of recent gradients or updates.
\item
  C. Only the current test accuracy.
\item
  D. A fixed random direction forever.
\end{itemize}

\begin{solution}[8mm]{3.3-M09}
% Write A, B, C, or D here. You may add a brief reason.

\end{solution}
\end{minipage}\par\medskip

% QUESTION 3.3-M10
\noindent\begin{minipage}{\linewidth}

\subsubsection{3.3-M10 · MC}\label{m10-mc-9}

\emph{Source:
\href{../Lecture\%20Slides/3.3.\%20Deep\%20Learning\%20Fundamentals.pdf\#page=40}{PDF
pp.~40--42}.}

Adam combines momentum-like information with:

\begin{itemize}
\tightlist
\item
  A. Exact Hessian inversion at every step.
\item
  B. A rule that all parameters share identical gradients.
\item
  C. Per-parameter scaling based on squared-gradient statistics.
\item
  D. Removal of the learning rate.
\end{itemize}

\begin{solution}[8mm]{3.3-M10}
% Write A, B, C, or D here. You may add a brief reason.

\end{solution}
\end{minipage}\par\medskip

\subsection{Short answers}\label{short-answers-9}

% QUESTION 3.3-S01
\noindent\begin{minipage}{\linewidth}

\subsubsection{3.3-S01 · Short}\label{s01-short-9}

\emph{Source:
\href{../Lecture\%20Slides/3.3.\%20Deep\%20Learning\%20Fundamentals.pdf\#page=8}{PDF
pp.~8--14}.}

Why does universal approximation by shallow networks not make deeper
architectures useless?

\begin{solution}[20mm]{3.3-S01}
% Write your answer here, usually 1--3 sentences.

\end{solution}
\end{minipage}\par\medskip

% QUESTION 3.3-S02
\noindent\begin{minipage}{\linewidth}

\subsubsection{3.3-S02 · Short}\label{s02-short-9}

\emph{Source:
\href{../Lecture\%20Slides/3.3.\%20Deep\%20Learning\%20Fundamentals.pdf\#page=17}{PDF
pp.~17--27}.}

Distinguish a score function, a loss function, and an optimizer in one
phrase each.

\begin{solution}[20mm]{3.3-S02}
% Write your answer here, usually 1--3 sentences.

\end{solution}
\end{minipage}\par\medskip

% QUESTION 3.3-S03
\noindent\begin{minipage}{\linewidth}

\subsubsection{3.3-S03 · Short}\label{s03-short-9}

\emph{Source:
\href{../Lecture\%20Slides/3.3.\%20Deep\%20Learning\%20Fundamentals.pdf\#page=23}{PDF
pp.~23--26}.}

Why can two predictions have the same correct top-1 class but different
cross-entropy losses?

\begin{solution}[20mm]{3.3-S03}
% Write your answer here, usually 1--3 sentences.

\end{solution}
\end{minipage}\par\medskip

% QUESTION 3.3-S04
\noindent\begin{minipage}{\linewidth}

\subsubsection{3.3-S04 · Short}\label{s04-short-9}

\emph{Source:
\href{../Lecture\%20Slides/3.3.\%20Deep\%20Learning\%20Fundamentals.pdf\#page=27}{PDF
pp.~27--28}.}

Why should the output activation and loss match the task? Suggest a
suitable loss for continuous-target regression.

\begin{solution}[20mm]{3.3-S04}
% Write your answer here, usually 1--3 sentences.

\end{solution}
\end{minipage}\par\medskip

% QUESTION 3.3-S05
\noindent\begin{minipage}{\linewidth}

\subsubsection{3.3-S05 · Short}\label{s05-short-9}

\emph{Source:
\href{../Lecture\%20Slides/3.3.\%20Deep\%20Learning\%20Fundamentals.pdf\#page=32}{PDF
pp.~32--37}.}

What can happen if the learning rate is too high? What if it is
extremely low?

\begin{solution}[20mm]{3.3-S05}
% Write your answer here, usually 1--3 sentences.

\end{solution}
\end{minipage}\par\medskip

% QUESTION 3.3-S06
\noindent\begin{minipage}{\linewidth}

\subsubsection{3.3-S06 · Short}\label{s06-short-9}

\emph{Source:
\href{../Lecture\%20Slides/3.3.\%20Deep\%20Learning\%20Fundamentals.pdf\#page=32}{PDF
pp.~32--35}.}

Distinguish an epoch from a mini-batch update. Why can their counts
differ substantially?

\begin{solution}[20mm]{3.3-S06}
% Write your answer here, usually 1--3 sentences.

\end{solution}
\end{minipage}\par\medskip

% QUESTION 3.3-S07
\noindent\begin{minipage}{\linewidth}

\subsubsection{3.3-S07 · Short}\label{s07-short-9}

\emph{Source:
\href{../Lecture\%20Slides/3.3.\%20Deep\%20Learning\%20Fundamentals.pdf\#page=38}{PDF
pp.~38--44}.}

How can momentum reduce oscillation during optimization? Does using Adam
remove the need to tune hyperparameters?

\begin{solution}[20mm]{3.3-S07}
% Write your answer here, usually 1--3 sentences.

\end{solution}
\end{minipage}\par\medskip

\subsection{Easy calculations}\label{easy-calculations-9}

% QUESTION 3.3-C01
\noindent\begin{minipage}{\linewidth}

\subsubsection{3.3-C01 · Calculation}\label{c01-calculation-9}

\emph{Source:
\href{../Lecture\%20Slides/3.3.\%20Deep\%20Learning\%20Fundamentals.pdf\#page=8}{PDF
pp.~8--10}.}

A fully connected network has 3 inputs, one hidden layer of 4 units, and
2 outputs. Include one bias per hidden/output unit. Count all trainable
parameters.

\begin{solution}[26mm]{3.3-C01}
% Write your short calculation and final result here. Use $...$ or \[...\] for maths.

\end{solution}
\end{minipage}\par\medskip

% QUESTION 3.3-C02
\noindent\begin{minipage}{\linewidth}

\subsubsection{3.3-C02 · Calculation}\label{c02-calculation-9}

\emph{Source:
\href{../Lecture\%20Slides/3.3.\%20Deep\%20Learning\%20Fundamentals.pdf\#page=24}{PDF
pp.~24--26}.}

Three logits have exponentials \((1,2,1)\); class 2 is correct. Compute
its softmax probability and the cross-entropy. Use \(\ln2=0.693\).

\begin{solution}[26mm]{3.3-C02}
% Write your short calculation and final result here. Use $...$ or \[...\] for maths.

\end{solution}
\end{minipage}\par\medskip

% QUESTION 3.3-C03
\noindent\begin{minipage}{\linewidth}

\subsubsection{3.3-C03 · Calculation}\label{c03-calculation-9}

\emph{Source:
\href{../Lecture\%20Slides/3.3.\%20Deep\%20Learning\%20Fundamentals.pdf\#page=23}{PDF
pp.~23, 25}.}

For one example, the correct-class score is 2 and the two incorrect
scores are 1 and 4. With margin 1, compute
\(L=\sum_{j\ne y}\max(0,s_j-s_y+1)\).

\begin{solution}[26mm]{3.3-C03}
% Write your short calculation and final result here. Use $...$ or \[...\] for maths.

\end{solution}
\end{minipage}\par\medskip

\clearpage

\section{3.4. Deep Learning Fundamentals
II}\label{deep-learning-fundamentals-ii}

\textbf{Slides:}
\href{../Lecture\%20Slides/3.4.\%20Deep\%20Learning\%20Fundamentals\%20II.pdf}{3.4.
Deep Learning Fundamentals II.pdf}. \textbf{20 questions: 10 MC · 7
short answers · 3 easy calculations.}

\subsection{Multiple choice}\label{multiple-choice-10}

% QUESTION 3.4-M01
\noindent\begin{minipage}{\linewidth}

\subsubsection{3.4-M01 · MC}\label{m01-mc-10}

\emph{Source:
\href{../Lecture\%20Slides/3.4.\%20Deep\%20Learning\%20Fundamentals\%20II.pdf\#page=7}{PDF
pp.~7--21}.}

Backpropagation computes:

\begin{itemize}
\tightlist
\item
  A. Test labels from training labels.
\item
  B. The optimal architecture without training.
\item
  C. Loss gradients with respect to parameters using the chain rule.
\item
  D. Parameter updates without an optimizer.
\end{itemize}

\begin{solution}[8mm]{3.4-M01}
% Write A, B, C, or D here. You may add a brief reason.

\end{solution}
\end{minipage}\par\medskip

% QUESTION 3.4-M02
\noindent\begin{minipage}{\linewidth}

\subsubsection{3.4-M02 · MC}\label{m02-mc-10}

\emph{Source:
\href{../Lecture\%20Slides/3.4.\%20Deep\%20Learning\%20Fundamentals\%20II.pdf\#page=7}{PDF
pp.~7, 19--21}.}

Why store intermediate activations during the forward pass?

\begin{itemize}
\tightlist
\item
  A. They can be needed to compute local derivatives during the backward
  pass.
\item
  B. They become new permanent parameters.
\item
  C. They guarantee zero training error.
\item
  D. They replace all training labels.
\end{itemize}

\begin{solution}[8mm]{3.4-M02}
% Write A, B, C, or D here. You may add a brief reason.

\end{solution}
\end{minipage}\par\medskip

% QUESTION 3.4-M03
\noindent\begin{minipage}{\linewidth}

\subsubsection{3.4-M03 · MC}\label{m03-mc-10}

\emph{Source:
\href{../Lecture\%20Slides/3.4.\%20Deep\%20Learning\%20Fundamentals\%20II.pdf\#page=10}{PDF
pp.~10--16}.}

If a parameter influences loss through two computational paths, its
gradient receives:

\begin{itemize}
\tightlist
\item
  A. Only the larger contribution.
\item
  B. Only the first path's contribution.
\item
  C. The sum of contributions from both paths.
\item
  D. Their average regardless of the computation.
\end{itemize}

\begin{solution}[8mm]{3.4-M03}
% Write A, B, C, or D here. You may add a brief reason.

\end{solution}
\end{minipage}\par\medskip

% QUESTION 3.4-M04
\noindent\begin{minipage}{\linewidth}

\subsubsection{3.4-M04 · MC}\label{m04-mc-10}

\emph{Source:
\href{../Lecture\%20Slides/3.4.\%20Deep\%20Learning\%20Fundamentals\%20II.pdf\#page=23}{PDF
pp.~23--25}.}

Why avoid identical zero initialization of hidden units in one layer?

\begin{itemize}
\tightlist
\item
  A. Zero is an invalid parameter value.
\item
  B. Biases must always be positive.
\item
  C. Units can remain symmetric and learn the same features.
\item
  D. The forward pass becomes undefined.
\end{itemize}

\begin{solution}[8mm]{3.4-M04}
% Write A, B, C, or D here. You may add a brief reason.

\end{solution}
\end{minipage}\par\medskip

% QUESTION 3.4-M05
\noindent\begin{minipage}{\linewidth}

\subsubsection{3.4-M05 · MC}\label{m05-mc-10}

\emph{Source:
\href{../Lecture\%20Slides/3.4.\%20Deep\%20Learning\%20Fundamentals\%20II.pdf\#page=24}{PDF
pp.~24--25}.}

Initialization schemes such as He initialization are intended partly to:

\begin{itemize}
\tightlist
\item
  A. Make all neurons identical.
\item
  B. Control activation and gradient scales across layers.
\item
  C. Fit test labels before training.
\item
  D. Eliminate the need for nonlinearities.
\end{itemize}

\begin{solution}[8mm]{3.4-M05}
% Write A, B, C, or D here. You may add a brief reason.

\end{solution}
\end{minipage}\par\medskip

% QUESTION 3.4-M06
\noindent\begin{minipage}{\linewidth}

\subsubsection{3.4-M06 · MC}\label{m06-mc-10}

\emph{Source:
\href{../Lecture\%20Slides/3.4.\%20Deep\%20Learning\%20Fundamentals\%20II.pdf\#page=27}{PDF
pp.~27--29}.}

An explicit L2 regularizer directly encourages:

\begin{itemize}
\tightlist
\item
  A. Higher output entropy in every case.
\item
  B. Larger training batches.
\item
  C. Smaller weight magnitudes.
\item
  D. More hidden layers.
\end{itemize}

\begin{solution}[8mm]{3.4-M06}
% Write A, B, C, or D here. You may add a brief reason.

\end{solution}
\end{minipage}\par\medskip

% QUESTION 3.4-M07
\noindent\begin{minipage}{\linewidth}

\subsubsection{3.4-M07 · MC}\label{m07-mc-10}

\emph{Source:
\href{../Lecture\%20Slides/3.4.\%20Deep\%20Learning\%20Fundamentals\%20II.pdf\#page=31}{PDF
pp.~31}.}

Early stopping normally selects a checkpoint using:

\begin{itemize}
\tightlist
\item
  A. The smallest training loss regardless of validation.
\item
  B. Repeated test-set optimization.
\item
  C. Held-out validation performance.
\item
  D. Only the wall-clock time of an epoch.
\end{itemize}

\begin{solution}[8mm]{3.4-M07}
% Write A, B, C, or D here. You may add a brief reason.

\end{solution}
\end{minipage}\par\medskip

% QUESTION 3.4-M08
\noindent\begin{minipage}{\linewidth}

\subsubsection{3.4-M08 · MC}\label{m08-mc-10}

\emph{Source:
\href{../Lecture\%20Slides/3.4.\%20Deep\%20Learning\%20Fundamentals\%20II.pdf\#page=32}{PDF
pp.~32, 34}.}

Standard dropout during training:

\begin{itemize}
\tightlist
\item
  A. Drops only incorrect training examples.
\item
  B. Is normally applied unchanged during ordinary inference.
\item
  C. Randomly suppresses some activations and uses an appropriate
  scaling convention.
\item
  D. Permanently deletes a fixed subset of units.
\end{itemize}

\begin{solution}[8mm]{3.4-M08}
% Write A, B, C, or D here. You may add a brief reason.

\end{solution}
\end{minipage}\par\medskip

% QUESTION 3.4-M09
\noindent\begin{minipage}{\linewidth}

\subsubsection{3.4-M09 · MC}\label{m09-mc-10}

\emph{Source:
\href{../Lecture\%20Slides/3.4.\%20Deep\%20Learning\%20Fundamentals\%20II.pdf\#page=33}{PDF
pp.~33--34}.}

A suitable supervised data augmentation should usually:

\begin{itemize}
\tightlist
\item
  A. Make every training example identical.
\item
  B. Preserve the intended target meaning while varying the input.
\item
  C. Change the true class arbitrarily.
\item
  D. Use test labels to select transformations.
\end{itemize}

\begin{solution}[8mm]{3.4-M09}
% Write A, B, C, or D here. You may add a brief reason.

\end{solution}
\end{minipage}\par\medskip

% QUESTION 3.4-M10
\noindent\begin{minipage}{\linewidth}

\subsubsection{3.4-M10 · MC}\label{m10-mc-10}

\emph{Source:
\href{../Lecture\%20Slides/3.4.\%20Deep\%20Learning\%20Fundamentals\%20II.pdf\#page=37}{PDF
pp.~37--38}.}

Double descent describes a possible pattern where test error:

\begin{itemize}
\tightlist
\item
  A. Is independent of model capacity.
\item
  B. Always decreases monotonically with model capacity.
\item
  C. Rises near an interpolation threshold and then falls again with
  further capacity.
\item
  D. Equals training error throughout training.
\end{itemize}

\begin{solution}[8mm]{3.4-M10}
% Write A, B, C, or D here. You may add a brief reason.

\end{solution}
\end{minipage}\par\medskip

\subsection{Short answers}\label{short-answers-10}

% QUESTION 3.4-S01
\noindent\begin{minipage}{\linewidth}

\subsubsection{3.4-S01 · Short}\label{s01-short-10}

\emph{Source:
\href{../Lecture\%20Slides/3.4.\%20Deep\%20Learning\%20Fundamentals\%20II.pdf\#page=7}{PDF
pp.~7--22}.}

Distinguish the forward pass, backward pass, and optimizer step. Which
computes predictions and which changes weights?

\begin{solution}[20mm]{3.4-S01}
% Write your answer here, usually 1--3 sentences.

\end{solution}
\end{minipage}\par\medskip

% QUESTION 3.4-S02
\noindent\begin{minipage}{\linewidth}

\subsubsection{3.4-S02 · Short}\label{s02-short-10}

\emph{Source:
\href{../Lecture\%20Slides/3.4.\%20Deep\%20Learning\%20Fundamentals\%20II.pdf\#page=10}{PDF
pp.~10--16}.}

Why does the chain rule make deep-network gradients tractable without
expanding the entire network into one enormous formula?

\begin{solution}[20mm]{3.4-S02}
% Write your answer here, usually 1--3 sentences.

\end{solution}
\end{minipage}\par\medskip

% QUESTION 3.4-S03
\noindent\begin{minipage}{\linewidth}

\subsubsection{3.4-S03 · Short}\label{s03-short-10}

\emph{Source:
\href{../Lecture\%20Slides/3.4.\%20Deep\%20Learning\%20Fundamentals\%20II.pdf\#page=24}{PDF
pp.~24--25}.}

What do vanishing and exploding gradients mean for learning in early
layers?

\begin{solution}[20mm]{3.4-S03}
% Write your answer here, usually 1--3 sentences.

\end{solution}
\end{minipage}\par\medskip

% QUESTION 3.4-S04
\noindent\begin{minipage}{\linewidth}

\subsubsection{3.4-S04 · Short}\label{s04-short-10}

\emph{Source:
\href{../Lecture\%20Slides/3.4.\%20Deep\%20Learning\%20Fundamentals\%20II.pdf\#page=26}{PDF
pp.~26--31}.}

Distinguish explicit regularization from an implicit regularizing effect
of the training procedure. Give one example of each.

\begin{solution}[20mm]{3.4-S04}
% Write your answer here, usually 1--3 sentences.

\end{solution}
\end{minipage}\par\medskip

% QUESTION 3.4-S05
\noindent\begin{minipage}{\linewidth}

\subsubsection{3.4-S05 · Short}\label{s05-short-10}

\emph{Source:
\href{../Lecture\%20Slides/3.4.\%20Deep\%20Learning\%20Fundamentals\%20II.pdf\#page=32}{PDF
pp.~32--34}.}

Why is dropout normally disabled during ordinary inference? What would
happen to repeatability if random dropout masks remained active?

\begin{solution}[20mm]{3.4-S05}
% Write your answer here, usually 1--3 sentences.

\end{solution}
\end{minipage}\par\medskip

% QUESTION 3.4-S06
\noindent\begin{minipage}{\linewidth}

\subsubsection{3.4-S06 · Short}\label{s06-short-10}

\emph{Source:
\href{../Lecture\%20Slides/3.4.\%20Deep\%20Learning\%20Fundamentals\%20II.pdf\#page=33}{PDF
pp.~33--34}.}

Give one label-preserving augmentation for an object-recognition task
and one transformation that might change the label in a
digit-recognition task.

\begin{solution}[20mm]{3.4-S06}
% Write your answer here, usually 1--3 sentences.

\end{solution}
\end{minipage}\par\medskip

% QUESTION 3.4-S07
\noindent\begin{minipage}{\linewidth}

\subsubsection{3.4-S07 · Short}\label{s07-short-10}

\emph{Source:
\href{../Lecture\%20Slides/3.4.\%20Deep\%20Learning\%20Fundamentals\%20II.pdf\#page=26}{PDF
pp.~26, 31, 35--36}.}

Training loss decreases while validation loss rises. What does this
suggest, and how would early stopping use these curves?

\begin{solution}[20mm]{3.4-S07}
% Write your answer here, usually 1--3 sentences.

\end{solution}
\end{minipage}\par\medskip

\subsection{Easy calculations}\label{easy-calculations-10}

% QUESTION 3.4-C01
\noindent\begin{minipage}{\linewidth}

\subsubsection{3.4-C01 · Calculation}\label{c01-calculation-10}

\emph{Source:
\href{../Lecture\%20Slides/3.4.\%20Deep\%20Learning\%20Fundamentals\%20II.pdf\#page=10}{PDF
pp.~10--16}.}

At a point in a graph, \(\partial L/\partial h=3\) and
\(\partial h/\partial w=2\). There is no other path from \(w\) to \(L\).
Compute \(\partial L/\partial w\).

\begin{solution}[26mm]{3.4-C01}
% Write your short calculation and final result here. Use $...$ or \[...\] for maths.

\end{solution}
\end{minipage}\par\medskip

% QUESTION 3.4-C02
\noindent\begin{minipage}{\linewidth}

\subsubsection{3.4-C02 · Calculation}\label{c02-calculation-10}

\emph{Source:
\href{../Lecture\%20Slides/3.4.\%20Deep\%20Learning\%20Fundamentals\%20II.pdf\#page=16}{PDF
pp.~16, 20--21}.}

A weight is \(w=2\), its gradient is 3, and the learning rate is 0.1.
Compute one gradient-descent update.

\begin{solution}[26mm]{3.4-C02}
% Write your short calculation and final result here. Use $...$ or \[...\] for maths.

\end{solution}
\end{minipage}\par\medskip

% QUESTION 3.4-C03
\noindent\begin{minipage}{\linewidth}

\subsubsection{3.4-C03 · Calculation}\label{c03-calculation-10}

\emph{Source:
\href{../Lecture\%20Slides/3.4.\%20Deep\%20Learning\%20Fundamentals\%20II.pdf\#page=27}{PDF
pp.~27--29}.}

For \(J=L+\lambda w^2\), the data-loss gradient is
\(\partial L/\partial w=3\), \(w=2\), and \(\lambda=0.1\). Using
\(\partial J/\partial w=\partial L/\partial w+2\lambda w\), compute the
total gradient.

\begin{solution}[26mm]{3.4-C03}
% Write your short calculation and final result here. Use $...$ or \[...\] for maths.

\end{solution}
\end{minipage}\par\medskip

\clearpage

\section{4.1. CNNs}\label{cnns}

\textbf{Slides:} \href{../Lecture\%20Slides/4.1.\%20CNNs.pdf}{4.1.
CNNs.pdf}. \textbf{20 questions: 10 MC · 7 short answers · 3 easy
calculations.}

\subsection{Multiple choice}\label{multiple-choice-11}

% QUESTION 4.1-M01
\noindent\begin{minipage}{\linewidth}

\subsubsection{4.1-M01 · MC}\label{m01-mc-11}

\emph{Source: \href{../Lecture\%20Slides/4.1.\%20CNNs.pdf\#page=10}{PDF
pp.~10--11}.}

Which properties reduce a convolutional layer's parameter count relative
to dense image connections?

\begin{itemize}
\tightlist
\item
  A. A separate filter for every pixel.
\item
  B. Local connectivity and spatial weight sharing.
\item
  C. Global connectivity without sharing.
\item
  D. More image pixels and more class labels.
\end{itemize}

\begin{solution}[8mm]{4.1-M01}
% Write A, B, C, or D here. You may add a brief reason.

\end{solution}
\end{minipage}\par\medskip

% QUESTION 4.1-M02
\noindent\begin{minipage}{\linewidth}

\subsubsection{4.1-M02 · MC}\label{m02-mc-11}

\emph{Source: \href{../Lecture\%20Slides/4.1.\%20CNNs.pdf\#page=19}{PDF
pp.~19, 29--30}.}

In a standard convolution, the output-channel count equals:

\begin{itemize}
\tightlist
\item
  A. The number of filters.
\item
  B. The number of spatial output locations.
\item
  C. The stride times the input width.
\item
  D. The kernel width.
\end{itemize}

\begin{solution}[8mm]{4.1-M02}
% Write A, B, C, or D here. You may add a brief reason.

\end{solution}
\end{minipage}\par\medskip

% QUESTION 4.1-M03
\noindent\begin{minipage}{\linewidth}

\subsubsection{4.1-M03 · MC}\label{m03-mc-11}

\emph{Source: \href{../Lecture\%20Slides/4.1.\%20CNNs.pdf\#page=15}{PDF
pp.~15--17}.}

Increasing convolution stride, with other settings fixed, generally:

\begin{itemize}
\tightlist
\item
  A. Increases spatial output resolution.
\item
  B. Reduces spatial output resolution.
\item
  C. Adds learned filter weights.
\item
  D. Makes padding unnecessary for every input.
\end{itemize}

\begin{solution}[8mm]{4.1-M03}
% Write A, B, C, or D here. You may add a brief reason.

\end{solution}
\end{minipage}\par\medskip

% QUESTION 4.1-M04
\noindent\begin{minipage}{\linewidth}

\subsubsection{4.1-M04 · MC}\label{m04-mc-11}

\emph{Source: \href{../Lecture\%20Slides/4.1.\%20CNNs.pdf\#page=15}{PDF
pp.~15--16, 20--23}.}

Increasing dilation while keeping the kernel's coefficient count fixed:

\begin{itemize}
\tightlist
\item
  A. Spreads sampled positions over a wider spatial region.
\item
  B. Makes all coefficients identical.
\item
  C. Changes only the output-channel count.
\item
  D. Adds a new coefficient at every skipped location.
\end{itemize}

\begin{solution}[8mm]{4.1-M04}
% Write A, B, C, or D here. You may add a brief reason.

\end{solution}
\end{minipage}\par\medskip

% QUESTION 4.1-M05
\noindent\begin{minipage}{\linewidth}

\subsubsection{4.1-M05 · MC}\label{m05-mc-11}

\emph{Source: \href{../Lecture\%20Slides/4.1.\%20CNNs.pdf\#page=32}{PDF
pp.~32--35}.}

An ordinary max-pooling layer:

\begin{itemize}
\tightlist
\item
  A. Always increases image resolution.
\item
  B. Selects local maxima and has no learned filter weights.
\item
  C. Learns a separate weight for every input pixel.
\item
  D. Computes softmax class probabilities.
\end{itemize}

\begin{solution}[8mm]{4.1-M05}
% Write A, B, C, or D here. You may add a brief reason.

\end{solution}
\end{minipage}\par\medskip

% QUESTION 4.1-M06
\noindent\begin{minipage}{\linewidth}

\subsubsection{4.1-M06 · MC}\label{m06-mc-11}

\emph{Source: \href{../Lecture\%20Slides/4.1.\%20CNNs.pdf\#page=36}{PDF
pp.~36}.}

A stride-1 \(1\times1\) convolution primarily:

\begin{itemize}
\tightlist
\item
  A. Expands the spatial receptive field to \(3\times3\).
\item
  B. Mixes channels at each location and can change channel count.
\item
  C. Removes all channel information.
\item
  D. Pools neighbouring pixels.
\end{itemize}

\begin{solution}[8mm]{4.1-M06}
% Write A, B, C, or D here. You may add a brief reason.

\end{solution}
\end{minipage}\par\medskip

% QUESTION 4.1-M07
\noindent\begin{minipage}{\linewidth}

\subsubsection{4.1-M07 · MC}\label{m07-mc-11}

\emph{Source: \href{../Lecture\%20Slides/4.1.\%20CNNs.pdf\#page=52}{PDF
pp.~52--54}.}

In \(y=x+F(x)\), setting \(F(x)=0\) gives:

\begin{itemize}
\tightlist
\item
  A. Twice the input.
\item
  B. The identity mapping.
\item
  C. Zero output.
\item
  D. Half the input channels.
\end{itemize}

\begin{solution}[8mm]{4.1-M07}
% Write A, B, C, or D here. You may add a brief reason.

\end{solution}
\end{minipage}\par\medskip

% QUESTION 4.1-M08
\noindent\begin{minipage}{\linewidth}

\subsubsection{4.1-M08 · MC}\label{m08-mc-11}

\emph{Source: \href{../Lecture\%20Slides/4.1.\%20CNNs.pdf\#page=55}{PDF
pp.~55--59}.}

Standard batch normalization with running statistics normally uses:

\begin{itemize}
\tightlist
\item
  A. The current evaluation batch for every deployment prediction.
\item
  B. No learned scale or shift.
\item
  C. Test-set statistics to fit its training parameters.
\item
  D. Batch statistics during training and stored statistics during
  evaluation.
\end{itemize}

\begin{solution}[8mm]{4.1-M08}
% Write A, B, C, or D here. You may add a brief reason.

\end{solution}
\end{minipage}\par\medskip

% QUESTION 4.1-M09
\noindent\begin{minipage}{\linewidth}

\subsubsection{4.1-M09 · MC}\label{m09-mc-11}

\emph{Source: \href{../Lecture\%20Slides/4.1.\%20CNNs.pdf\#page=67}{PDF
pp.~67--73}.}

For a small dataset similar to pretraining data, a sensible
transfer-learning starting point is:

\begin{itemize}
\tightlist
\item
  A. Keep the old labels even when the task changes.
\item
  B. Tune on the final test set.
\item
  C. Discard all pretrained weights immediately.
\item
  D. Freeze pretrained features and train a new task head.
\end{itemize}

\begin{solution}[8mm]{4.1-M09}
% Write A, B, C, or D here. You may add a brief reason.

\end{solution}
\end{minipage}\par\medskip

% QUESTION 4.1-M10
\noindent\begin{minipage}{\linewidth}

\subsubsection{4.1-M10 · MC}\label{m10-mc-11}

\emph{Source: \href{../Lecture\%20Slides/4.1.\%20CNNs.pdf\#page=90}{PDF
pp.~90--93}.}

Instance segmentation differs from semantic segmentation by:

\begin{itemize}
\tightlist
\item
  A. Assigning no class labels to objects.
\item
  B. Avoiding pixel-level outputs.
\item
  C. Predicting only one class for the whole image.
\item
  D. Separating distinct objects even when they have the same class.
\end{itemize}

\begin{solution}[8mm]{4.1-M10}
% Write A, B, C, or D here. You may add a brief reason.

\end{solution}
\end{minipage}\par\medskip

\subsection{Short answers}\label{short-answers-11}

% QUESTION 4.1-S01
\noindent\begin{minipage}{\linewidth}

\subsubsection{4.1-S01 · Short}\label{s01-short-11}

\emph{Source: \href{../Lecture\%20Slides/4.1.\%20CNNs.pdf\#page=20}{PDF
pp.~20--23}.}

Define a unit's receptive field. Why can deeper convolutional layers use
information from a larger input region?

\begin{solution}[20mm]{4.1-S01}
% Write your answer here, usually 1--3 sentences.

\end{solution}
\end{minipage}\par\medskip

% QUESTION 4.1-S02
\noindent\begin{minipage}{\linewidth}

\subsubsection{4.1-S02 · Short}\label{s02-short-11}

\emph{Source: \href{../Lecture\%20Slides/4.1.\%20CNNs.pdf\#page=52}{PDF
pp.~52--54, 61}.}

What shape condition is needed to add a residual shortcut to a block
output? How can a projection shortcut help when shapes differ?

\begin{solution}[20mm]{4.1-S02}
% Write your answer here, usually 1--3 sentences.

\end{solution}
\end{minipage}\par\medskip

% QUESTION 4.1-S03
\noindent\begin{minipage}{\linewidth}

\subsubsection{4.1-S03 · Short}\label{s03-short-11}

\emph{Source: \href{../Lecture\%20Slides/4.1.\%20CNNs.pdf\#page=55}{PDF
pp.~55--59}.}

Contrast batch normalization and layer normalization in terms of which
values contribute to their normalization statistics.

\begin{solution}[20mm]{4.1-S03}
% Write your answer here, usually 1--3 sentences.

\end{solution}
\end{minipage}\par\medskip

% QUESTION 4.1-S04
\noindent\begin{minipage}{\linewidth}

\subsubsection{4.1-S04 · Short}\label{s04-short-11}

\emph{Source: \href{../Lecture\%20Slides/4.1.\%20CNNs.pdf\#page=64}{PDF
pp.~64}.}

What are the two operations in a depthwise-separable convolution? Which
operation mixes channels?

\begin{solution}[20mm]{4.1-S04}
% Write your answer here, usually 1--3 sentences.

\end{solution}
\end{minipage}\par\medskip

% QUESTION 4.1-S05
\noindent\begin{minipage}{\linewidth}

\subsubsection{4.1-S05 · Short}\label{s05-short-11}

\emph{Source: \href{../Lecture\%20Slides/4.1.\%20CNNs.pdf\#page=67}{PDF
pp.~67--73}.}

Distinguish frozen feature extraction from fine-tuning. Which weights
change in each case?

\begin{solution}[20mm]{4.1-S05}
% Write your answer here, usually 1--3 sentences.

\end{solution}
\end{minipage}\par\medskip

% QUESTION 4.1-S06
\noindent\begin{minipage}{\linewidth}

\subsubsection{4.1-S06 · Short}\label{s06-short-11}

\emph{Source: \href{../Lecture\%20Slides/4.1.\%20CNNs.pdf\#page=51}{PDF
pp.~51--62, 75--79}.}

A deeper plain CNN has higher training error than a shallower one. Why
is this alone not evidence of overfitting? How can residual connections
help?

\begin{solution}[20mm]{4.1-S06}
% Write your answer here, usually 1--3 sentences.

\end{solution}
\end{minipage}\par\medskip

% QUESTION 4.1-S07
\noindent\begin{minipage}{\linewidth}

\subsubsection{4.1-S07 · Short}\label{s07-short-11}

\emph{Source: \href{../Lecture\%20Slides/4.1.\%20CNNs.pdf\#page=83}{PDF
pp.~83--93}.}

Distinguish the expected outputs of image classification, object
detection, and semantic segmentation.

\begin{solution}[20mm]{4.1-S07}
% Write your answer here, usually 1--3 sentences.

\end{solution}
\end{minipage}\par\medskip

\subsection{Easy calculations}\label{easy-calculations-11}

% QUESTION 4.1-C01
\noindent\begin{minipage}{\linewidth}

\subsubsection{4.1-C01 · Calculation}\label{c01-calculation-11}

\emph{Source: \href{../Lecture\%20Slides/4.1.\%20CNNs.pdf\#page=15}{PDF
pp.~15--17}.}

A \(7\times7\) input is convolved with a \(3\times3\) kernel, stride 2,
padding 1 per side, and dilation 1. Use
\(W_{out}=\lfloor(W-K+2P)/S\rfloor+1\) to compute output height and
width.

\begin{solution}[26mm]{4.1-C01}
% Write your short calculation and final result here. Use $...$ or \[...\] for maths.

\end{solution}
\end{minipage}\par\medskip

% QUESTION 4.1-C02
\noindent\begin{minipage}{\linewidth}

\subsubsection{4.1-C02 · Calculation}\label{c02-calculation-11}

\emph{Source: \href{../Lecture\%20Slides/4.1.\%20CNNs.pdf\#page=29}{PDF
pp.~29--30}.}

A standard \(3\times3\) convolution has 2 input channels and 4 output
channels, with one bias per output channel. Count the trainable scalar
parameters.

\begin{solution}[26mm]{4.1-C02}
% Write your short calculation and final result here. Use $...$ or \[...\] for maths.

\end{solution}
\end{minipage}\par\medskip

% QUESTION 4.1-C03
\noindent\begin{minipage}{\linewidth}

\subsubsection{4.1-C03 · Calculation}\label{c03-calculation-11}

\emph{Source: \href{../Lecture\%20Slides/4.1.\%20CNNs.pdf\#page=20}{PDF
pp.~20--23}.}

Two successive \(3\times3\) convolutions both have stride 1 and dilation
1. For an interior output position, what is the receptive-field height
and width in the original input?

\begin{solution}[26mm]{4.1-C03}
% Write your short calculation and final result here. Use $...$ or \[...\] for maths.

\end{solution}
\end{minipage}\par\medskip

\clearpage

\section{5.1. Sequence Models}\label{sequence-models}

\textbf{Slides:}
\href{../Lecture\%20Slides/5.1.\%20Sequence\%20Models.pdf}{5.1. Sequence
Models.pdf}. \textbf{20 questions: 10 MC · 7 short answers · 3 easy
calculations.}

\subsection{Multiple choice}\label{multiple-choice-12}

% QUESTION 5.1-M01
\noindent\begin{minipage}{\linewidth}

\subsubsection{5.1-M01 · MC}\label{m01-mc-12}

\emph{Source:
\href{../Lecture\%20Slides/5.1.\%20Sequence\%20Models.pdf\#page=4}{PDF
pp.~4--6}.}

What makes sequential data different from an unordered collection?

\begin{itemize}
\tightlist
\item
  A. Temporal dependencies are always absent.
\item
  B. Ordering can change the meaning or prediction task.
\item
  C. Sequence values must be categorical.
\item
  D. Every sequence has the same length.
\end{itemize}

\begin{solution}[8mm]{5.1-M01}
% Write A, B, C, or D here. You may add a brief reason.

\end{solution}
\end{minipage}\par\medskip

% QUESTION 5.1-M02
\noindent\begin{minipage}{\linewidth}

\subsubsection{5.1-M02 · MC}\label{m02-mc-12}

\emph{Source:
\href{../Lecture\%20Slides/5.1.\%20Sequence\%20Models.pdf\#page=6}{PDF
pp.~6--9}.}

As a trained RNN processes a sequence at inference time:

\begin{itemize}
\tightlist
\item
  A. It learns independent weights for each position.
\item
  B. Its hidden state stays constant while weights change.
\item
  C. Its hidden state changes while learned weights are shared across
  steps.
\item
  D. Its hidden dimension grows with every token.
\end{itemize}

\begin{solution}[8mm]{5.1-M02}
% Write A, B, C, or D here. You may add a brief reason.

\end{solution}
\end{minipage}\par\medskip

% QUESTION 5.1-M03
\noindent\begin{minipage}{\linewidth}

\subsubsection{5.1-M03 · MC}\label{m03-mc-12}

\emph{Source:
\href{../Lecture\%20Slides/5.1.\%20Sequence\%20Models.pdf\#page=7}{PDF
pp.~7--9}.}

The recurrent term involving \(h_{t-1}\) allows the current computation
to depend on:

\begin{itemize}
\tightlist
\item
  A. Only the current input.
\item
  B. Only the future target label.
\item
  C. Information carried from previous steps.
\item
  D. A separate parameter file for every sequence.
\end{itemize}

\begin{solution}[8mm]{5.1-M03}
% Write A, B, C, or D here. You may add a brief reason.

\end{solution}
\end{minipage}\par\medskip

% QUESTION 5.1-M04
\noindent\begin{minipage}{\linewidth}

\subsubsection{5.1-M04 · MC}\label{m04-mc-12}

\emph{Source:
\href{../Lecture\%20Slides/5.1.\%20Sequence\%20Models.pdf\#page=13}{PDF
pp.~13--15}.}

Vanishing gradients primarily make it difficult to:

\begin{itemize}
\tightlist
\item
  A. Learn dependencies spanning many recurrent steps.
\item
  B. Calculate the number of tokens.
\item
  C. Store the output vocabulary.
\item
  D. Evaluate the current input's dimensions.
\end{itemize}

\begin{solution}[8mm]{5.1-M04}
% Write A, B, C, or D here. You may add a brief reason.

\end{solution}
\end{minipage}\par\medskip

% QUESTION 5.1-M05
\noindent\begin{minipage}{\linewidth}

\subsubsection{5.1-M05 · MC}\label{m05-mc-12}

\emph{Source:
\href{../Lecture\%20Slides/5.1.\%20Sequence\%20Models.pdf\#page=16}{PDF
pp.~16}.}

Norm-based gradient clipping directly addresses:

\begin{itemize}
\tightlist
\item
  A. Every source of vanishing gradients.
\item
  B. Excessively large gradient norms.
\item
  C. Missing target labels.
\item
  D. Insufficient vocabulary size.
\end{itemize}

\begin{solution}[8mm]{5.1-M05}
% Write A, B, C, or D here. You may add a brief reason.

\end{solution}
\end{minipage}\par\medskip

% QUESTION 5.1-M06
\noindent\begin{minipage}{\linewidth}

\subsubsection{5.1-M06 · MC}\label{m06-mc-12}

\emph{Source:
\href{../Lecture\%20Slides/5.1.\%20Sequence\%20Models.pdf\#page=17}{PDF
pp.~17--21}.}

LSTM gating and its cell-state path help by:

\begin{itemize}
\tightlist
\item
  A. Making every hidden state equal to the input.
\item
  B. Requiring an entire sequence to have one label.
\item
  C. Controlling retention and updates through an additive memory path.
\item
  D. Removing all trainable recurrent weights.
\end{itemize}

\begin{solution}[8mm]{5.1-M06}
% Write A, B, C, or D here. You may add a brief reason.

\end{solution}
\end{minipage}\par\medskip

% QUESTION 5.1-M07
\noindent\begin{minipage}{\linewidth}

\subsubsection{5.1-M07 · MC}\label{m07-mc-12}

\emph{Source:
\href{../Lecture\%20Slides/5.1.\%20Sequence\%20Models.pdf\#page=21}{PDF
pp.~21}.}

Relative to an LSTM, the usual GRU:

\begin{itemize}
\tightlist
\item
  A. Has an additional separate cell state.
\item
  B. Cannot process variable-length sequences.
\item
  C. Has update/reset gates and no separate LSTM-style cell state.
\item
  D. Uses no gates.
\end{itemize}

\begin{solution}[8mm]{5.1-M07}
% Write A, B, C, or D here. You may add a brief reason.

\end{solution}
\end{minipage}\par\medskip

% QUESTION 5.1-M08
\noindent\begin{minipage}{\linewidth}

\subsubsection{5.1-M08 · MC}\label{m08-mc-12}

\emph{Source:
\href{../Lecture\%20Slides/5.1.\%20Sequence\%20Models.pdf\#page=22}{PDF
pp.~22}.}

A bidirectional RNN typically needs:

\begin{itemize}
\tightlist
\item
  A. Only previously observed inputs.
\item
  B. A causal mask for its forward branch only.
\item
  C. No recurrent state.
\item
  D. Access to the sequence in both directions.
\end{itemize}

\begin{solution}[8mm]{5.1-M08}
% Write A, B, C, or D here. You may add a brief reason.

\end{solution}
\end{minipage}\par\medskip

% QUESTION 5.1-M09
\noindent\begin{minipage}{\linewidth}

\subsubsection{5.1-M09 · MC}\label{m09-mc-12}

\emph{Source:
\href{../Lecture\%20Slides/5.1.\%20Sequence\%20Models.pdf\#page=28}{PDF
pp.~28--29}.}

Teacher forcing feeds the decoder:

\begin{itemize}
\tightlist
\item
  A. The final target sentence as a single class label.
\item
  B. Only its own mistakes during training.
\item
  C. The ground-truth next token at inference time.
\item
  D. The ground-truth previous token during training.
\end{itemize}

\begin{solution}[8mm]{5.1-M09}
% Write A, B, C, or D here. You may add a brief reason.

\end{solution}
\end{minipage}\par\medskip

% QUESTION 5.1-M10
\noindent\begin{minipage}{\linewidth}

\subsubsection{5.1-M10 · MC}\label{m10-mc-12}

\emph{Source:
\href{../Lecture\%20Slides/5.1.\%20Sequence\%20Models.pdf\#page=30}{PDF
pp.~30}.}

Attention addresses the basic seq2seq bottleneck by:

\begin{itemize}
\tightlist
\item
  A. Increasing the vocabulary indefinitely.
\item
  B. Removing the input sequence.
\item
  C. Forcing every decoder step to use identical context weights.
\item
  D. Allowing output steps to access different weighted combinations of
  encoder states.
\end{itemize}

\begin{solution}[8mm]{5.1-M10}
% Write A, B, C, or D here. You may add a brief reason.

\end{solution}
\end{minipage}\par\medskip

\subsection{Short answers}\label{short-answers-12}

% QUESTION 5.1-S01
\noindent\begin{minipage}{\linewidth}

\subsubsection{5.1-S01 · Short}\label{s01-short-12}

\emph{Source:
\href{../Lecture\%20Slides/5.1.\%20Sequence\%20Models.pdf\#page=6}{PDF
pp.~6--9}.}

Why does processing a longer sequence not by itself increase the number
of trainable RNN parameters?

\begin{solution}[20mm]{5.1-S01}
% Write your answer here, usually 1--3 sentences.

\end{solution}
\end{minipage}\par\medskip

% QUESTION 5.1-S02
\noindent\begin{minipage}{\linewidth}

\subsubsection{5.1-S02 · Short}\label{s02-short-12}

\emph{Source:
\href{../Lecture\%20Slides/5.1.\%20Sequence\%20Models.pdf\#page=14}{PDF
pp.~14}.}

What is backpropagation through time? What does truncating it trade away
to save computation and memory?

\begin{solution}[20mm]{5.1-S02}
% Write your answer here, usually 1--3 sentences.

\end{solution}
\end{minipage}\par\medskip

% QUESTION 5.1-S03
\noindent\begin{minipage}{\linewidth}

\subsubsection{5.1-S03 · Short}\label{s03-short-12}

\emph{Source:
\href{../Lecture\%20Slides/5.1.\%20Sequence\%20Models.pdf\#page=15}{PDF
pp.~15--16}.}

Why can repeated products of derivatives cause vanishing or exploding
gradients? Does clipping large gradients fix the vanishing case?

\begin{solution}[20mm]{5.1-S03}
% Write your answer here, usually 1--3 sentences.

\end{solution}
\end{minipage}\par\medskip

% QUESTION 5.1-S04
\noindent\begin{minipage}{\linewidth}

\subsubsection{5.1-S04 · Short}\label{s04-short-12}

\emph{Source:
\href{../Lecture\%20Slides/5.1.\%20Sequence\%20Models.pdf\#page=20}{PDF
pp.~20--21}.}

State the roles of the LSTM forget, input, and output gates in short
phrases.

\begin{solution}[20mm]{5.1-S04}
% Write your answer here, usually 1--3 sentences.

\end{solution}
\end{minipage}\par\medskip

% QUESTION 5.1-S05
\noindent\begin{minipage}{\linewidth}

\subsubsection{5.1-S05 · Short}\label{s05-short-12}

\emph{Source:
\href{../Lecture\%20Slides/5.1.\%20Sequence\%20Models.pdf\#page=25}{PDF
pp.~25--27}.}

What makes a temporal convolution causal? How can dilation expand access
to earlier time steps?

\begin{solution}[20mm]{5.1-S05}
% Write your answer here, usually 1--3 sentences.

\end{solution}
\end{minipage}\par\medskip

% QUESTION 5.1-S06
\noindent\begin{minipage}{\linewidth}

\subsubsection{5.1-S06 · Short}\label{s06-short-12}

\emph{Source:
\href{../Lecture\%20Slides/5.1.\%20Sequence\%20Models.pdf\#page=29}{PDF
pp.~29}.}

What is exposure bias in a teacher-forced decoder? Why can an early
inference-time mistake affect later outputs?

\begin{solution}[20mm]{5.1-S06}
% Write your answer here, usually 1--3 sentences.

\end{solution}
\end{minipage}\par\medskip

% QUESTION 5.1-S07
\noindent\begin{minipage}{\linewidth}

\subsubsection{5.1-S07 · Short}\label{s07-short-12}

\emph{Source:
\href{../Lecture\%20Slides/5.1.\%20Sequence\%20Models.pdf\#page=28}{PDF
pp.~28--30}.}

Why can a single fixed-size context vector become a bottleneck for long
input sequences?

\begin{solution}[20mm]{5.1-S07}
% Write your answer here, usually 1--3 sentences.

\end{solution}
\end{minipage}\par\medskip

\subsection{Easy calculations}\label{easy-calculations-12}

% QUESTION 5.1-C01
\noindent\begin{minipage}{\linewidth}

\subsubsection{5.1-C01 · Calculation}\label{c01-calculation-12}

\emph{Source:
\href{../Lecture\%20Slides/5.1.\%20Sequence\%20Models.pdf\#page=6}{PDF
pp.~6--9}.}

A vanilla RNN has input dimension 2 and hidden dimension 3. Count
parameters in \(h_t=\tanh(W_xx_t+W_hh_{t-1}+b)\), using one hidden bias
vector. Exclude any output layer.

\begin{solution}[26mm]{5.1-C01}
% Write your short calculation and final result here. Use $...$ or \[...\] for maths.

\end{solution}
\end{minipage}\par\medskip

% QUESTION 5.1-C02
\noindent\begin{minipage}{\linewidth}

\subsubsection{5.1-C02 · Calculation}\label{c02-calculation-12}

\emph{Source:
\href{../Lecture\%20Slides/5.1.\%20Sequence\%20Models.pdf\#page=16}{PDF
pp.~16}.}

A gradient is \((3,4)\) with norm 5. Norm clipping limits it to norm 2.
Compute the common rescaling factor and the clipped vector.

\begin{solution}[26mm]{5.1-C02}
% Write your short calculation and final result here. Use $...$ or \[...\] for maths.

\end{solution}
\end{minipage}\par\medskip

% QUESTION 5.1-C03
\noindent\begin{minipage}{\linewidth}

\subsubsection{5.1-C03 · Calculation}\label{c03-calculation-12}

\emph{Source:
\href{../Lecture\%20Slides/5.1.\%20Sequence\%20Models.pdf\#page=25}{PDF
pp.~25--27}.}

A causal 1D convolution has kernel size 3, dilation 2, and stride 1.
Using span \(=1+(K-1)d\), how many time positions wide is its receptive
field? How many positions are actually sampled by its three
coefficients?

\begin{solution}[26mm]{5.1-C03}
% Write your short calculation and final result here. Use $...$ or \[...\] for maths.

\end{solution}
\end{minipage}\par\medskip

\clearpage

\section{5.1. Extra Material}\label{extra-material}

\textbf{Slides:}
\href{../Lecture\%20Slides/5.1.\%20Extra\%20Material.pdf}{5.1. Extra
Material.pdf}. \textbf{20 questions: 10 MC · 7 short answers · 3 easy
calculations.}

\subsection{Multiple choice}\label{multiple-choice-13}

% QUESTION 5.1E-M01
\noindent\begin{minipage}{\linewidth}

\subsubsection{5.1E-M01 · MC}\label{e-m01-mc}

\emph{Source:
\href{../Lecture\%20Slides/5.1.\%20Extra\%20Material.pdf\#page=5}{PDF
pp.~5--6}.}

In the tiny LSTM example, vocabulary size \(V\) determines:

\begin{itemize}
\tightlist
\item
  A. The number of forget gates per hidden unit.
\item
  B. The hidden-state size regardless of the architecture choice.
\item
  C. The number of time steps in every sequence.
\item
  D. The length of a one-hot input vector.
\end{itemize}

\begin{solution}[8mm]{5.1E-M01}
% Write A, B, C, or D here. You may add a brief reason.

\end{solution}
\end{minipage}\par\medskip

% QUESTION 5.1E-M02
\noindent\begin{minipage}{\linewidth}

\subsubsection{5.1E-M02 · MC}\label{e-m02-mc}

\emph{Source:
\href{../Lecture\%20Slides/5.1.\%20Extra\%20Material.pdf\#page=5}{PDF
pp.~5--6}.}

With hidden size \(H\), the LSTM cell state and hidden state have:

\begin{itemize}
\tightlist
\item
  A. \(4H\) components each.
\item
  B. \(H\) components each.
\item
  C. One component per vocabulary item.
\item
  D. One component per training sequence.
\end{itemize}

\begin{solution}[8mm]{5.1E-M02}
% Write A, B, C, or D here. You may add a brief reason.

\end{solution}
\end{minipage}\par\medskip

% QUESTION 5.1E-M03
\noindent\begin{minipage}{\linewidth}

\subsubsection{5.1E-M03 · MC}\label{e-m03-mc}

\emph{Source:
\href{../Lecture\%20Slides/5.1.\%20Extra\%20Material.pdf\#page=5}{PDF
pp.~5, 7}.}

For column-vector inputs, a single gate's input matrix \(W_{xf}\) has
shape:

\begin{itemize}
\tightlist
\item
  A. \(V\times V\).
\item
  B. \(H\times V\).
\item
  C. \(H\times H\).
\item
  D. \(V\times H\) when used directly as \(W_{xf}x\).
\end{itemize}

\begin{solution}[8mm]{5.1E-M03}
% Write A, B, C, or D here. You may add a brief reason.

\end{solution}
\end{minipage}\par\medskip

% QUESTION 5.1E-M04
\noindent\begin{minipage}{\linewidth}

\subsubsection{5.1E-M04 · MC}\label{e-m04-mc}

\emph{Source:
\href{../Lecture\%20Slides/5.1.\%20Extra\%20Material.pdf\#page=5}{PDF
pp.~5, 7}.}

A single gate's recurrent matrix \(W_{hf}\) has shape:

\begin{itemize}
\tightlist
\item
  A. \(4V\times H\).
\item
  B. \(V\times H\).
\item
  C. \(V\times V\).
\item
  D. \(H\times H\).
\end{itemize}

\begin{solution}[8mm]{5.1E-M04}
% Write A, B, C, or D here. You may add a brief reason.

\end{solution}
\end{minipage}\par\medskip

% QUESTION 5.1E-M05
\noindent\begin{minipage}{\linewidth}

\subsubsection{5.1E-M05 · MC}\label{e-m05-mc}

\emph{Source:
\href{../Lecture\%20Slides/5.1.\%20Extra\%20Material.pdf\#page=6}{PDF
pp.~6--7}.}

Multiplying an input matrix by the one-hot vector for letter b:

\begin{itemize}
\tightlist
\item
  A. Selects the column corresponding to b.
\item
  B. Averages all vocabulary columns.
\item
  C. Selects all recurrent weights at once.
\item
  D. Selects the row for the predicted next letter.
\end{itemize}

\begin{solution}[8mm]{5.1E-M05}
% Write A, B, C, or D here. You may add a brief reason.

\end{solution}
\end{minipage}\par\medskip

% QUESTION 5.1E-M06
\noindent\begin{minipage}{\linewidth}

\subsubsection{5.1E-M06 · MC}\label{e-m06-mc}

\emph{Source:
\href{../Lecture\%20Slides/5.1.\%20Extra\%20Material.pdf\#page=2}{PDF
pp.~2--4, 7}.}

If a forget-gate component is close to 0, it:

\begin{itemize}
\tightlist
\item
  A. Retains nearly all the previous cell component.
\item
  B. Forces the output gate to 1.
\item
  C. Suppresses that previous cell-state contribution.
\item
  D. Removes the corresponding learned weights permanently.
\end{itemize}

\begin{solution}[8mm]{5.1E-M06}
% Write A, B, C, or D here. You may add a brief reason.

\end{solution}
\end{minipage}\par\medskip

% QUESTION 5.1E-M07
\noindent\begin{minipage}{\linewidth}

\subsubsection{5.1E-M07 · MC}\label{e-m07-mc}

\emph{Source:
\href{../Lecture\%20Slides/5.1.\%20Extra\%20Material.pdf\#page=2}{PDF
pp.~2--4, 7}.}

If an input-gate component is zero, the candidate's contribution to that
cell component is:

\begin{itemize}
\tightlist
\item
  A. Added without scaling.
\item
  B. Doubled.
\item
  C. Used as the output vocabulary size.
\item
  D. Blocked.
\end{itemize}

\begin{solution}[8mm]{5.1E-M07}
% Write A, B, C, or D here. You may add a brief reason.

\end{solution}
\end{minipage}\par\medskip

% QUESTION 5.1E-M08
\noindent\begin{minipage}{\linewidth}

\subsubsection{5.1E-M08 · MC}\label{e-m08-mc}

\emph{Source:
\href{../Lecture\%20Slides/5.1.\%20Extra\%20Material.pdf\#page=2}{PDF
pp.~2--4, 7}.}

The output gate primarily controls:

\begin{itemize}
\tightlist
\item
  A. The size of the vocabulary.
\item
  B. The number of stored training sequences.
\item
  C. How much of the transformed cell state is exposed as hidden output.
\item
  D. Which input column is selected.
\end{itemize}

\begin{solution}[8mm]{5.1E-M08}
% Write A, B, C, or D here. You may add a brief reason.

\end{solution}
\end{minipage}\par\medskip

% QUESTION 5.1E-M09
\noindent\begin{minipage}{\linewidth}

\subsubsection{5.1E-M09 · MC}\label{e-m09-mc}

\emph{Source:
\href{../Lecture\%20Slides/5.1.\%20Extra\%20Material.pdf\#page=6}{PDF
pp.~6}.}

Which quantities normally change while a trained LSTM reads a sequence
without training updates?

\begin{itemize}
\tightlist
\item
  A. Vocabulary size at every step.
\item
  B. All layer dimensions at every step.
\item
  C. Only the weights.
\item
  D. \(c_t\) and \(h_t\), while learned weights stay fixed.
\end{itemize}

\begin{solution}[8mm]{5.1E-M09}
% Write A, B, C, or D here. You may add a brief reason.

\end{solution}
\end{minipage}\par\medskip

% QUESTION 5.1E-M10
\noindent\begin{minipage}{\linewidth}

\subsubsection{5.1E-M10 · MC}\label{e-m10-mc}

\emph{Source:
\href{../Lecture\%20Slides/5.1.\%20Extra\%20Material.pdf\#page=5}{PDF
pp.~5--7}.}

The next-token output matrix maps hidden size \(H\) to vocabulary size
\(V\). Its output before softmax consists of:

\begin{itemize}
\tightlist
\item
  A. One probability automatically guaranteed to sum to 1.
\item
  B. \(V\) logits.
\item
  C. \(4H\) forget-gate values.
\item
  D. \(H\) independent cell memories.
\end{itemize}

\begin{solution}[8mm]{5.1E-M10}
% Write A, B, C, or D here. You may add a brief reason.

\end{solution}
\end{minipage}\par\medskip

\subsection{Short answers}\label{short-answers-13}

% QUESTION 5.1E-S01
\noindent\begin{minipage}{\linewidth}

\subsubsection{5.1E-S01 · Short}\label{e-s01-short}

\emph{Source:
\href{../Lecture\%20Slides/5.1.\%20Extra\%20Material.pdf\#page=2}{PDF
pp.~2--6}.}

Distinguish the roles of cell state \(c_t\) and hidden state \(h_t\).

\begin{solution}[20mm]{5.1E-S01}
% Write your answer here, usually 1--3 sentences.

\end{solution}
\end{minipage}\par\medskip

% QUESTION 5.1E-S02
\noindent\begin{minipage}{\linewidth}

\subsubsection{5.1E-S02 · Short}\label{e-s02-short}

\emph{Source:
\href{../Lecture\%20Slides/5.1.\%20Extra\%20Material.pdf\#page=6}{PDF
pp.~6}.}

Why is it incorrect to say that the cell state necessarily contains one
entry per letter in the vocabulary?

\begin{solution}[20mm]{5.1E-S02}
% Write your answer here, usually 1--3 sentences.

\end{solution}
\end{minipage}\par\medskip

% QUESTION 5.1E-S03
\noindent\begin{minipage}{\linewidth}

\subsubsection{5.1E-S03 · Short}\label{e-s03-short}

\emph{Source:
\href{../Lecture\%20Slides/5.1.\%20Extra\%20Material.pdf\#page=5}{PDF
pp.~5--6}.}

If a new vocabulary symbol is added while \(H\) stays fixed, which
input/output matrix dimensions change? Do \(c_t\) and \(h_t\) become
longer?

\begin{solution}[20mm]{5.1E-S03}
% Write your answer here, usually 1--3 sentences.

\end{solution}
\end{minipage}\par\medskip

% QUESTION 5.1E-S04
\noindent\begin{minipage}{\linewidth}

\subsubsection{5.1E-S04 · Short}\label{e-s04-short}

\emph{Source:
\href{../Lecture\%20Slides/5.1.\%20Extra\%20Material.pdf\#page=5}{PDF
pp.~5--6}.}

If \(H\) increases while the vocabulary stays fixed, which families of
parameters grow: input-gate matrices, recurrent matrices, output
projection, or all three?

\begin{solution}[20mm]{5.1E-S04}
% Write your answer here, usually 1--3 sentences.

\end{solution}
\end{minipage}\par\medskip

% QUESTION 5.1E-S05
\noindent\begin{minipage}{\linewidth}

\subsubsection{5.1E-S05 · Short}\label{e-s05-short}

\emph{Source:
\href{../Lecture\%20Slides/5.1.\%20Extra\%20Material.pdf\#page=6}{PDF
pp.~6--7}.}

Where is the learned input representation of a letter stored in a
one-hot-input LSTM? Why can the same letter produce different hidden
outputs in different contexts?

\begin{solution}[20mm]{5.1E-S05}
% Write your answer here, usually 1--3 sentences.

\end{solution}
\end{minipage}\par\medskip

% QUESTION 5.1E-S06
\noindent\begin{minipage}{\linewidth}

\subsubsection{5.1E-S06 · Short}\label{e-s06-short}

\emph{Source:
\href{../Lecture\%20Slides/5.1.\%20Extra\%20Material.pdf\#page=2}{PDF
pp.~2--4, 7}.}

For \(c_t=f_t\odot c_{t-1}+i_t\odot g_t\), what gate settings would
approximately preserve the previous cell state without writing new
content?

\begin{solution}[20mm]{5.1E-S06}
% Write your answer here, usually 1--3 sentences.

\end{solution}
\end{minipage}\par\medskip

% QUESTION 5.1E-S07
\noindent\begin{minipage}{\linewidth}

\subsubsection{5.1E-S07 · Short}\label{e-s07-short}

\emph{Source:
\href{../Lecture\%20Slides/5.1.\%20Extra\%20Material.pdf\#page=2}{PDF
pp.~2--4, 7}.}

Why use sigmoid outputs for the forget/input/output gates, but allow the
candidate update to contain negative as well as positive values?

\begin{solution}[20mm]{5.1E-S07}
% Write your answer here, usually 1--3 sentences.

\end{solution}
\end{minipage}\par\medskip

\subsection{Easy calculations}\label{easy-calculations-13}

% QUESTION 5.1E-C01
\noindent\begin{minipage}{\linewidth}

\subsubsection{5.1E-C01 · Calculation}\label{e-c01-calculation}

\emph{Source:
\href{../Lecture\%20Slides/5.1.\%20Extra\%20Material.pdf\#page=5}{PDF
pp.~5}.}

Count the parameters of an LSTM with input vocabulary size \(V=3\) and
hidden size \(H=2\). Use four input matrices, four recurrent matrices,
and one bias vector per gate: \(4(HV+H^2+H)\). Exclude the output
projection.

\begin{solution}[26mm]{5.1E-C01}
% Write your short calculation and final result here. Use $...$ or \[...\] for maths.

\end{solution}
\end{minipage}\par\medskip

% QUESTION 5.1E-C02
\noindent\begin{minipage}{\linewidth}

\subsubsection{5.1E-C02 · Calculation}\label{e-c02-calculation}

\emph{Source:
\href{../Lecture\%20Slides/5.1.\%20Extra\%20Material.pdf\#page=7}{PDF
pp.~7}.}

For one cell component, \(c_{t-1}=2\), \(f_t=0.5\), \(i_t=0.25\), and
\(g_t=0.8\). Compute \(c_t=f_tc_{t-1}+i_tg_t\).

\begin{solution}[26mm]{5.1E-C02}
% Write your short calculation and final result here. Use $...$ or \[...\] for maths.

\end{solution}
\end{minipage}\par\medskip

% QUESTION 5.1E-C03
\noindent\begin{minipage}{\linewidth}

\subsubsection{5.1E-C03 · Calculation}\label{e-c03-calculation}

\emph{Source:
\href{../Lecture\%20Slides/5.1.\%20Extra\%20Material.pdf\#page=7}{PDF
pp.~7}.}

For one hidden component, the output gate is \(o_t=0.5\) and
\(\tanh(c_t)=0.8\). Compute \(h_t=o_t\tanh(c_t)\). What would \(h_t\) be
if \(o_t=0\)?

\begin{solution}[26mm]{5.1E-C03}
% Write your short calculation and final result here. Use $...$ or \[...\] for maths.

\end{solution}
\end{minipage}\par\medskip

\clearpage

\section{5.2. Transformers}\label{transformers}

\textbf{Slides:}
\href{../Lecture\%20Slides/5.2.\%20Transformers.pdf}{5.2.
Transformers.pdf}. \textbf{20 questions: 10 MC · 7 short answers · 3
easy calculations.}

\subsection{Multiple choice}\label{multiple-choice-14}

% QUESTION 5.2-M01
\noindent\begin{minipage}{\linewidth}

\subsubsection{5.2-M01 · MC}\label{m01-mc-13}

\emph{Source:
\href{../Lecture\%20Slides/5.2.\%20Transformers.pdf\#page=8}{PDF
pp.~8--13}.}

Before softmax, attention scores are computed from:

\begin{itemize}
\tightlist
\item
  A. Ground-truth class labels.
\item
  B. Only the value-vector norms.
\item
  C. Only token indices, without content.
\item
  D. Query--key similarities.
\end{itemize}

\begin{solution}[8mm]{5.2-M01}
% Write A, B, C, or D here. You may add a brief reason.

\end{solution}
\end{minipage}\par\medskip

% QUESTION 5.2-M02
\noindent\begin{minipage}{\linewidth}

\subsubsection{5.2-M02 · MC}\label{m02-mc-13}

\emph{Source:
\href{../Lecture\%20Slides/5.2.\%20Transformers.pdf\#page=9}{PDF
pp.~9--13}.}

The output for one query is:

\begin{itemize}
\tightlist
\item
  A. An unweighted sum of keys.
\item
  B. A fixed learned vector independent of the inputs.
\item
  C. The query with its largest component removed.
\item
  D. A weighted sum of value vectors.
\end{itemize}

\begin{solution}[8mm]{5.2-M02}
% Write A, B, C, or D here. You may add a brief reason.

\end{solution}
\end{minipage}\par\medskip

% QUESTION 5.2-M03
\noindent\begin{minipage}{\linewidth}

\subsubsection{5.2-M03 · MC}\label{m03-mc-13}

\emph{Source:
\href{../Lecture\%20Slides/5.2.\%20Transformers.pdf\#page=13}{PDF
pp.~13}.}

Dividing dot products by \(\sqrt{d_k}\) helps:

\begin{itemize}
\tightlist
\item
  A. Replace positional information.
\item
  B. Reduce the number of tokens.
\item
  C. Control score scale and avoid excessive softmax saturation.
\item
  D. Make all attention weights equal.
\end{itemize}

\begin{solution}[8mm]{5.2-M03}
% Write A, B, C, or D here. You may add a brief reason.

\end{solution}
\end{minipage}\par\medskip

% QUESTION 5.2-M04
\noindent\begin{minipage}{\linewidth}

\subsubsection{5.2-M04 · MC}\label{m04-mc-13}

\emph{Source:
\href{../Lecture\%20Slides/5.2.\%20Transformers.pdf\#page=11}{PDF
pp.~11, 15--17}.}

Why supply positional information to self-attention?

\begin{itemize}
\tightlist
\item
  A. To remove all learned embeddings.
\item
  B. To convert every token into a class label.
\item
  C. Unmasked self-attention alone does not identify sequence positions.
\item
  D. To prevent GPU parallelism.
\end{itemize}

\begin{solution}[8mm]{5.2-M04}
% Write A, B, C, or D here. You may add a brief reason.

\end{solution}
\end{minipage}\par\medskip

% QUESTION 5.2-M05
\noindent\begin{minipage}{\linewidth}

\subsubsection{5.2-M05 · MC}\label{m05-mc-13}

\emph{Source:
\href{../Lecture\%20Slides/5.2.\%20Transformers.pdf\#page=19}{PDF
pp.~19--20}.}

Subword tokenization is useful partly because it:

\begin{itemize}
\tightlist
\item
  A. Assigns a separate vocabulary to every sentence.
\item
  B. Represents rare words using smaller reusable pieces.
\item
  C. Requires each word to be a unique vocabulary entry.
\item
  D. Eliminates all sequence-length limits.
\end{itemize}

\begin{solution}[8mm]{5.2-M05}
% Write A, B, C, or D here. You may add a brief reason.

\end{solution}
\end{minipage}\par\medskip

% QUESTION 5.2-M06
\noindent\begin{minipage}{\linewidth}

\subsubsection{5.2-M06 · MC}\label{m06-mc-13}

\emph{Source:
\href{../Lecture\%20Slides/5.2.\%20Transformers.pdf\#page=20}{PDF
pp.~20--22}.}

A contextual embedding differs from a basic token lookup because it:

\begin{itemize}
\tightlist
\item
  A. Is necessarily identical for every occurrence of a token.
\item
  B. Can reflect the surrounding tokens available to the model.
\item
  C. Has no learned parameters anywhere in its computation.
\item
  D. Must always be one-hot.
\end{itemize}

\begin{solution}[8mm]{5.2-M06}
% Write A, B, C, or D here. You may add a brief reason.

\end{solution}
\end{minipage}\par\medskip

% QUESTION 5.2-M07
\noindent\begin{minipage}{\linewidth}

\subsubsection{5.2-M07 · MC}\label{m07-mc-13}

\emph{Source:
\href{../Lecture\%20Slides/5.2.\%20Transformers.pdf\#page=24}{PDF
pp.~24}.}

Multi-head attention uses:

\begin{itemize}
\tightlist
\item
  A. A separate target label for every head.
\item
  B. Only one shared scalar weight for all token pairs.
\item
  C. A different vocabulary for each head.
\item
  D. Several learned projection sets to model relationships in parallel.
\end{itemize}

\begin{solution}[8mm]{5.2-M07}
% Write A, B, C, or D here. You may add a brief reason.

\end{solution}
\end{minipage}\par\medskip

% QUESTION 5.2-M08
\noindent\begin{minipage}{\linewidth}

\subsubsection{5.2-M08 · MC}\label{m08-mc-13}

\emph{Source:
\href{../Lecture\%20Slides/5.2.\%20Transformers.pdf\#page=30}{PDF
pp.~30--31}.}

At input position \(t\) predicting token \(t+1\), the usual causal mask
permits attention to:

\begin{itemize}
\tightlist
\item
  A. Only position \(t+1\).
\item
  B. Only future positions.
\item
  C. Input positions \(1,\ldots,t\).
\item
  D. Every position except \(t\).
\end{itemize}

\begin{solution}[8mm]{5.2-M08}
% Write A, B, C, or D here. You may add a brief reason.

\end{solution}
\end{minipage}\par\medskip

% QUESTION 5.2-M09
\noindent\begin{minipage}{\linewidth}

\subsubsection{5.2-M09 · MC}\label{m09-mc-13}

\emph{Source:
\href{../Lecture\%20Slides/5.2.\%20Transformers.pdf\#page=32}{PDF
pp.~32--35}.}

In encoder--decoder cross-attention:

\begin{itemize}
\tightlist
\item
  A. Queries come from decoder representations; keys and values come
  from encoder output.
\item
  B. Queries, keys, and values must all come from the same sequence.
\item
  C. The encoder queries future decoder outputs.
\item
  D. Keys are the future target labels.
\end{itemize}

\begin{solution}[8mm]{5.2-M09}
% Write A, B, C, or D here. You may add a brief reason.

\end{solution}
\end{minipage}\par\medskip

% QUESTION 5.2-M10
\noindent\begin{minipage}{\linewidth}

\subsubsection{5.2-M10 · MC}\label{m10-mc-13}

\emph{Source:
\href{../Lecture\%20Slides/5.2.\%20Transformers.pdf\#page=36}{PDF
pp.~36}.}

For fixed dimensions, doubling token count multiplies the number of
entries in a full attention score matrix by:

\begin{itemize}
\tightlist
\item
  A. Two.
\item
  B. One: it stays unchanged.
\item
  C. Eight.
\item
  D. Four.
\end{itemize}

\begin{solution}[8mm]{5.2-M10}
% Write A, B, C, or D here. You may add a brief reason.

\end{solution}
\end{minipage}\par\medskip

\subsection{Short answers}\label{short-answers-14}

% QUESTION 5.2-S01
\noindent\begin{minipage}{\linewidth}

\subsubsection{5.2-S01 · Short}\label{s01-short-13}

\emph{Source:
\href{../Lecture\%20Slides/5.2.\%20Transformers.pdf\#page=23}{PDF
pp.~23--27}.}

What different jobs are performed by self-attention and the per-token
feed-forward MLP in a transformer block?

\begin{solution}[20mm]{5.2-S01}
% Write your answer here, usually 1--3 sentences.

\end{solution}
\end{minipage}\par\medskip

% QUESTION 5.2-S02
\noindent\begin{minipage}{\linewidth}

\subsubsection{5.2-S02 · Short}\label{s02-short-13}

\emph{Source:
\href{../Lecture\%20Slides/5.2.\%20Transformers.pdf\#page=25}{PDF
pp.~25, 27}.}

What do the residual connection and layer normalization contribute in an
Add \& Norm component?

\begin{solution}[20mm]{5.2-S02}
% Write your answer here, usually 1--3 sentences.

\end{solution}
\end{minipage}\par\medskip

% QUESTION 5.2-S03
\noindent\begin{minipage}{\linewidth}

\subsubsection{5.2-S03 · Short}\label{s03-short-13}

\emph{Source:
\href{../Lecture\%20Slides/5.2.\%20Transformers.pdf\#page=28}{PDF
pp.~28--35}.}

Contrast encoder-only, decoder-only, and encoder--decoder models by
their attention access and one typical task each.

\begin{solution}[20mm]{5.2-S03}
% Write your answer here, usually 1--3 sentences.

\end{solution}
\end{minipage}\par\medskip

% QUESTION 5.2-S04
\noindent\begin{minipage}{\linewidth}

\subsubsection{5.2-S04 · Short}\label{s04-short-13}

\emph{Source:
\href{../Lecture\%20Slides/5.2.\%20Transformers.pdf\#page=30}{PDF
pp.~30--31}.}

Why can next-token training evaluate all known sequence positions in
parallel while ordinary autoregressive generation produces new tokens
sequentially?

\begin{solution}[20mm]{5.2-S04}
% Write your answer here, usually 1--3 sentences.

\end{solution}
\end{minipage}\par\medskip

% QUESTION 5.2-S05
\noindent\begin{minipage}{\linewidth}

\subsubsection{5.2-S05 · Short}\label{s05-short-13}

\emph{Source:
\href{../Lecture\%20Slides/5.2.\%20Transformers.pdf\#page=16}{PDF
pp.~16--17}.}

Contrast learned absolute position embeddings with RoPE in terms of
where positional information is introduced.

\begin{solution}[20mm]{5.2-S05}
% Write your answer here, usually 1--3 sentences.

\end{solution}
\end{minipage}\par\medskip

% QUESTION 5.2-S06
\noindent\begin{minipage}{\linewidth}

\subsubsection{5.2-S06 · Short}\label{s06-short-13}

\emph{Source:
\href{../Lecture\%20Slides/5.2.\%20Transformers.pdf\#page=36}{PDF
pp.~36}.}

What does a KV cache store during autoregressive generation, and what
repeated computation does it avoid?

\begin{solution}[20mm]{5.2-S06}
% Write your answer here, usually 1--3 sentences.

\end{solution}
\end{minipage}\par\medskip

% QUESTION 5.2-S07
\noindent\begin{minipage}{\linewidth}

\subsubsection{5.2-S07 · Short}\label{s07-short-13}

\emph{Source:
\href{../Lecture\%20Slides/5.2.\%20Transformers.pdf\#page=37}{PDF
pp.~37--40}.}

How does a Vision Transformer turn an image into a sequence? Why must
spatial positions also be represented?

\begin{solution}[20mm]{5.2-S07}
% Write your answer here, usually 1--3 sentences.

\end{solution}
\end{minipage}\par\medskip

\subsection{Easy calculations}\label{easy-calculations-14}

% QUESTION 5.2-C01
\noindent\begin{minipage}{\linewidth}

\subsubsection{5.2-C01 · Calculation}\label{c01-calculation-13}

\emph{Source:
\href{../Lecture\%20Slides/5.2.\%20Transformers.pdf\#page=9}{PDF
pp.~9--13}.}

Attention weights are \((0.25,0.75)\) over scalar values 2 and 6.
Compute the attention output.

\begin{solution}[26mm]{5.2-C01}
% Write your short calculation and final result here. Use $...$ or \[...\] for maths.

\end{solution}
\end{minipage}\par\medskip

% QUESTION 5.2-C02
\noindent\begin{minipage}{\linewidth}

\subsubsection{5.2-C02 · Calculation}\label{c02-calculation-13}

\emph{Source:
\href{../Lecture\%20Slides/5.2.\%20Transformers.pdf\#page=13}{PDF
pp.~13, 24}.}

One attention head maps model dimension 4 into query, key, and value
dimensions of 2 each. Count the parameters in its three projection
matrices. Assume no biases and exclude the output projection.

\begin{solution}[26mm]{5.2-C02}
% Write your short calculation and final result here. Use $...$ or \[...\] for maths.

\end{solution}
\end{minipage}\par\medskip

% QUESTION 5.2-C03
\noindent\begin{minipage}{\linewidth}

\subsubsection{5.2-C03 · Calculation}\label{c03-calculation-13}

\emph{Source:
\href{../Lecture\%20Slides/5.2.\%20Transformers.pdf\#page=37}{PDF
pp.~37--39}.}

A \(32\times32\) image is divided into non-overlapping \(8\times8\)
patches. How many patch tokens are produced? How many tokens are there
after adding one CLS token?

\begin{solution}[26mm]{5.2-C03}
% Write your short calculation and final result here. Use $...$ or \[...\] for maths.

\end{solution}
\end{minipage}\par\medskip

\clearpage

\section{5.3. Foundation models}\label{foundation-models}

\textbf{Slides:}
\href{../Lecture\%20Slides/5.3.\%20Foundation\%20models.pdf}{5.3.
Foundation models.pdf}. \textbf{20 questions: 10 MC · 7 short answers ·
3 easy calculations.}

\subsection{Multiple choice}\label{multiple-choice-15}

% QUESTION 5.3-M01
\noindent\begin{minipage}{\linewidth}

\subsubsection{5.3-M01 · MC}\label{m01-mc-14}

\emph{Source:
\href{../Lecture\%20Slides/5.3.\%20Foundation\%20models.pdf\#page=9}{PDF
pp.~9--13}.}

A foundation model is primarily characterized by:

\begin{itemize}
\tightlist
\item
  A. A fixed minimum layer count.
\item
  B. Broad pretraining and adaptability to many downstream tasks.
\item
  C. Always using one particular architecture.
\item
  D. Being trained separately from scratch for each input.
\end{itemize}

\begin{solution}[8mm]{5.3-M01}
% Write A, B, C, or D here. You may add a brief reason.

\end{solution}
\end{minipage}\par\medskip

% QUESTION 5.3-M02
\noindent\begin{minipage}{\linewidth}

\subsubsection{5.3-M02 · MC}\label{m02-mc-14}

\emph{Source:
\href{../Lecture\%20Slides/5.3.\%20Foundation\%20models.pdf\#page=13}{PDF
pp.~13--15}.}

Self-supervised pretraining obtains targets:

\begin{itemize}
\tightlist
\item
  A. Only from interactive rewards.
\item
  B. From the training data itself.
\item
  C. From held-out test labels.
\item
  D. Only from manually annotated classes.
\end{itemize}

\begin{solution}[8mm]{5.3-M02}
% Write A, B, C, or D here. You may add a brief reason.

\end{solution}
\end{minipage}\par\medskip

% QUESTION 5.3-M03
\noindent\begin{minipage}{\linewidth}

\subsubsection{5.3-M03 · MC}\label{m03-mc-14}

\emph{Source:
\href{../Lecture\%20Slides/5.3.\%20Foundation\%20models.pdf\#page=14}{PDF
pp.~14}.}

Supervised fine-tuning primarily learns from:

\begin{itemize}
\tightlist
\item
  A. An automatic guarantee that responses are correct.
\item
  B. Demonstrations or supplied input--target examples.
\item
  C. Only the original pretraining vocabulary size.
\item
  D. Random actions without targets.
\end{itemize}

\begin{solution}[8mm]{5.3-M03}
% Write A, B, C, or D here. You may add a brief reason.

\end{solution}
\end{minipage}\par\medskip

% QUESTION 5.3-M04
\noindent\begin{minipage}{\linewidth}

\subsubsection{5.3-M04 · MC}\label{m04-mc-14}

\emph{Source:
\href{../Lecture\%20Slides/5.3.\%20Foundation\%20models.pdf\#page=14}{PDF
pp.~14}.}

Preference-based RL uses feedback that concerns:

\begin{itemize}
\tightlist
\item
  A. A required correct token label for every position.
\item
  B. Only the number of network layers.
\item
  C. Only how often a token occurs in raw text.
\item
  D. Which sampled outputs are preferred or receive greater reward.
\end{itemize}

\begin{solution}[8mm]{5.3-M04}
% Write A, B, C, or D here. You may add a brief reason.

\end{solution}
\end{minipage}\par\medskip

% QUESTION 5.3-M05
\noindent\begin{minipage}{\linewidth}

\subsubsection{5.3-M05 · MC}\label{m05-mc-14}

\emph{Source:
\href{../Lecture\%20Slides/5.3.\%20Foundation\%20models.pdf\#page=14}{PDF
pp.~14}.}

Which is an example of a verifiable reward?

\begin{itemize}
\tightlist
\item
  A. A random score independent of the output.
\item
  B. The number of training parameters alone.
\item
  C. Whether generated code passes specified unit tests.
\item
  D. A manually supplied next-token label by definition.
\end{itemize}

\begin{solution}[8mm]{5.3-M05}
% Write A, B, C, or D here. You may add a brief reason.

\end{solution}
\end{minipage}\par\medskip

% QUESTION 5.3-M06
\noindent\begin{minipage}{\linewidth}

\subsubsection{5.3-M06 · MC}\label{m06-mc-14}

\emph{Source:
\href{../Lecture\%20Slides/5.3.\%20Foundation\%20models.pdf\#page=19}{PDF
pp.~19--20}.}

CLIP learns:

\begin{itemize}
\tightlist
\item
  A. A separate encoder for every future class name.
\item
  B. A transition model for an RL environment.
\item
  C. A shared representation space for images and text.
\item
  D. Only a pixel reconstruction model.
\end{itemize}

\begin{solution}[8mm]{5.3-M06}
% Write A, B, C, or D here. You may add a brief reason.

\end{solution}
\end{minipage}\par\medskip

% QUESTION 5.3-M07
\noindent\begin{minipage}{\linewidth}

\subsubsection{5.3-M07 · MC}\label{m07-mc-14}

\emph{Source:
\href{../Lecture\%20Slides/5.3.\%20Foundation\%20models.pdf\#page=19}{PDF
pp.~19--20}.}

Zero-shot CLIP-style classification can compare an image with:

\begin{itemize}
\tightlist
\item
  A. Text embeddings of candidate class descriptions.
\item
  B. Only its original image filename.
\item
  C. Only classes represented by newly trained classifier weights.
\item
  D. A table of true test labels.
\end{itemize}

\begin{solution}[8mm]{5.3-M07}
% Write A, B, C, or D here. You may add a brief reason.

\end{solution}
\end{minipage}\par\medskip

% QUESTION 5.3-M08
\noindent\begin{minipage}{\linewidth}

\subsubsection{5.3-M08 · MC}\label{m08-mc-14}

\emph{Source:
\href{../Lecture\%20Slides/5.3.\%20Foundation\%20models.pdf\#page=21}{PDF
pp.~21}.}

DINOv2 is presented as learning visual representations through:

\begin{itemize}
\tightlist
\item
  A. Predicting text tokens only.
\item
  B. Fully labelled segmentation masks for every image.
\item
  C. Self-supervised teacher--student learning from image views.
\item
  D. An image generator competing with a discriminator.
\end{itemize}

\begin{solution}[8mm]{5.3-M08}
% Write A, B, C, or D here. You may add a brief reason.

\end{solution}
\end{minipage}\par\medskip

% QUESTION 5.3-M09
\noindent\begin{minipage}{\linewidth}

\subsubsection{5.3-M09 · MC}\label{m09-mc-14}

\emph{Source:
\href{../Lecture\%20Slides/5.3.\%20Foundation\%20models.pdf\#page=22}{PDF
pp.~22}.}

SAM is associated with:

\begin{itemize}
\tightlist
\item
  A. Segmentation.
\item
  B. Document-topic discovery.
\item
  C. Polynomial regression.
\item
  D. Tabular Q-learning.
\end{itemize}

\begin{solution}[8mm]{5.3-M09}
% Write A, B, C, or D here. You may add a brief reason.

\end{solution}
\end{minipage}\par\medskip

% QUESTION 5.3-M10
\noindent\begin{minipage}{\linewidth}

\subsubsection{5.3-M10 · MC}\label{m10-mc-14}

\emph{Source:
\href{../Lecture\%20Slides/5.3.\%20Foundation\%20models.pdf\#page=14}{PDF
pp.~14, 17, 19--25}.}

In this context, multimodal learning combines information from:

\begin{itemize}
\tightlist
\item
  A. Only different train/test splits.
\item
  B. Different modalities, such as images and text.
\item
  C. Only different learning rates.
\item
  D. Multiple random seeds for the same model.
\end{itemize}

\begin{solution}[8mm]{5.3-M10}
% Write A, B, C, or D here. You may add a brief reason.

\end{solution}
\end{minipage}\par\medskip

\subsection{Short answers}\label{short-answers-15}

% QUESTION 5.3-S01
\noindent\begin{minipage}{\linewidth}

\subsubsection{5.3-S01 · Short}\label{s01-short-14}

\emph{Source:
\href{../Lecture\%20Slides/5.3.\%20Foundation\%20models.pdf\#page=10}{PDF
pp.~10--14}.}

Distinguish broad pretraining from adaptation to a downstream task. Why
can reusing pretrained features reduce the need for task-specific
labelled data?

\begin{solution}[20mm]{5.3-S01}
% Write your answer here, usually 1--3 sentences.

\end{solution}
\end{minipage}\par\medskip

% QUESTION 5.3-S02
\noindent\begin{minipage}{\linewidth}

\subsubsection{5.3-S02 · Short}\label{s02-short-14}

\emph{Source:
\href{../Lecture\%20Slides/5.3.\%20Foundation\%20models.pdf\#page=12}{PDF
pp.~12, 26}.}

Why can biases in a widely reused foundation model affect many
downstream applications?

\begin{solution}[20mm]{5.3-S02}
% Write your answer here, usually 1--3 sentences.

\end{solution}
\end{minipage}\par\medskip

% QUESTION 5.3-S03
\noindent\begin{minipage}{\linewidth}

\subsubsection{5.3-S03 · Short}\label{s03-short-14}

\emph{Source:
\href{../Lecture\%20Slides/5.3.\%20Foundation\%20models.pdf\#page=12}{PDF
pp.~12}.}

For a mixture-of-experts model, why can total parameter count differ
from the number of parameters active for one token? Which count still
matters for storing all experts?

\begin{solution}[20mm]{5.3-S03}
% Write your answer here, usually 1--3 sentences.

\end{solution}
\end{minipage}\par\medskip

% QUESTION 5.3-S04
\noindent\begin{minipage}{\linewidth}

\subsubsection{5.3-S04 · Short}\label{s04-short-14}

\emph{Source:
\href{../Lecture\%20Slides/5.3.\%20Foundation\%20models.pdf\#page=14}{PDF
pp.~14}.}

Distinguish human demonstrations, preference rankings, and automatically
verifiable rewards as training signals.

\begin{solution}[20mm]{5.3-S04}
% Write your answer here, usually 1--3 sentences.

\end{solution}
\end{minipage}\par\medskip

% QUESTION 5.3-S05
\noindent\begin{minipage}{\linewidth}

\subsubsection{5.3-S05 · Short}\label{s05-short-14}

\emph{Source:
\href{../Lecture\%20Slides/5.3.\%20Foundation\%20models.pdf\#page=19}{PDF
pp.~19--20}.}

Explain CLIP-style zero-shot classification in two steps, from candidate
class descriptions to a predicted class.

\begin{solution}[20mm]{5.3-S05}
% Write your answer here, usually 1--3 sentences.

\end{solution}
\end{minipage}\par\medskip

% QUESTION 5.3-S06
\noindent\begin{minipage}{\linewidth}

\subsubsection{5.3-S06 · Short}\label{s06-short-14}

\emph{Source:
\href{../Lecture\%20Slides/5.3.\%20Foundation\%20models.pdf\#page=21}{PDF
pp.~21}.}

How can a visual encoder learn useful features without semantic class
labels? Describe the role of different views of one image in
teacher--student training.

\begin{solution}[20mm]{5.3-S06}
% Write your answer here, usually 1--3 sentences.

\end{solution}
\end{minipage}\par\medskip

% QUESTION 5.3-S07
\noindent\begin{minipage}{\linewidth}

\subsubsection{5.3-S07 · Short}\label{s07-short-14}

\emph{Source:
\href{../Lecture\%20Slides/5.3.\%20Foundation\%20models.pdf\#page=24}{PDF
pp.~24--25}.}

In a text-conditioned image generator, distinguish the role of the text
encoder from that of the image-generation model.

\begin{solution}[20mm]{5.3-S07}
% Write your answer here, usually 1--3 sentences.

\end{solution}
\end{minipage}\par\medskip

\subsection{Easy calculations}\label{easy-calculations-15}

% QUESTION 5.3-C01
\noindent\begin{minipage}{\linewidth}

\subsubsection{5.3-C01 · Calculation}\label{c01-calculation-14}

\emph{Source:
\href{../Lecture\%20Slides/5.3.\%20Foundation\%20models.pdf\#page=12}{PDF
pp.~12}.}

A toy model has 10 million parameters stored in FP32, using 4 bytes per
parameter. Ignoring all other tensors and metadata, compute its weight
storage in decimal MB (\(1\,\mathrm{MB}=10^6\) bytes).

\begin{solution}[26mm]{5.3-C01}
% Write your short calculation and final result here. Use $...$ or \[...\] for maths.

\end{solution}
\end{minipage}\par\medskip

% QUESTION 5.3-C02
\noindent\begin{minipage}{\linewidth}

\subsubsection{5.3-C02 · Calculation}\label{c02-calculation-14}

\emph{Source:
\href{../Lecture\%20Slides/5.3.\%20Foundation\%20models.pdf\#page=19}{PDF
pp.~19--20}.}

An image has unit embedding \((1,0)\). Two candidate descriptions have
unit embeddings \(t_A=(0.8,0.6)\) and \(t_B=(0,1)\). Compute both
dot-product similarities and select the higher-similarity description.

\begin{solution}[26mm]{5.3-C02}
% Write your short calculation and final result here. Use $...$ or \[...\] for maths.

\end{solution}
\end{minipage}\par\medskip

% QUESTION 5.3-C03
\noindent\begin{minipage}{\linewidth}

\subsubsection{5.3-C03 · Calculation}\label{c03-calculation-14}

\emph{Source:
\href{../Lecture\%20Slides/5.3.\%20Foundation\%20models.pdf\#page=12}{PDF
pp.~12}.}

A toy mixture-of-experts contains 8 experts with 4 million parameters
each and activates 2 experts per token. Ignore shared parameters.
Compute the total expert parameters and the active expert parameters per
token.

\begin{solution}[26mm]{5.3-C03}
% Write your short calculation and final result here. Use $...$ or \[...\] for maths.

\end{solution}
\end{minipage}\par\medskip

\clearpage

\section{5.4. Unsupervised learning}\label{unsupervised-learning}

\textbf{Slides:}
\href{../Lecture\%20Slides/5.4.\%20Unsupervised\%20learning.pdf}{5.4.
Unsupervised learning.pdf}. \textbf{20 questions: 10 MC · 7 short
answers · 3 easy calculations.}

\subsection{Multiple choice}\label{multiple-choice-16}

% QUESTION 5.4-M01
\noindent\begin{minipage}{\linewidth}

\subsubsection{5.4-M01 · MC}\label{m01-mc-15}

\emph{Source:
\href{../Lecture\%20Slides/5.4.\%20Unsupervised\%20learning.pdf\#page=6}{PDF
pp.~6--7}.}

Semi-supervised learning uses:

\begin{itemize}
\tightlist
\item
  A. Only a smaller supervised learning rate.
\item
  B. Only labels from the test set.
\item
  C. Labelled examples together with relevant unlabelled examples.
\item
  D. Only rewards from an interactive environment.
\end{itemize}

\begin{solution}[8mm]{5.4-M01}
% Write A, B, C, or D here. You may add a brief reason.

\end{solution}
\end{minipage}\par\medskip

% QUESTION 5.4-M02
\noindent\begin{minipage}{\linewidth}

\subsubsection{5.4-M02 · MC}\label{m02-mc-15}

\emph{Source:
\href{../Lecture\%20Slides/5.4.\%20Unsupervised\%20learning.pdf\#page=9}{PDF
pp.~9--11}.}

For centred data, the first principal component is the unit direction
that:

\begin{itemize}
\tightlist
\item
  A. Gives the greatest projected variance.
\item
  B. Maximizes labelled classification accuracy by definition.
\item
  C. Gives the smallest projected variance.
\item
  D. Selects the feature with the shortest name.
\end{itemize}

\begin{solution}[8mm]{5.4-M02}
% Write A, B, C, or D here. You may add a brief reason.

\end{solution}
\end{minipage}\par\medskip

% QUESTION 5.4-M03
\noindent\begin{minipage}{\linewidth}

\subsubsection{5.4-M03 · MC}\label{m03-mc-15}

\emph{Source:
\href{../Lecture\%20Slides/5.4.\%20Unsupervised\%20learning.pdf\#page=13}{PDF
pp.~13--18, 36}.}

A defensible interpretation of t-SNE is that it mainly aims to preserve:

\begin{itemize}
\tightlist
\item
  A. All pairwise distances exactly.
\item
  B. Local neighbourhood structure in a low-dimensional visualization.
\item
  C. Cluster sizes as exact population proportions.
\item
  D. Only the covariance eigenvalues.
\end{itemize}

\begin{solution}[8mm]{5.4-M03}
% Write A, B, C, or D here. You may add a brief reason.

\end{solution}
\end{minipage}\par\medskip

% QUESTION 5.4-M04
\noindent\begin{minipage}{\linewidth}

\subsubsection{5.4-M04 · MC}\label{m04-mc-15}

\emph{Source:
\href{../Lecture\%20Slides/5.4.\%20Unsupervised\%20learning.pdf\#page=15}{PDF
pp.~15}.}

The t-SNE perplexity hyperparameter is related to:

\begin{itemize}
\tightlist
\item
  A. The neural-network layer count.
\item
  B. The effective neighbourhood scale used by the method.
\item
  C. The number of labels per image.
\item
  D. The number of semantic classes, fixed exactly.
\end{itemize}

\begin{solution}[8mm]{5.4-M04}
% Write A, B, C, or D here. You may add a brief reason.

\end{solution}
\end{minipage}\par\medskip

% QUESTION 5.4-M05
\noindent\begin{minipage}{\linewidth}

\subsubsection{5.4-M05 · MC}\label{m05-mc-15}

\emph{Source:
\href{../Lecture\%20Slides/5.4.\%20Unsupervised\%20learning.pdf\#page=20}{PDF
pp.~20}.}

Standard k-means alternates:

\begin{itemize}
\tightlist
\item
  A. Maximizing variance and sorting eigenvectors.
\item
  B. Class-label annotation and supervised backpropagation.
\item
  C. Nearest-centre assignments and recomputation of centres as means.
\item
  D. Only merging existing clusters, never moving centres.
\end{itemize}

\begin{solution}[8mm]{5.4-M05}
% Write A, B, C, or D here. You may add a brief reason.

\end{solution}
\end{minipage}\par\medskip

% QUESTION 5.4-M06
\noindent\begin{minipage}{\linewidth}

\subsubsection{5.4-M06 · MC}\label{m06-mc-15}

\emph{Source:
\href{../Lecture\%20Slides/5.4.\%20Unsupervised\%20learning.pdf\#page=21}{PDF
pp.~21}.}

Agglomerative hierarchical clustering starts with:

\begin{itemize}
\tightlist
\item
  A. Individual points as clusters and progressively merges them.
\item
  B. Exactly one cluster per known class label.
\item
  C. A trained decoder generating new points.
\item
  D. All points permanently in one cluster.
\end{itemize}

\begin{solution}[8mm]{5.4-M06}
% Write A, B, C, or D here. You may add a brief reason.

\end{solution}
\end{minipage}\par\medskip

% QUESTION 5.4-M07
\noindent\begin{minipage}{\linewidth}

\subsubsection{5.4-M07 · MC}\label{m07-mc-15}

\emph{Source:
\href{../Lecture\%20Slides/5.4.\%20Unsupervised\%20learning.pdf\#page=22}{PDF
pp.~22}.}

Latent Dirichlet allocation models a document as:

\begin{itemize}
\tightlist
\item
  A. A single principal component.
\item
  B. A mixture of topics.
\item
  C. One mandatory hard cluster assignment only.
\item
  D. A sequence of RL states.
\end{itemize}

\begin{solution}[8mm]{5.4-M07}
% Write A, B, C, or D here. You may add a brief reason.

\end{solution}
\end{minipage}\par\medskip

% QUESTION 5.4-M08
\noindent\begin{minipage}{\linewidth}

\subsubsection{5.4-M08 · MC}\label{m08-mc-15}

\emph{Source:
\href{../Lecture\%20Slides/5.4.\%20Unsupervised\%20learning.pdf\#page=26}{PDF
pp.~26--27}.}

Which is a self-supervised pretext task?

\begin{itemize}
\tightlist
\item
  A. Predicting an image rotation that was applied automatically.
\item
  B. Asking a human to label every object category.
\item
  C. Learning solely from externally assigned rewards.
\item
  D. Selecting a model using test labels.
\end{itemize}

\begin{solution}[8mm]{5.4-M08}
% Write A, B, C, or D here. You may add a brief reason.

\end{solution}
\end{minipage}\par\medskip

% QUESTION 5.4-M09
\noindent\begin{minipage}{\linewidth}

\subsubsection{5.4-M09 · MC}\label{m09-mc-15}

\emph{Source:
\href{../Lecture\%20Slides/5.4.\%20Unsupervised\%20learning.pdf\#page=28}{PDF
pp.~28--29}.}

With positive and negative pairs, a contrastive objective generally
encourages:

\begin{itemize}
\tightlist
\item
  A. Similar representations for positives and separation from
  negatives.
\item
  B. Identical representations for every input.
\item
  C. Greater distance between all positive pairs.
\item
  D. Raw pixel reconstruction as its only objective.
\end{itemize}

\begin{solution}[8mm]{5.4-M09}
% Write A, B, C, or D here. You may add a brief reason.

\end{solution}
\end{minipage}\par\medskip

% QUESTION 5.4-M10
\noindent\begin{minipage}{\linewidth}

\subsubsection{5.4-M10 · MC}\label{m10-mc-15}

\emph{Source:
\href{../Lecture\%20Slides/5.4.\%20Unsupervised\%20learning.pdf\#page=30}{PDF
pp.~30--31}.}

The basic BYOL approach is notable for:

\begin{itemize}
\tightlist
\item
  A. Needing a separate classifier for every image.
\item
  B. Learning from two views using online/target networks without
  explicit negatives.
\item
  C. Requiring manually labelled negative examples.
\item
  D. Replacing the encoder with k-means.
\end{itemize}

\begin{solution}[8mm]{5.4-M10}
% Write A, B, C, or D here. You may add a brief reason.

\end{solution}
\end{minipage}\par\medskip

\subsection{Short answers}\label{short-answers-16}

% QUESTION 5.4-S01
\noindent\begin{minipage}{\linewidth}

\subsubsection{5.4-S01 · Short}\label{s01-short-15}

\emph{Source:
\href{../Lecture\%20Slides/5.4.\%20Unsupervised\%20learning.pdf\#page=8}{PDF
pp.~8--23}.}

Distinguish dimensionality reduction from clustering. What is the output
of each?

\begin{solution}[20mm]{5.4-S01}
% Write your answer here, usually 1--3 sentences.

\end{solution}
\end{minipage}\par\medskip

% QUESTION 5.4-S02
\noindent\begin{minipage}{\linewidth}

\subsubsection{5.4-S02 · Short}\label{s02-short-15}

\emph{Source:
\href{../Lecture\%20Slides/5.4.\%20Unsupervised\%20learning.pdf\#page=10}{PDF
pp.~10--11}.}

What does retaining 90\% of variance with PCA mean? Does it imply 90\%
classification accuracy?

\begin{solution}[20mm]{5.4-S02}
% Write your answer here, usually 1--3 sentences.

\end{solution}
\end{minipage}\par\medskip

% QUESTION 5.4-S03
\noindent\begin{minipage}{\linewidth}

\subsubsection{5.4-S03 · Short}\label{s03-short-15}

\emph{Source:
\href{../Lecture\%20Slides/5.4.\%20Unsupervised\%20learning.pdf\#page=13}{PDF
pp.~13--15}.}

Why should a visible gap between groups in a t-SNE plot not
automatically be interpreted as proof of separate true classes?

\begin{solution}[20mm]{5.4-S03}
% Write your answer here, usually 1--3 sentences.

\end{solution}
\end{minipage}\par\medskip

% QUESTION 5.4-S04
\noindent\begin{minipage}{\linewidth}

\subsubsection{5.4-S04 · Short}\label{s04-short-15}

\emph{Source:
\href{../Lecture\%20Slides/5.4.\%20Unsupervised\%20learning.pdf\#page=20}{PDF
pp.~20}.}

Why can k-means results depend on the initial centres? Does one run
guarantee the globally best clustering?

\begin{solution}[20mm]{5.4-S04}
% Write your answer here, usually 1--3 sentences.

\end{solution}
\end{minipage}\par\medskip

% QUESTION 5.4-S05
\noindent\begin{minipage}{\linewidth}

\subsubsection{5.4-S05 · Short}\label{s05-short-15}

\emph{Source:
\href{../Lecture\%20Slides/5.4.\%20Unsupervised\%20learning.pdf\#page=25}{PDF
pp.~25--28}.}

Why can self-supervised learning have a supervised-looking loss even
though semantic class labels are not manually supplied?

\begin{solution}[20mm]{5.4-S05}
% Write your answer here, usually 1--3 sentences.

\end{solution}
\end{minipage}\par\medskip

% QUESTION 5.4-S06
\noindent\begin{minipage}{\linewidth}

\subsubsection{5.4-S06 · Short}\label{s06-short-15}

\emph{Source:
\href{../Lecture\%20Slides/5.4.\%20Unsupervised\%20learning.pdf\#page=28}{PDF
pp.~28--31}.}

What is representation collapse? Why would mapping every input to the
same vector be unhelpful for downstream tasks?

\begin{solution}[20mm]{5.4-S06}
% Write your answer here, usually 1--3 sentences.

\end{solution}
\end{minipage}\par\medskip

% QUESTION 5.4-S07
\noindent\begin{minipage}{\linewidth}

\subsubsection{5.4-S07 · Short}\label{s07-short-15}

\emph{Source:
\href{../Lecture\%20Slides/5.4.\%20Unsupervised\%20learning.pdf\#page=6}{PDF
pp.~6--7, 33}.}

Why might unlabelled examples from an unrelated distribution fail to
help a semi-supervised task?

\begin{solution}[20mm]{5.4-S07}
% Write your answer here, usually 1--3 sentences.

\end{solution}
\end{minipage}\par\medskip

\subsection{Easy calculations}\label{easy-calculations-16}

% QUESTION 5.4-C01
\noindent\begin{minipage}{\linewidth}

\subsubsection{5.4-C01 · Calculation}\label{c01-calculation-15}

\emph{Source:
\href{../Lecture\%20Slides/5.4.\%20Unsupervised\%20learning.pdf\#page=10}{PDF
pp.~10--11}.}

Covariance eigenvalues are 6, 3, and 1. What fraction of total variance
is retained by the first two principal components?

\begin{solution}[26mm]{5.4-C01}
% Write your short calculation and final result here. Use $...$ or \[...\] for maths.

\end{solution}
\end{minipage}\par\medskip

% QUESTION 5.4-C02
\noindent\begin{minipage}{\linewidth}

\subsubsection{5.4-C02 · Calculation}\label{c02-calculation-15}

\emph{Source:
\href{../Lecture\%20Slides/5.4.\%20Unsupervised\%20learning.pdf\#page=20}{PDF
pp.~20}.}

Three 1D points, 2, 4, and 9, are assigned to the same k-means cluster.
Compute its updated centre.

\begin{solution}[26mm]{5.4-C02}
% Write your short calculation and final result here. Use $...$ or \[...\] for maths.

\end{solution}
\end{minipage}\par\medskip

% QUESTION 5.4-C03
\noindent\begin{minipage}{\linewidth}

\subsubsection{5.4-C03 · Calculation}\label{c03-calculation-15}

\emph{Source:
\href{../Lecture\%20Slides/5.4.\%20Unsupervised\%20learning.pdf\#page=20}{PDF
pp.~20}.}

Two 1D k-means centres are 0 and 10. For a point at 7, compute its
distance to each centre and state its next assignment.

\begin{solution}[26mm]{5.4-C03}
% Write your short calculation and final result here. Use $...$ or \[...\] for maths.

\end{solution}
\end{minipage}\par\medskip

\clearpage

\section{5.5. Generative Models}\label{generative-models}

\textbf{Slides:}
\href{../Lecture\%20Slides/5.5.\%20Generative\%20Models.pdf}{5.5.
Generative Models.pdf}. \textbf{20 questions: 10 MC · 7 short answers ·
3 easy calculations.}

\subsection{Multiple choice}\label{multiple-choice-17}

% QUESTION 5.5-M01
\noindent\begin{minipage}{\linewidth}

\subsubsection{5.5-M01 · MC}\label{m01-mc-16}

\emph{Source:
\href{../Lecture\%20Slides/5.5.\%20Generative\%20Models.pdf\#page=3}{PDF
pp.~3--6}.}

Generative modelling primarily aims to:

\begin{itemize}
\tightlist
\item
  A. Only select a class boundary.
\item
  B. Store every possible input without a model.
\item
  C. Remove all uncertainty from observed data.
\item
  D. Learn a data distribution or sampling mechanism from examples.
\end{itemize}

\begin{solution}[8mm]{5.5-M01}
% Write A, B, C, or D here. You may add a brief reason.

\end{solution}
\end{minipage}\par\medskip

% QUESTION 5.5-M02
\noindent\begin{minipage}{\linewidth}

\subsubsection{5.5-M02 · MC}\label{m02-mc-16}

\emph{Source:
\href{../Lecture\%20Slides/5.5.\%20Generative\%20Models.pdf\#page=7}{PDF
pp.~7}.}

A basic autoencoder is normally trained to:

\begin{itemize}
\tightlist
\item
  A. Sort examples by filename.
\item
  B. Predict an unrelated class label.
\item
  C. Reconstruct its input through an internal representation.
\item
  D. Estimate a reward for an action.
\end{itemize}

\begin{solution}[8mm]{5.5-M02}
% Write A, B, C, or D here. You may add a brief reason.

\end{solution}
\end{minipage}\par\medskip

% QUESTION 5.5-M03
\noindent\begin{minipage}{\linewidth}

\subsubsection{5.5-M03 · MC}\label{m03-mc-16}

\emph{Source:
\href{../Lecture\%20Slides/5.5.\%20Generative\%20Models.pdf\#page=7}{PDF
pp.~7--8}.}

The latent vector in a basic autoencoder is:

\begin{itemize}
\tightlist
\item
  A. Necessarily a one-hot vocabulary vector.
\item
  B. The optimizer's learning-rate schedule.
\item
  C. Necessarily the true semantic class label.
\item
  D. An encoded representation used by the decoder.
\end{itemize}

\begin{solution}[8mm]{5.5-M03}
% Write A, B, C, or D here. You may add a brief reason.

\end{solution}
\end{minipage}\par\medskip

% QUESTION 5.5-M04
\noindent\begin{minipage}{\linewidth}

\subsubsection{5.5-M04 · MC}\label{m04-mc-16}

\emph{Source:
\href{../Lecture\%20Slides/5.5.\%20Generative\%20Models.pdf\#page=8}{PDF
pp.~8}.}

Compared with a deterministic autoencoder, a VAE typically encodes:

\begin{itemize}
\tightlist
\item
  A. Parameters of a distribution over latent variables.
\item
  B. Only an adversarial real/fake decision.
\item
  C. Only one hard class index.
\item
  D. Only the reconstruction error.
\end{itemize}

\begin{solution}[8mm]{5.5-M04}
% Write A, B, C, or D here. You may add a brief reason.

\end{solution}
\end{minipage}\par\medskip

% QUESTION 5.5-M05
\noindent\begin{minipage}{\linewidth}

\subsubsection{5.5-M05 · MC}\label{m05-mc-16}

\emph{Source:
\href{../Lecture\%20Slides/5.5.\%20Generative\%20Models.pdf\#page=8}{PDF
pp.~8}.}

The KL regularization term in a VAE encourages:

\begin{itemize}
\tightlist
\item
  A. Removal of the reconstruction objective.
\item
  B. Every reconstruction to be identical.
\item
  C. Infinite latent variance for every input.
\item
  D. The latent distribution to agree with a chosen prior.
\end{itemize}

\begin{solution}[8mm]{5.5-M05}
% Write A, B, C, or D here. You may add a brief reason.

\end{solution}
\end{minipage}\par\medskip

% QUESTION 5.5-M06
\noindent\begin{minipage}{\linewidth}

\subsubsection{5.5-M06 · MC}\label{m06-mc-16}

\emph{Source:
\href{../Lecture\%20Slides/5.5.\%20Generative\%20Models.pdf\#page=9}{PDF
pp.~9--12}.}

In a GAN, the discriminator learns to:

\begin{itemize}
\tightlist
\item
  A. Distinguish real data from generated samples.
\item
  B. Generate images from noise directly.
\item
  C. Compute PCA directions.
\item
  D. Assign a unique latent vector to every real image.
\end{itemize}

\begin{solution}[8mm]{5.5-M06}
% Write A, B, C, or D here. You may add a brief reason.

\end{solution}
\end{minipage}\par\medskip

% QUESTION 5.5-M07
\noindent\begin{minipage}{\linewidth}

\subsubsection{5.5-M07 · MC}\label{m07-mc-16}

\emph{Source:
\href{../Lecture\%20Slides/5.5.\%20Generative\%20Models.pdf\#page=11}{PDF
pp.~11--12}.}

The GAN generator is trained to:

\begin{itemize}
\tightlist
\item
  A. Output only its noise input unchanged.
\item
  B. Replace every real training example with a label.
\item
  C. Produce samples the discriminator judges as real.
\item
  D. Maximize the discriminator's ability to reject its samples.
\end{itemize}

\begin{solution}[8mm]{5.5-M07}
% Write A, B, C, or D here. You may add a brief reason.

\end{solution}
\end{minipage}\par\medskip

% QUESTION 5.5-M08
\noindent\begin{minipage}{\linewidth}

\subsubsection{5.5-M08 · MC}\label{m08-mc-16}

\emph{Source:
\href{../Lecture\%20Slides/5.5.\%20Generative\%20Models.pdf\#page=13}{PDF
pp.~13}.}

A generator produces only a few similar outputs despite varied noise
inputs. This is:

\begin{itemize}
\tightlist
\item
  A. Feature normalization.
\item
  B. Gradient clipping.
\item
  C. Teacher forcing.
\item
  D. Mode collapse.
\end{itemize}

\begin{solution}[8mm]{5.5-M08}
% Write A, B, C, or D here. You may add a brief reason.

\end{solution}
\end{minipage}\par\medskip

% QUESTION 5.5-M09
\noindent\begin{minipage}{\linewidth}

\subsubsection{5.5-M09 · MC}\label{m09-mc-16}

\emph{Source:
\href{../Lecture\%20Slides/5.5.\%20Generative\%20Models.pdf\#page=17}{PDF
pp.~17--18}.}

The forward process in a basic diffusion model:

\begin{itemize}
\tightlist
\item
  A. Removes all noise in one exact inverse operation.
\item
  B. Progressively adds noise to data.
\item
  C. Learns to classify data into semantic categories.
\item
  D. Uses a discriminator to reject samples.
\end{itemize}

\begin{solution}[8mm]{5.5-M09}
% Write A, B, C, or D here. You may add a brief reason.

\end{solution}
\end{minipage}\par\medskip

% QUESTION 5.5-M10
\noindent\begin{minipage}{\linewidth}

\subsubsection{5.5-M10 · MC}\label{m10-mc-16}

\emph{Source:
\href{../Lecture\%20Slides/5.5.\%20Generative\%20Models.pdf\#page=18}{PDF
pp.~18--21}.}

Sampling with a standard diffusion model starts from noise and uses:

\begin{itemize}
\tightlist
\item
  A. A fixed lookup table of training images.
\item
  B. Only an image classifier's class index.
\item
  C. Only the forward noise-addition process.
\item
  D. A learned sequence of denoising steps.
\end{itemize}

\begin{solution}[8mm]{5.5-M10}
% Write A, B, C, or D here. You may add a brief reason.

\end{solution}
\end{minipage}\par\medskip

\subsection{Short answers}\label{short-answers-17}

% QUESTION 5.5-S01
\noindent\begin{minipage}{\linewidth}

\subsubsection{5.5-S01 · Short}\label{s01-short-16}

\emph{Source:
\href{../Lecture\%20Slides/5.5.\%20Generative\%20Models.pdf\#page=3}{PDF
pp.~3--6}.}

Contrast learning \(p(y\mid x)\) for classification with learning a
model that can generate data. Can generative models be conditioned on
labels or prompts?

\begin{solution}[20mm]{5.5-S01}
% Write your answer here, usually 1--3 sentences.

\end{solution}
\end{minipage}\par\medskip

% QUESTION 5.5-S02
\noindent\begin{minipage}{\linewidth}

\subsubsection{5.5-S02 · Short}\label{s02-short-16}

\emph{Source:
\href{../Lecture\%20Slides/5.5.\%20Generative\%20Models.pdf\#page=7}{PDF
pp.~7--8}.}

Why can reconstructing inputs through a constrained latent
representation encourage useful feature learning?

\begin{solution}[20mm]{5.5-S02}
% Write your answer here, usually 1--3 sentences.

\end{solution}
\end{minipage}\par\medskip

% QUESTION 5.5-S03
\noindent\begin{minipage}{\linewidth}

\subsubsection{5.5-S03 · Short}\label{s03-short-16}

\emph{Source:
\href{../Lecture\%20Slides/5.5.\%20Generative\%20Models.pdf\#page=8}{PDF
pp.~8}.}

Why might a basic autoencoder decode an arbitrary random latent vector
poorly? What changes in a VAE make sampling more structured?

\begin{solution}[20mm]{5.5-S03}
% Write your answer here, usually 1--3 sentences.

\end{solution}
\end{minipage}\par\medskip

% QUESTION 5.5-S04
\noindent\begin{minipage}{\linewidth}

\subsubsection{5.5-S04 · Short}\label{s04-short-16}

\emph{Source:
\href{../Lecture\%20Slides/5.5.\%20Generative\%20Models.pdf\#page=8}{PDF
pp.~8}.}

What trade-off exists between reconstruction quality and latent-prior
regularization in a VAE objective?

\begin{solution}[20mm]{5.5-S04}
% Write your answer here, usually 1--3 sentences.

\end{solution}
\end{minipage}\par\medskip

% QUESTION 5.5-S05
\noindent\begin{minipage}{\linewidth}

\subsubsection{5.5-S05 · Short}\label{s05-short-16}

\emph{Source:
\href{../Lecture\%20Slides/5.5.\%20Generative\%20Models.pdf\#page=11}{PDF
pp.~11--13}.}

Why can GAN training be unstable when the generator and discriminator
learn at different rates?

\begin{solution}[20mm]{5.5-S05}
% Write your answer here, usually 1--3 sentences.

\end{solution}
\end{minipage}\par\medskip

% QUESTION 5.5-S06
\noindent\begin{minipage}{\linewidth}

\subsubsection{5.5-S06 · Short}\label{s06-short-16}

\emph{Source:
\href{../Lecture\%20Slides/5.5.\%20Generative\%20Models.pdf\#page=17}{PDF
pp.~17--18}.}

Distinguish the fixed forward noising process from the learned reverse
process in diffusion.

\begin{solution}[20mm]{5.5-S06}
% Write your answer here, usually 1--3 sentences.

\end{solution}
\end{minipage}\par\medskip

% QUESTION 5.5-S07
\noindent\begin{minipage}{\linewidth}

\subsubsection{5.5-S07 · Short}\label{s07-short-16}

\emph{Source:
\href{../Lecture\%20Slides/5.5.\%20Generative\%20Models.pdf\#page=5}{PDF
pp.~5, 21--22}.}

Contrast next-token prediction with diffusion denoising. Why can
standard iterative diffusion sampling require many model evaluations per
generated image?

\begin{solution}[20mm]{5.5-S07}
% Write your answer here, usually 1--3 sentences.

\end{solution}
\end{minipage}\par\medskip

\subsection{Easy calculations}\label{easy-calculations-17}

% QUESTION 5.5-C01
\noindent\begin{minipage}{\linewidth}

\subsubsection{5.5-C01 · Calculation}\label{c01-calculation-16}

\emph{Source:
\href{../Lecture\%20Slides/5.5.\%20Generative\%20Models.pdf\#page=7}{PDF
pp.~7}.}

An autoencoder reconstructs input \((1,3)\) as \((2,1)\). Compute the
mean squared reconstruction error over these two components.

\begin{solution}[26mm]{5.5-C01}
% Write your short calculation and final result here. Use $...$ or \[...\] for maths.

\end{solution}
\end{minipage}\par\medskip

% QUESTION 5.5-C02
\noindent\begin{minipage}{\linewidth}

\subsubsection{5.5-C02 · Calculation}\label{c02-calculation-16}

\emph{Source:
\href{../Lecture\%20Slides/5.5.\%20Generative\%20Models.pdf\#page=8}{PDF
pp.~8}.}

A toy VAE objective is \(L=L_{recon}+\beta D_{KL}\). If \(L_{recon}=3\),
\(D_{KL}=2\), and \(\beta=0.5\), compute \(L\). What happens to this
fixed example's loss if \(\beta\) increases?

\begin{solution}[26mm]{5.5-C02}
% Write your short calculation and final result here. Use $...$ or \[...\] for maths.

\end{solution}
\end{minipage}\par\medskip

% QUESTION 5.5-C03
\noindent\begin{minipage}{\linewidth}

\subsubsection{5.5-C03 · Calculation}\label{c03-calculation-16}

\emph{Source:
\href{../Lecture\%20Slides/5.5.\%20Generative\%20Models.pdf\#page=17}{PDF
pp.~17}.}

For a toy scalar forward-noising step, use \(x_t=0.8x_0+0.6\epsilon\).
With clean value \(x_0=2\) and sampled noise \(\epsilon=-1\), compute
\(x_t\). The coefficients and sampled noise are supplied; no
distribution calculation is needed.

\begin{solution}[26mm]{5.5-C03}
% Write your short calculation and final result here. Use $...$ or \[...\] for maths.

\end{solution}
\end{minipage}\par\medskip

\clearpage

\section{Lab 4. manolo\_presentation --- Explainability and model
compression}\label{lab-4.-manolo_presentation-explainability-and-model-compression}

\textbf{Slides:}
\href{../Lab\%20Sessions/Lab4/manolo_presentation.pdf}{manolo\_presentation.pdf}.
\textbf{20 questions: 10 MC · 7 short answers · 3 easy calculations.}

Supplement from the Lab 4 briefing. This is separate from the numbered
lecture PDFs.

\subsection{Multiple choice}\label{multiple-choice-18}

% QUESTION L4-M01
\noindent\begin{minipage}{\linewidth}

\subsubsection{L4-M01 · MC}\label{l4-m01-mc}

\emph{Source:
\href{../Lab\%20Sessions/Lab4/manolo_presentation.pdf\#page=2}{PDF
pp.~2, 7}.}

Grad-CAM highlights:

\begin{itemize}
\tightlist
\item
  A. Only the brightest input pixels.
\item
  B. The ground-truth segmentation mask automatically.
\item
  C. Regions supporting a selected class score according to a
  gradient-based attribution.
\item
  D. The exact causal explanation for every decision.
\end{itemize}

\begin{solution}[8mm]{L4-M01}
% Write A, B, C, or D here. You may add a brief reason.

\end{solution}
\end{minipage}\par\medskip

% QUESTION L4-M02
\noindent\begin{minipage}{\linewidth}

\subsubsection{L4-M02 · MC}\label{l4-m02-mc}

\emph{Source:
\href{../Lab\%20Sessions/Lab4/manolo_presentation.pdf\#page=7}{PDF
pp.~7}.}

A plausible-looking attribution heatmap is:

\begin{itemize}
\tightlist
\item
  A. Proof that the model learned the intended causal mechanism.
\item
  B. Evidence to investigate, with limitations and possible failure
  modes.
\item
  C. A substitute for evaluating predictions.
\item
  D. A guarantee of fairness across all groups.
\end{itemize}

\begin{solution}[8mm]{L4-M02}
% Write A, B, C, or D here. You may add a brief reason.

\end{solution}
\end{minipage}\par\medskip

% QUESTION L4-M03
\noindent\begin{minipage}{\linewidth}

\subsubsection{L4-M03 · MC}\label{l4-m03-mc}

\emph{Source:
\href{../Lab\%20Sessions/Lab4/manolo_presentation.pdf\#page=8}{PDF
pp.~8--9}.}

FP32-to-INT8 quantization primarily changes:

\begin{itemize}
\tightlist
\item
  A. The definition of the target labels.
\item
  B. The number of output classes.
\item
  C. Every convolution's spatial dimensions.
\item
  D. The numerical representation and precision of model values.
\end{itemize}

\begin{solution}[8mm]{L4-M03}
% Write A, B, C, or D here. You may add a brief reason.

\end{solution}
\end{minipage}\par\medskip

% QUESTION L4-M04
\noindent\begin{minipage}{\linewidth}

\subsubsection{L4-M04 · MC}\label{l4-m04-mc}

\emph{Source:
\href{../Lab\%20Sessions/Lab4/manolo_presentation.pdf\#page=9}{PDF
pp.~9}.}

Calibration for post-training static quantization is used to:

\begin{itemize}
\tightlist
\item
  A. Delete the smallest weights.
\item
  B. Learn a new semantic label for each neuron.
\item
  C. Create new output classes.
\item
  D. Estimate representative activation ranges for quantization
  settings.
\end{itemize}

\begin{solution}[8mm]{L4-M04}
% Write A, B, C, or D here. You may add a brief reason.

\end{solution}
\end{minipage}\par\medskip

% QUESTION L4-M05
\noindent\begin{minipage}{\linewidth}

\subsubsection{L4-M05 · MC}\label{l4-m05-mc}

\emph{Source:
\href{../Lab\%20Sessions/Lab4/manolo_presentation.pdf\#page=10}{PDF
pp.~10}.}

Magnitude-based pruning typically targets:

\begin{itemize}
\tightlist
\item
  A. All biases regardless of importance.
\item
  B. Weights with small absolute values.
\item
  C. Only the largest positive weights.
\item
  D. Only the training examples with large loss.
\end{itemize}

\begin{solution}[8mm]{L4-M05}
% Write A, B, C, or D here. You may add a brief reason.

\end{solution}
\end{minipage}\par\medskip

% QUESTION L4-M06
\noindent\begin{minipage}{\linewidth}

\subsubsection{L4-M06 · MC}\label{l4-m06-mc}

\emph{Source:
\href{../Lab\%20Sessions/Lab4/manolo_presentation.pdf\#page=10}{PDF
pp.~10}.}

Unstructured pruning may give little speed-up with dense kernels
because:

\begin{itemize}
\tightlist
\item
  A. Zero weights always use more storage than nonzero weights.
\item
  B. The dense implementation may still store and process zero entries.
\item
  C. It prevents the model from running on a CPU.
\item
  D. It increases the number of channels by definition.
\end{itemize}

\begin{solution}[8mm]{L4-M06}
% Write A, B, C, or D here. You may add a brief reason.

\end{solution}
\end{minipage}\par\medskip

% QUESTION L4-M07
\noindent\begin{minipage}{\linewidth}

\subsubsection{L4-M07 · MC}\label{l4-m07-mc}

\emph{Source:
\href{../Lab\%20Sessions/Lab4/manolo_presentation.pdf\#page=10}{PDF
pp.~10}.}

Structured pruning can remove:

\begin{itemize}
\tightlist
\item
  A. Only the learned learning rate.
\item
  B. Only the validation labels.
\item
  C. Entire channels or filters, shrinking tensor dimensions.
\item
  D. Only isolated scalar weights.
\end{itemize}

\begin{solution}[8mm]{L4-M07}
% Write A, B, C, or D here. You may add a brief reason.

\end{solution}
\end{minipage}\par\medskip

% QUESTION L4-M08
\noindent\begin{minipage}{\linewidth}

\subsubsection{L4-M08 · MC}\label{l4-m08-mc}

\emph{Source:
\href{../Lab\%20Sessions/Lab4/manolo_presentation.pdf\#page=8}{PDF
pp.~8}.}

Knowledge distillation trains:

\begin{itemize}
\tightlist
\item
  A. A new vocabulary for the teacher.
\item
  B. A discriminator to reject synthetic images.
\item
  C. A larger model to memorize every student error.
\item
  D. A student to imitate a teacher's outputs or representations.
\end{itemize}

\begin{solution}[8mm]{L4-M08}
% Write A, B, C, or D here. You may add a brief reason.

\end{solution}
\end{minipage}\par\medskip

% QUESTION L4-M09
\noindent\begin{minipage}{\linewidth}

\subsubsection{L4-M09 · MC}\label{l4-m09-mc}

\emph{Source:
\href{../Lab\%20Sessions/Lab4/manolo_presentation.pdf\#page=5}{PDF
pp.~5, 11}.}

Unchanged average accuracy after compression:

\begin{itemize}
\tightlist
\item
  A. Establishes that latency improved.
\item
  B. Proves unchanged explanations for every image.
\item
  C. Can still hide deterioration for rare classes or subgroups.
\item
  D. Guarantees unchanged performance for every subgroup.
\end{itemize}

\begin{solution}[8mm]{L4-M09}
% Write A, B, C, or D here. You may add a brief reason.

\end{solution}
\end{minipage}\par\medskip

% QUESTION L4-M10
\noindent\begin{minipage}{\linewidth}

\subsubsection{L4-M10 · MC}\label{l4-m10-mc}

\emph{Source:
\href{../Lab\%20Sessions/Lab4/manolo_presentation.pdf\#page=6}{PDF
pp.~6}.}

One reason to run a model on an edge device is:

\begin{itemize}
\tightlist
\item
  A. Guaranteed better accuracy than the same model on a server.
\item
  B. Unlimited memory and energy.
\item
  C. No need to measure inference latency.
\item
  D. Reduced reliance on network connectivity and remote data transfer.
\end{itemize}

\begin{solution}[8mm]{L4-M10}
% Write A, B, C, or D here. You may add a brief reason.

\end{solution}
\end{minipage}\par\medskip

\subsection{Short answers}\label{short-answers-18}

% QUESTION L4-S01
\noindent\begin{minipage}{\linewidth}

\subsubsection{L4-S01 · Short}\label{l4-s01-short}

\emph{Source:
\href{../Lab\%20Sessions/Lab4/manolo_presentation.pdf\#page=2}{PDF
pp.~2, 7}.}

Why is Grad-CAM usually spatially coarse? How is this related to the
resolution of the selected feature maps?

\begin{solution}[20mm]{L4-S01}
% Write your answer here, usually 1--3 sentences.

\end{solution}
\end{minipage}\par\medskip

% QUESTION L4-S02
\noindent\begin{minipage}{\linewidth}

\subsubsection{L4-S02 · Short}\label{l4-s02-short}

\emph{Source:
\href{../Lab\%20Sessions/Lab4/manolo_presentation.pdf\#page=5}{PDF
pp.~5, 7}.}

A classifier's heatmap focuses on a watermark rather than the object.
What concern does this raise, and what targeted data check could you
make?

\begin{solution}[20mm]{L4-S02}
% Write your answer here, usually 1--3 sentences.

\end{solution}
\end{minipage}\par\medskip

% QUESTION L4-S03
\noindent\begin{minipage}{\linewidth}

\subsubsection{L4-S03 · Short}\label{l4-s03-short}

\emph{Source:
\href{../Lab\%20Sessions/Lab4/manolo_presentation.pdf\#page=8}{PDF
pp.~8--10}.}

Distinguish quantization from pruning. What does each change in the
model's numerical representation or structure?

\begin{solution}[20mm]{L4-S03}
% Write your answer here, usually 1--3 sentences.

\end{solution}
\end{minipage}\par\medskip

% QUESTION L4-S04
\noindent\begin{minipage}{\linewidth}

\subsubsection{L4-S04 · Short}\label{l4-s04-short}

\emph{Source:
\href{../Lab\%20Sessions/Lab4/manolo_presentation.pdf\#page=9}{PDF
pp.~9}.}

Why should quantization calibration inputs resemble the intended
deployment data?

\begin{solution}[20mm]{L4-S04}
% Write your answer here, usually 1--3 sentences.

\end{solution}
\end{minipage}\par\medskip

% QUESTION L4-S05
\noindent\begin{minipage}{\linewidth}

\subsubsection{L4-S05 · Short}\label{l4-s05-short}

\emph{Source:
\href{../Lab\%20Sessions/Lab4/manolo_presentation.pdf\#page=10}{PDF
pp.~10}.}

Why can fine-tuning after pruning help recover performance? Does high
sparsity alone guarantee a smaller dense model file?

\begin{solution}[20mm]{L4-S05}
% Write your answer here, usually 1--3 sentences.

\end{solution}
\end{minipage}\par\medskip

% QUESTION L4-S06
\noindent\begin{minipage}{\linewidth}

\subsubsection{L4-S06 · Short}\label{l4-s06-short}

\emph{Source:
\href{../Lab\%20Sessions/Lab4/manolo_presentation.pdf\#page=10}{PDF
pp.~10--11}.}

Name three measurements, besides overall accuracy, that help assess a
compressed model's usefulness.

\begin{solution}[20mm]{L4-S06}
% Write your answer here, usually 1--3 sentences.

\end{solution}
\end{minipage}\par\medskip

% QUESTION L4-S07
\noindent\begin{minipage}{\linewidth}

\subsubsection{L4-S07 · Short}\label{l4-s07-short}

\emph{Source:
\href{../Lab\%20Sessions/Lab4/manolo_presentation.pdf\#page=5}{PDF
pp.~5, 11}.}

Why re-evaluate per-class performance and inspect explanations after
compression, even if overall accuracy barely changes?

\begin{solution}[20mm]{L4-S07}
% Write your answer here, usually 1--3 sentences.

\end{solution}
\end{minipage}\par\medskip

\subsection{Easy calculations}\label{easy-calculations-18}

% QUESTION L4-C01
\noindent\begin{minipage}{\linewidth}

\subsubsection{L4-C01 · Calculation}\label{l4-c01-calculation}

\emph{Source:
\href{../Lab\%20Sessions/Lab4/manolo_presentation.pdf\#page=9}{PDF
pp.~9}.}

One million weights are stored using either 32 or 8 bits per weight.
Ignoring metadata, compute both storage sizes in decimal MB and the
reduction factor.

\begin{solution}[26mm]{L4-C01}
% Write your short calculation and final result here. Use $...$ or \[...\] for maths.

\end{solution}
\end{minipage}\par\medskip

% QUESTION L4-C02
\noindent\begin{minipage}{\linewidth}

\subsubsection{L4-C02 · Calculation}\label{l4-c02-calculation}

\emph{Source:
\href{../Lab\%20Sessions/Lab4/manolo_presentation.pdf\#page=10}{PDF
pp.~10}.}

Pruning sets 600 of 1,000 weights to zero. Compute sparsity and the
remaining nonzero-weight count.

\begin{solution}[26mm]{L4-C02}
% Write your short calculation and final result here. Use $...$ or \[...\] for maths.

\end{solution}
\end{minipage}\par\medskip

% QUESTION L4-C03
\noindent\begin{minipage}{\linewidth}

\subsubsection{L4-C03 · Calculation}\label{l4-c03-calculation}

\emph{Source:
\href{../Lab\%20Sessions/Lab4/manolo_presentation.pdf\#page=6}{PDF
pp.~6, 11}.}

Measured inference latency falls from 40 ms to 10 ms after compression.
Compute the speed-up factor and the percentage reduction in latency.

\begin{solution}[26mm]{L4-C03}
% Write your short calculation and final result here. Use $...$ or \[...\] for maths.

\end{solution}
\end{minipage}\par\medskip

\end{document}